% CellSim — Full Technical Reference
% Auto-authored from an exhaustive source-code scan (2026-04-20).
% Companion to docs/equations/cellsim_equations.tex.
% Build: pdflatex cellsim_full_reference.tex (twice for TOC)

\documentclass[11pt,letterpaper,twoside,openany]{book}
\usepackage[margin=2.2cm,includeheadfoot]{geometry}
\usepackage{amsmath,amssymb,amsthm}
\usepackage{bm}
\usepackage{siunitx}
\usepackage{booktabs}
\usepackage{longtable}
\usepackage{array}
\usepackage{tabularx}
\usepackage{hyperref}
\usepackage{xcolor}
\usepackage{listings}
\usepackage{fancyhdr}
\usepackage{enumitem}
\usepackage[T1]{fontenc}
\usepackage{lmodern}
\usepackage{microtype}
\usepackage{makeidx}

\hypersetup{colorlinks=true,linkcolor=blue!60!black,urlcolor=blue!60!black,citecolor=blue!60!black}
\setlist{nosep}
\makeindex

% ── Listings style for C++ code ──────────────────────────────────────
\definecolor{lstkw}{rgb}{0.13,0.29,0.53}
\definecolor{lstcm}{rgb}{0.45,0.45,0.45}
\definecolor{lststr}{rgb}{0.64,0.08,0.08}
\definecolor{lstbg}{rgb}{0.98,0.98,0.96}
\lstset{
  basicstyle=\ttfamily\footnotesize,
  keywordstyle=\color{lstkw}\bfseries,
  commentstyle=\color{lstcm}\itshape,
  stringstyle=\color{lststr},
  backgroundcolor=\color{lstbg},
  frame=single,
  framesep=3pt,
  rulecolor=\color{lstbg},
  columns=flexible,
  keepspaces=true,
  showstringspaces=false,
  numbers=left,
  numberstyle=\tiny\color{lstcm},
  numbersep=6pt,
  breaklines=true,
  breakatwhitespace=true,
  language=C++,
  morekeywords={simd_float3,simd_float4,simd_float4x4,uint8_t,uint32_t,int32_t,MolTag,ParticleType,Compartment,GenePolymerase,VirionStage,BacteriumStage,ApoPhase,SimMode,MitosisPhase,ChemicalEntityKind}
}

% ── Headers / footers ────────────────────────────────────────────────
\pagestyle{fancy}
\fancyhf{}
\fancyhead[LE,RO]{\thepage}
\fancyhead[LO]{\small\textsc{\nouppercase{\rightmark}}}
\fancyhead[RE]{\small\textsc{CellSim --- Full Technical Reference}}
\renewcommand{\headrulewidth}{0.4pt}

% ── Custom commands ──────────────────────────────────────────────────
\newcommand{\code}[1]{\texttt{#1}}
\newcommand{\file}[1]{\texttt{\small#1}}
\newcommand{\sym}[1]{\ensuremath{\boldsymbol{#1}}}
\newcommand{\unitof}[1]{\,\text{#1}}
\newcommand{\bioh}{\text{bio-h}}
\newcommand{\biomin}{\text{bio-min}}
\newcommand{\biosec}{\text{bio-s}}
\newcommand{\realsec}{\text{real-s}}
\newcommand{\wu}{\text{wu}}
\newcommand{\um}{\ensuremath{\mu\text{m}}}
\newcommand{\cellref}[1]{\textsc{\scriptsize[#1]}}

% ── Title page ───────────────────────────────────────────────────────
\title{%
  \Huge\textbf{CellSim}\\[0.4em]
  \LARGE Full Technical Reference\\[0.5em]
  \Large \textsl{Every Subsystem, Every Function, Every Equation}\\[1.0em]
  \large A complete, bottom-up walkthrough of the simulation engine ---\\
  from constants and geometry, through the central dogma,\\
  to the chemistry-first drug system, pathogen infection chain,\\
  and the Metal-native rendering pipeline.
}
\author{%
  Generated from an exhaustive scan of the source tree\\
  (\file{src/main.mm}, \file{src/simulation/*}, \file{src/core/Constants.h},\\
  14 biochem modules, test harnesses, and the Phase 8A additions)\\[0.5em]
  Companion to \file{docs/equations/cellsim\_equations.tex}
}
\date{2026-04-20}

\begin{document}
\frontmatter
\maketitle

% ── Preface ──────────────────────────────────────────────────────────
\chapter*{Preface}
\addcontentsline{toc}{chapter}{Preface}

\noindent This document is a deep, \emph{bottom-up} walkthrough of
every subsystem inside \textsc{CellSim} --- the Metal-native interactive
cell biology simulator in \file{/Users/henry/CellSim}. Unlike the
\emph{equations reference} in \file{docs/equations/cellsim\_equations.tex},
which catalogues only the mathematical content, this book traces the
\emph{code}: it names every important struct, every function that
mutates cell state, every hook between modules, and every literature
source anchoring the constants.

It is written so a reader who has never opened the repository can
nonetheless follow the simulator stage by stage. Each chapter begins
with an \emph{in-a-paragraph} summary, then proceeds into the
structures, then the algorithms, then the interactions with the rest
of the engine. Lines of real code are quoted freely (with file-line
references such as \file{src/simulation/Simulation.h}:1042); equations
are stated with units and named constants; pitfalls the author has
already hit (e.g.\ mass-balance invariants, mode gating) are called
out explicitly so readers can reason about why a given line of code
exists.

\paragraph{Organisation.} The twelve parts follow the natural
dependency order of the engine. Part~\textsc{I} (Architecture) is the
coordinate system and frame loop everything else plugs into. Parts
\textsc{II}--\textsc{IV} handle intrinsic cell biology: membrane,
cycle, mitosis, apoptosis, the central dogma, metabolism, membrane
transport, vesicular traffic, cytoskeleton. Part~\textsc{V} is the
dish-level \code{MediumField} PDE. Part~\textsc{VI} is the discrete
particle / Langevin physics inside a focused cell. Part~\textsc{VII}
is the chemistry-first drug pipeline --- the first demonstration of
the ``no tropism flag, no MOA label'' discipline. Part~\textsc{VIII}
is the pathogen infection chain (virion + bacterium). Part~\textsc{IX}
is the immune system scaffold. Part~\textsc{X} is the renderer.
Part~\textsc{XI} covers UI panels, export, and the headless test
framework. Part~\textsc{XII} documents the Phase 8A multi-gene
Central Dogma upgrade that landed on 2026-04-20.

\paragraph{Discipline of the engine.} Two cross-cutting invariants
thread the whole codebase and are called out as they appear:
\begin{itemize}
  \item \textbf{Mass conservation.} Every pool that holds mass ---
        cytosol biomass, membrane mass, receptor mass, ATP, medium
        species, intracellular metabolites --- must be sourced and
        sinked by functions that move mass between pools. Nothing is
        created or destroyed except by well-defined flux terms. The
        \code{MediumField::checkBalance()} call drifts less than
        $10^{-3}$ across a 46 bio-h simulation.
  \item \textbf{No prescription, only emergence.} Drugs, viruses, and
        bacteria carry \emph{only} structural data (SMILES,
        pharmacophore, descriptor vector). The effect of applying any
        of them on a cell is computed by the simulator
        --- not encoded by a ``mechanism of action''
        flag. A scrambled descriptor yields no binding; a matched
        descriptor does. Negative controls are run in the headless
        test suite to prove this.
\end{itemize}

\paragraph{How to read the code excerpts.} Listings cite the file
and (when relevant) line number. Bold-face types such as
\code{SimCell} or \code{CentralDogmaState} are clickable in the PDF;
equation labels of the form \eqref{eq:fick} are reused across
chapters.

\paragraph{A word on length.} The source tree is roughly
$23\,000$ lines of C\texttt{++}, which is why this document is long.
The engine lives at the intersection of biology, chemistry, graphics,
and numerical methods; cutting corners in documentation produces a
reader who cannot debug the simulator. This book is meant to be the
artefact a new collaborator reads, cover to cover, before they make
their first change.

\tableofcontents
\listoftables
\listoffigures
\mainmatter

% ═════════════════════════════════════════════════════════════════════
\part{Architecture and Foundations}
% ═════════════════════════════════════════════════════════════════════

\chapter{The Shape of the Simulator}
\label{ch:arch}

\textsc{CellSim} is organised as three concentric layers. At the core
sits \code{Simulation} (\file{src/simulation/Simulation.h}, 2826
lines), a self-contained C++ class that owns the population of cells
(\code{std::vector<SimCell> cells}), the dish chemistry (a
\code{MediumField} object named \code{nutrients}), the per-tick stats,
and every update routine that advances biology by one step. Around
this core sits the \emph{biochem module layer} in
\file{src/simulation/biochem/*} --- twenty-odd headers that define
individual sub-processes (apoptosis, central dogma, metabolism,
vesicular traffic, membrane transport, signalling, cytoskeleton,
gene regulation, DNA repair, pathogens, immunity, proteome, etc.).
These are self-contained: each could be dropped into another code
base with its forward-declared \code{SimCell} reference. Wrapping
everything is the \emph{application shell} in \file{src/main.mm}
(9159 lines), which carries the window, Metal rendering pipelines,
ImGui UI panels, particle-level intracellular renderer, CSV export,
and every interactive input.

\section{The coordinate system}
\label{sec:coord}

Every spatial constant in the simulator descends from a single anchor
in \file{src/core/Constants.h}:
\begin{lstlisting}
constexpr float MICROMETERS_PER_WORLD_UNIT = 5.0f;
constexpr float UM_TO_WU = 1.0f / MICROMETERS_PER_WORLD_UNIT;
\end{lstlisting}
One world unit (\wu) equals 5\,\um. A HeLa cell, whose median radius
is 10\,\um (Alberts MBoC Ch 1; BioNumbers BNID 100434; our own
\file{data/reference/hela\_reference.csv}), therefore has a simulated
radius of 2\,\wu when \code{CELL\_RADIUS\_BASE = 2.0f}. The scene
bound is $\pm\,55\,\wu$, corresponding to a rectangular sub-region
$275\,\um \times 550\,\um$ on a virtual dish floor at
\code{FLOOR\_Y = -7.0f}. This is deliberately a \emph{magnified
sub-region} of a real 10\,cm petri dish rather than the whole plate:
it lets the user see individual cell mechanics at a legible scale
without pretending to simulate a 100\,000-cell colony in full.

The dish chemistry lives on a
$\code{NUTRIENT\_GRID\_N} = 64$ square grid; each voxel therefore
measures $55 \cdot 2 / 64 \approx 1.72\,\wu \approx 8.6\,\um$ on a
side. This resolution is fine enough that a cell (diameter
$\sim 4\,\wu$) spans multiple voxels, so local depletion effects
around a single cell are captured, while still remaining cheap
($64^2 = 4\,096$ cells per species, twelve species).

\section{The time base}
\label{sec:time}

Biological time and wall-clock time are decoupled through a single
constant, \code{BIO\_TIME\_SCALE = 180.0f} (bio-seconds per real-
second). One real second of wall-clock, at the default
\code{timeScale = 1.0f}, produces 180\,bio-s of simulated time;
one real minute is $3\,\bioh$. The runtime scaling slider
(\code{gSim.timeScale}) multiplies this further. Because some
intracellular processes are stiff at large time steps, several
subsystems compress further through additional scales:
\begin{align*}
  \code{FAST\_DT\_SCALE}   &= 1.0         & \text{fast dynamics (receptor kinetics)} \\
  \code{MEDIUM\_DT\_SCALE} &= 0.20096     & \text{mid dynamics (cycle, fate)}        \\
  \code{SLOW\_DT\_SCALE}   &= 0.055       & \text{slow dynamics (metabolism, drug PK)} \\
  \code{CDK\_DT\_SCALE}    &= 0.8         & \text{CDK/cyclin ODE stiffness control}
\end{align*}
These values are the result of 2026-04 calibration against the
CTC Fluo-N2DL-HeLa \texttt{seq02} reference dataset; changing them
moves the mean relative error bar on cell-count growth, which is
currently held at $5.1\%$ over 46\,\bioh.

\section{The main update loop}
\label{sec:mainloop}

The heart of the simulator is \code{Simulation::update(float dt)}
(\file{src/simulation/Simulation.h} around line 1160). Given a
real-time step \code{dt}, it advances every cell's biology by a
\emph{scaled} time \code{dt} $\times$ \code{timeScale}. To keep the
biology stable regardless of how bursty the host frame rate is, the
function uses \emph{dt accumulation} and \emph{per-frame capping}:

\begin{lstlisting}
float requestedDt = dt * timeScale;
pendingScaledDt += requestedDt;
float responsiveCap = MAX_SCALED_DT_PER_FRAME_COLONY;
if (mode == MODE_SINGLE_CELL) {
    responsiveCap = (cells.size() <= 2)
        ? MAX_SCALED_DT_PER_FRAME_SINGLE
        : MAX_SCALED_DT_PER_FRAME_SINGLE_LINEAGE;
}
float totalDt = fminf(pendingScaledDt, responsiveCap);
if (totalDt <= 0.0f) return;
pendingScaledDt -= totalDt;
float subDt = fminf(totalDt, MAX_SUBSTEP_DT);
int steps = (int)ceilf(totalDt / fmaxf(subDt, 1e-6f));
subDt = totalDt / fmaxf((float)steps, 1.0f);
\end{lstlisting}

Each inner step then drives the physics, per-cell metabolism, cycle,
mitosis, central dogma, fate, apoptosis, drug PK, pathogen update,
and medium diffusion, in that deliberate order:

\begin{lstlisting}
for (int s=0; s<steps; s++) {
    updatePhysics(subDt);
    for (auto& c : cells) { if(!c.alive) continue;
        updateMetabolism(c, subDt);
        updateMitochondria(c, subDt);
        updateAdhesion(c, subDt);
        updateCellCycle(c, subDt); updateMitosisProgram(c, subDt);
        c.program.ensureCDogmaInitialized();
        float dogmaDt = fminf(fmaxf(subDt, 1.0f/240.0f), 1.0f/30.0f);
        c.program.cdogma.update(dogmaDt, c.phase);
        updateFate(c, subDt); updateApoptosis(c, subDt);
    }
    updateDrugPK(subDt);
    statInfectedCells = 0;
    updatePathogens(subDt);
    nutrients.diffuse(subDt, envO2, envGlucose);
}
\end{lstlisting}

The order matters biologically. Physics integrates first so cell
positions are fresh before we compute crowding pressure and cell-cell
adhesion. Metabolism runs before the cycle because ATP availability
gates CDK activity. Central Dogma runs before fate because a
replication stall drives DNA-damage response, which feeds the fate
decision. Drug PK runs last \emph{per cell} so target occupancy
effects apply to next tick's biology (a one-tick delay, a common
stability trick in stiff integrators). Pathogen physics runs after
all per-cell mutations so a pathogen that just budded out is seen
with its final host-cell state. Finally the medium diffuses.

\section{Execution modes}
\label{sec:modes}

Two enumerated modes gate much of the higher-level behaviour:
\begin{lstlisting}
enum SimMode { MODE_SINGLE_CELL = 0, MODE_COLONY = 1 };
\end{lstlisting}
In \code{MODE\_SINGLE\_CELL}, the simulator initialises one fully-
detailed cell (hundreds of \code{IntraParticle} instances filling
the interior: ribosomes, tRNA, polymerases, Golgi, ER, mitochondria,
DNA helix, etc.) and renders its interior. In \code{MODE\_COLONY}, it
initialises many simpler cells (by default \code{INIT\_CELLS = 5}, or
user-specified) and skips the per-cell interior renderer entirely.
The two modes share the same biology (cycle, fate, apoptosis, drug
system, pathogens) but diverge on rendering cost. Phase 7A's
\emph{Test Mode (Lite)} (discussed in Ch.~\ref{ch:testmode}) actually
runs in colony mode with a single cell so a 16\,GB machine can render
it smoothly.

\chapter{Constants as the Single Source of Truth}
\label{ch:constants}

\file{src/core/Constants.h} (523 lines) is the \emph{only} place
where scientific and geometric constants live. Every other file
includes it. This discipline is load-bearing: if a constant needs
adjustment (for calibration, for a new cell type, for reproducing
an experiment), it must be tuned \emph{here} and \emph{only here}.
Reviewers can audit the file in isolation to verify scientific
fidelity.

\section{Geometry and capacity}

Key scale-setting values:
\begin{longtable}{lrl}
\toprule
\textbf{Constant} & \textbf{Value} & \textbf{Meaning}\\ \midrule
\code{MICROMETERS\_PER\_WORLD\_UNIT}     & $5.0$       & 1\,\wu = 5\,\um\\
\code{CELL\_RADIUS\_BASE}                 & $2.0$       & 10\,\um (HeLa)\\
\code{FLOOR\_Y}                           & $-7.0$      & dish floor in world\\
\code{SCENE\_BOUND}                       & $55.0$      & half-extent ($\pm 275\,\um$)\\
\code{MAX\_CELLS}                         & $2000$      & memory cap\\
\code{INIT\_CELLS}                        & $5$         & colony default\\
\code{SINGLE\_CELL\_COUNT}                & $1$         & single-cell default\\
\code{MONOLAYER\_CAPACITY}                & $600$       & soft cap triggering crowding effects\\
\code{MITO\_COUNT\_G1}/\code{G2}          & $3/6$       & mitochondrial fission across cycle\\
\code{GOLGI\_COUNT\_G1}/\code{G2}         & $1/2$       & Golgi fragmentation\\
\code{NUTRIENT\_GRID\_N}                  & $64$        & medium grid resolution\\
\bottomrule
\end{longtable}

\section{Time scales}

Three orders of magnitude separate metabolism from the cycle, which
motivates the time-scale fanout:
\begin{lstlisting}
constexpr float BIO_TIME_SCALE  = 180.0f;   // bio-s per real-s
constexpr float BIO_MIN_PER_SEC = 3.0f;     // 180/60
constexpr float FAST_DT_SCALE   = 1.0f;
constexpr float MEDIUM_DT_SCALE = 0.20096f;
constexpr float SLOW_DT_SCALE   = 0.055f;
constexpr float CDK_DT_SCALE    = 0.8f;
constexpr float LAG_DURATION    = 0.3f;
constexpr float LAG_BIOSEC      = 30.0f * 60.0f;
\end{lstlisting}

The calibration story of \code{SLOW\_DT\_SCALE} is instructive. In
2026-03 it was \(0.052\); after the CTC seq02 mean error analysis
converged on a sigmoid over-shoot of \(6.4\%\), raising it to
\(0.055\) and raising \code{MEDIUM\_DT\_SCALE} from \(0.190\) to
\(0.201\) reduced the error to \(5.1\%\) and kept all five
independent seeds within \(\pm 0.2\%\). These are not arbitrary
numbers; they are the output of a grid sweep documented in
\file{scripts/calibrate\_rate.sh}.

\section{Cycle targets}

The four cycle-phase durations (in bio-seconds) drive the cell-cycle
timer and, indirectly, the proliferation rate:
\begin{lstlisting}
constexpr float CYCLE_G1_TARGET_BIOSEC = 11.0f * 3600.0f;   // 11 bio-h
constexpr float CYCLE_S_TARGET_BIOSEC  =  8.0f * 3600.0f;   //  8 bio-h
constexpr float CYCLE_G2_TARGET_BIOSEC =  4.0f * 3600.0f;   //  4 bio-h
constexpr float CYCLE_M_TARGET_BIOSEC  =  1.0f * 3600.0f;   //  1 bio-h
\end{lstlisting}
Sum: 24\,\bioh. This matches the classical HeLa generation time
(doubling time 24\,h at 37$^\circ$C in DMEM + 10\% FBS; Chao 2019
supplementary Fig.~1). The individual phase durations are tuned so
that the S/G2 ratio reproduces the measured BrdU pulse-chase
distribution.

\section{Mechanotransduction}

\begin{lstlisting}
constexpr float MECH_P21_COUPLING    = 0.002f;
constexpr float CONTACT_INHIBIT_NBRS = 6.0f;
\end{lstlisting}
The \code{MECH\_P21\_COUPLING} constant is the coupling strength
between neighbour-count pressure and the p21 checkpoint, implementing
contact inhibition: above six neighbours, p21 rises at
$\code{MECH\_P21\_COUPLING} \cdot (\text{neighbours} -
\code{CONTACT\_INHIBIT\_NBRS})\cdot dt$. A single cell never rises
above this baseline because its neighbour count is zero.

\section{Apoptosis namespace}

All apoptosis-engine constants live in a nested namespace to keep
the global Constants.h tidy:
\begin{lstlisting}
namespace Apoptosis {
    enum Phase : uint8_t {
        ALIVE=0, PRIMED=1, MOMP=2, EXECUTION=3,
        FRAGMENTATION=4, BODIES=5, CLEARED=6
    };
    constexpr int   BODIES_PER_CELL_MIN = 5;
    constexpr int   BODIES_PER_CELL_MAX = 15;
    constexpr float MOMP_DURATION_BIOSEC          = 300.0f;
    constexpr float EXECUTION_DURATION_BIOSEC     = 1800.0f;
    constexpr float FRAGMENTATION_DURATION_BIOSEC = 1800.0f;
    // ... many more — see src/core/Constants.h
}
\end{lstlisting}

The two counts \code{BODIES\_PER\_CELL\_MIN/MAX} govern how many
apoptotic bodies a dying cell fragments into; each body drifts and
decomposes on its own. The three dwell times come from Ye 2017
\emph{Sci Rep} (MOMP), Spencer 2009 \emph{Nature} (execution phase),
and Silva 2010 \emph{FEBS Lett} (secondary necrosis).

% ═════════════════════════════════════════════════════════════════════
\part{Cell Biology}
% ═════════════════════════════════════════════════════════════════════

\chapter{The SimCell Structure}
\label{ch:simcell}

\code{SimCell} is the per-cell record. It lives in
\file{src/simulation/Simulation.h} lines 281--640 and carries every
scalar state variable, substructure, and flag that governs one cell's
trajectory. Logically it is divided into eight groups of fields,
summarised below and walked through in the next sections.

The struct is deliberately a plain-data ``bag of fields'' rather than
a class with encapsulated invariants. The decision is pragmatic: the
simulator's hot loop iterates a \code{std::vector<SimCell>} and reads
or writes arbitrary combinations of these fields dozens of times per
tick. Encapsulating behind getters/setters would force gratuitous
function calls; embedding sub-objects with their own invariants would
multiply the number of places a field can change in an emergency
debugging session. Instead, the discipline is: all mutations go
through named \code{Simulation::updateXxx(c, dt)} routines, each of
which is a free function in the \code{Simulation} class so it can be
grepped by name.

The total sizeof(SimCell) sits in the low kilobytes. With the
default \code{MAX\_CELLS = 2000} cap, the population fits within a
few megabytes --- small enough that an entire population can be
serialized and exported without streaming. That observation drives
the Export system's one-shot snapshot design (Ch.~\ref{ch:export}).

\section{Field-group overview}

\begin{longtable}{lp{9.5cm}}
\toprule
\textbf{Group} & \textbf{Fields}\\ \midrule
Spatial       & \code{position}, \code{velocity}, \code{radius}, \code{size}, \code{alive}\\
Bioenergetics & \code{ATP}, \code{biomass}, \code{mitoPotential}, \code{mitoHealth}, \code{glycolytic}, \code{warburgTimer}, \code{hypoxiaTimer}, \code{hypoxiaIntensity}, \code{necrotic}\\
Cycle         & \code{cdk} (struct), \code{phase}, \code{cycleTimer}, \code{cycleProgress}, \code{checkpointG1Passed}, \code{checkpointG2Passed}, \code{g1WaitTimer}, \code{g2WaitTimer}\\
Damage/stress & \code{stress}, \code{ROS}, \code{damageLevel}, \code{age}\\
Fate decision & \code{fate} (int), \code{fateScores[3]}, \code{fateTimer}, \code{fateLocked}, \code{atpDangerTimer}\\
Intrinsic     & \code{glycolysisBias}, \code{prolifBias}, \code{rosTolerance}, \code{repairRate}, \code{generation}, \code{cloneId}, \code{telomere}, \code{senescent}\\
Division      & \code{divisionPending}, \code{localPressure}, \code{motileAngle}, \code{motileSpeed}, \code{divisionCooldown}, \code{postDivisionRecovery}\\
Apoptosis     & \code{apo} (engine), \code{apoPhase}, \code{membraneMass\_bm}, \code{receptorMass\_bm}, \code{initialBiomassAtDeath}, \code{initialMembraneAtDeath}, \code{initialReceptorAtDeath}, \code{chronicPressureBioSec}, \code{senescenceBioSec}, \code{fasLExposure}, \code{lysosomalLoad\_{cyto,mem,rec}}, \code{damageSensedBioSec}, \code{bodiesSpawned}\\
Drug/path     & \code{drugIntra\_uM}, \code{targetBound\_uM}, \code{targetCount}, \code{targetOccupancy}, \code{intraVirions}, \code{intraBacteria}, \code{assembledVirions}, \code{pathogensReleased}\\
Central dogma & inside \code{program} (CellCycleProgram) which holds \code{cdogma}, \code{mitosis}, \code{events}\\
\bottomrule
\end{longtable}

\section{Initialisation}

\code{SimCell::init(simd\_float3 pos, int i)} sets sensible defaults
for every field. Key initialisations:

\begin{lstlisting}
float membraneFrac = Apoptosis::MEMBRANE_MASS_PER_BIOMASS;  // 0.20
float receptorFrac = Apoptosis::RECEPTOR_MASS_PER_BIOMASS;  // 0.07
membraneMass_bm = membraneFrac * biomass;
receptorMass_bm = receptorFrac * biomass;
\end{lstlisting}
The mass-conservation ledger is populated \emph{at init} so that when
the cell later dies, the fractions released into the medium sum to
the original biomass + membrane + receptor pool. Without this ledger,
the closed-system invariant would drift every time a cell lyses.

\section{CDK state}

\begin{lstlisting}
struct CDKState {
    float CycD, CycE, CycA, CycB;
    float Rb, p21, p53;
    float E2F;
    float APC_C;
    // derived flags
    bool readyForS() const;
    bool readyForM() const;
};
\end{lstlisting}
The cyclin-dependent kinase state is driven by the ODE in
\code{updateCellCycle} (Ch.~\ref{ch:cycle}): CycD rises during G1 in
response to growth factors, phosphorylates Rb, releases E2F, E2F
transcribes CycE, CycE + CycA drive S-phase, CycB peaks at mitosis
entry, APC/C degrades CycB at mitotic exit. The literature source is
Aldridge 2006 \emph{Nat Cell Biol} coupled to Kholodenko 2006
\emph{Nat Rev Mol Cell Biol} ERK--MAPK; we keep a lightweight
four-cyclin model rather than the full 25-species version because
the marginal accuracy is smaller than the calibration uncertainty of
the CTC reference.

\chapter{The Cell Cycle}
\label{ch:cycle}

Cells progress through G1 \(\to\) S \(\to\) G2 \(\to\) M with
literature-anchored durations (Ch.~\ref{ch:constants}). The driver is
\code{Simulation::updateCellCycle(SimCell\& c, float dt)}
(\file{src/simulation/Simulation.h} around line 1860).

\section{Phase advancement}

Phase is an integer \code{c.phase} in \{0,1,2,3\} for G1/S/G2/M.
Each tick the phase timer \code{c.cycleTimer} advances by the scaled
bio-second \code{dt * \code{MEDIUM\_DT\_SCALE} * 3600}. When it exceeds
the phase target, the cell requests the next phase via the
corresponding checkpoint. Checkpoints are gated in three helpers:

\begin{lstlisting}
static CheckpointResult evalG1S(const SimCell& c,
                                float envGF, const CentralDogmaState& cdogma);
static CheckpointResult evalG2M(const SimCell& c,
                                const CentralDogmaState& cdogma);
static CheckpointResult evalSAC(const MitosisState& mitosis);
\end{lstlisting}

\code{evalG1S} tests growth-factor signalling, energy charge, Rb
phosphorylation (\code{CycD$>$threshold}), p21 tolerance,
\code{damageLevel} under the DNA damage ceiling, and p53 activity.
\code{evalG2M} demands complete replication (\code{cdogma
.replicationReadyForM()}), low Chk1 signal, adequate
\code{replicationQuality}, and active CycB. \code{evalSAC} is the
spindle assembly checkpoint at metaphase. All three return
\code{\{bool passed; const char* reason;\}} so failure paths emit a
diagnostic string like \texttt{"CHK1\_HIGH"} or \texttt{"CYCD\_LOW"}
to the telemetry log for post-run inspection.

\section{Interaction with the Central Dogma}

The replication arm of the cycle is \emph{delegated} to
\code{CentralDogmaState::updateReplication(...)}: its internal
origin-firing, fork progression, and mismatch-repair logic runs in
parallel with the cycle timer. The cycle reads back
\code{replicationProgress} and \code{replicationReadyForM()} to gate
G2\(\to\)M. See Ch.~\ref{ch:replication} for the full mechanism.

\section{Growth-factor modulation}

\code{envGF} (environmental growth factor, reading
\code{nutrients.get(MS\_GROWTH\_F, x, z)}) directly raises CycD
transcription rate. When the user slides the growth-factor input
from 50\,ng/mL (DMEM + 10\% FBS default) down to 5\,ng/mL
(starvation template), G1 lengthens dramatically and division
frequency collapses. This is the mechanism behind the
\emph{Starvation} preset.

\chapter{Mitosis and Cytokinesis}
\label{ch:mitosis}

Mitosis is the period when the cell is committed to divide. Inside
\code{SimCell::program} sits \code{MitosisState} (declared in
\file{src/simulation/CellCycleProgram.h}), which tracks the fine
phase (\code{MITO\_NONE, PROPHASE, PROMETAPHASE, METAPHASE, ANAPHASE,
TELOPHASE, CYTOKINESIS, COMPLETE}) and a handful of floats:
\code{furrowDepth}, \code{poleSeparation}, \code{chromatinCondense},
\code{dnaDuplicationProgress}, \code{dnaCheckpointPassed},
\code{replicationQuality}, \code{particlesDuplicated}.

\section{Driver: \code{updateMitosisProgram}}

This routine (\file{src/simulation/Simulation.h} line 1275+) decides
when a given cell starts its mitosis program, when to advance the
phase, when the SAC fires, when cytokinesis completes, and when to
commit the division. The key invariant:
\begin{lstlisting}
if (gSimMode == MODE_SINGLE_CELL && isPrimaryCell(c)) return;
\end{lstlisting}
--- only background cells go through the per-cell mitosis program.
The primary focused cell in single-cell mode uses the \emph{visual}
mitosis pipeline inside \file{src/main.mm}, which renders condensed
chromatin, a spindle, chromosome alignment, and a visible furrow.
The per-cell program commits divisions via
\code{Simulation::processDivisions()} after
\code{updateMitosisProgram} runs.

\section{SAC (spindle assembly checkpoint)}

The SAC pauses anaphase entry until every chromosome is bi-oriented.
In the lightweight sim, ``bi-orientation'' is evaluated as follows:
\begin{lstlisting}
static CheckpointResult evalSAC(const MitosisState& m) {
    if (!m.active) return {false, "SAC_NOT_ACTIVE"};
    if (m.replicationQuality < 0.70f) return {false, "REPL_QUAL_LOW"};
    if (!m.dnaCheckpointPassed)       return {false, "DNA_CHK_OPEN"};
    if (m.furrowDepth > 0.05f)         return {false, "FURROW_TOO_EARLY"};
    return {true, "SAC_OK"};
}
\end{lstlisting}
SAC failure is a plausible cause of mitotic slippage and subsequent
apoptosis, especially under \code{TGT\_TUBULIN}-bound drugs like
paclitaxel. The TargetLibrary modulator for tubulin explicitly
inflates \code{c.damageLevel} only while \code{c.phase == 3} (M
phase), emulating SAC-dependent injection of stress.

\section{Division commit}

\code{Simulation::processDivisions()} (around line 850) consumes
\code{DivisionEvent} records queued by \code{updateMitosisProgram}
and rebuilds the daughter cells from an \emph{immutable} parent
snapshot. This is the single place where one cell becomes two:
\begin{lstlisting}
SimCell d = original;  // full copy from snapshot
d.cellUid = allocateCellUid();
d.cloneId = original.cloneId;
d.generation = original.generation + 1;
d.age = 0;
d.biomass *= 0.5f;
d.radius *= 0.87f;         // radius scales as cbrt(0.5)
d.size = 1.0f;             // reset visual scale
// ... apply SplitStats noise
\end{lstlisting}
The \code{SplitStats} struct carries per-split asymmetry:
biomass-split ratio, damage-inheritance partition, telomere clock
decrement, and any dividing bias. This asymmetry is how sibling
cells gradually diverge in fate over generations.

\chapter{Apoptosis}
\label{ch:apoptosis}

\textsc{CellSim} implements a multi-threshold apoptosis engine
(\file{src/simulation/biochem/Apoptosis.h}, 323 lines) whose state
transitions are driven by eleven independent trigger inputs and
whose mass-conservation discipline is baked into the architecture.
Each dying cell moves through seven phases and spills its mass into
the medium through well-defined flux terms.

\section{Phase machine}

\begin{lstlisting}
enum Phase : uint8_t {
    ALIVE         = 0,   // baseline
    PRIMED        = 1,   // pro-apoptotic pressure is building
    MOMP          = 2,   // mitochondrial outer membrane permeabilised
    EXECUTION     = 3,   // caspase-3 active, PS flip, shrinkage
    FRAGMENTATION = 4,   // cell breaks into bodies
    BODIES        = 5,   // drifting apoptotic bodies
    CLEARED       = 6    // engulfed / decomposed
};
\end{lstlisting}

Literature-anchored dwell times:
\begin{itemize}
  \item MOMP $\to$ EXECUTION: 300\,\biosec (5\,min) --- Ye 2017.
  \item EXECUTION $\to$ FRAGMENTATION: 1\,800\,\biosec (30\,min) --- Spencer 2009.
  \item FRAGMENTATION $\to$ BODIES: 1\,800\,\biosec (30\,min).
  \item BODIES $\to$ secondary necrosis: 6\,\bioh post-MOMP --- Silva 2010.
\end{itemize}

\section{Eleven input triggers}

\code{ApoTriggers} is populated each tick by
\code{Simulation::updateApoptosis}:
\begin{enumerate}
  \item \textbf{p53 DNA damage} $\gets$ \code{c.damageLevel}.
  \item \textbf{ROS} $\gets$ \code{c.ROS}.
  \item \textbf{ATP collapse} $\gets$ \code{c.ATP < thresh}.
  \item \textbf{Hypoxia} $\gets$ \code{c.hypoxiaIntensity > 0.5}.
  \item \textbf{Survival-factor loss} $\gets$ low growth-factor.
  \item \textbf{Anoikis} $\gets$ low \code{adhesionStrength}.
  \item \textbf{Crowding} $\gets$ \code{chronicPressureBioSec}.
  \item \textbf{Replicative} $\gets$ telomere / senescence.
  \item \textbf{Fas ligand} $\gets$ \code{fasLExposure}.
  \item \textbf{Mitochondrial dysfunction} $\gets$ \code{mitoPotential < 0.4}.
  \item \textbf{Drug pressure} $\gets$ cumulative \code{drug\_pro\_apop}.
\end{enumerate}
These feed the Bcl-2 / Bax / Bak / MOMP / caspase 3/8/9 / Smac / IAP
cascade ODE inside \code{ApoptosisEngine::step()}. The engine uses
sub-stepping (\code{DT\_MAX = 0.08\,\biosec}) because the caspase
amplification loop is stiff.

\section{Mass-conserved release}

At each phase transition, a literature-calibrated fraction of
cellular mass is released through \code{MediumField::exchange()}:
\begin{lstlisting}
void Simulation::releaseCytosol(SimCell& c, float fraction) {
    float dB = c.biomass * fraction;
    c.biomass -= dB;
    float flux[MS_COUNT] = {0};
    flux[MS_AA_POOL]  = dB * Apoptosis::REL_AA_PER_BIOMASS;
    flux[MS_IONS]     = dB * Apoptosis::REL_IONS_PER_BIOMASS;
    flux[MS_CALCIUM]  = dB * Apoptosis::REL_CALCIUM_PER_BIOMASS;
    flux[MS_PYRUVATE] = dB * Apoptosis::REL_PYRUVATE_PER_BIOMASS;
    flux[MS_LACTATE]  = dB * Apoptosis::REL_LACTATE_PER_BIOMASS;
    flux[MS_GLUCOSE]  = dB * Apoptosis::REL_GLUCOSE_PER_BIOMASS;
    flux[MS_WATER]    = dB * Apoptosis::REL_WATER_PER_BIOMASS;
    nutrients.exchange(c.position.x, c.position.z, flux, 1.0f);
}
\end{lstlisting}
\code{releaseMembrane} and \code{releaseReceptors} follow the same
pattern on their respective pools. The invariant is strict: mass
out of the cell equals mass into the medium voxel. No mass is ever
created.

\section{Apoptotic bodies in the rendered mode}

\file{src/main.mm} renders the apoptotic body swarm: each dying cell
spawns 5--15 bodies (\code{gApoBodies}) with partitioned mass ledgers
(cytosol, membrane, receptor). Bodies drift with the fluid current,
swell through osmotic imbalance, and burst at 6\,\bioh via the
nanopore-driven osmotic model of Chen 2016. Live neighbouring cells
within 4\,\um (efferocytosis radius) engulf bodies into lysosomal
pools at $85\%$ recycling efficiency (Ravichandran 2015,
Appelqvist 2013).

% ═════════════════════════════════════════════════════════════════════
\part{The Central Dogma}
% ═════════════════════════════════════════════════════════════════════

\chapter{DNA Replication}
\label{ch:replication}

The \code{CentralDogmaState} struct (\file{src/simulation/CentralDogma.h}
line 457, post-Phase 8A edits) holds a complete origin--fork--
mismatch pipeline for the HBB host gene. One S-phase produces one
replicated HBB copy per cell; a daughter inherits half.

\section{Gene geometry}

\code{HBB\_SEQUENCE[]} is 1558\,bp of real human beta-globin
(Ensembl GRCh38 chr11:5\,225\,464--5\,227\,021, minus strand;
NCBI Gene ID 3043):
\begin{itemize}
  \item Exon 1: 0--89 (90\,bp) --- 5' UTR + start of CDS.
  \item Intron 1: 90--218 (129\,bp).
  \item Exon 2: 219--441 (223\,bp).
  \item Intron 2: 442--1293 (852\,bp).
  \item Exon 3: 1294--1557 (264\,bp) --- includes 3' UTR.
\end{itemize}
The engine stores the raw sequence as a literal so \code{transcribe
(HBB\_SEQUENCE[bp])} is a pointer-to-literal access (zero allocation,
zero copying).

\section{Origin firing}

Three origins are hard-coded at $L/4$, $L/2$, $3L/4$. Origin 0 is
licensed and marked non-dormant; origins 1, 2 are licensed dormant
and fire only if the primary origin stalls:
\begin{lstlisting}
replicationOrigins[0].baseIndex = HBB_LENGTH / 2;
replicationOrigins[0].licensed = true;
replicationOrigins[0].dormant = false;

replicationOrigins[1].baseIndex = HBB_LENGTH / 4;
replicationOrigins[1].licensed = true;
replicationOrigins[1].dormant = true;

replicationOrigins[2].baseIndex = (HBB_LENGTH * 3) / 4;
replicationOrigins[2].licensed = true;
replicationOrigins[2].dormant = true;
\end{lstlisting}
On fire, each origin produces two oppositely-directed replication
forks. Fork velocity is capped at 24\,bp/s
(\code{CDOGMA\_LIT\_FORK\_SPEED\_BP\_PER\_SEC}), the mammalian
single-molecule DNA-combing rate.

\section{Mismatch pipeline}

Each fork integrates a per-base mismatch probability that combines
four factors: a homopolymer run-length amplifier, a local GC fraction
penalty, a CpG transition hotspot, and a sequence-specific risk
factor. On error, the algorithm first tries proofreading
(Pol\(\epsilon\) exonuclease, \(\sim\)90\% efficient), then mismatch
repair (MMR, \(\sim\)95\% efficient on the remaining errors), and
finally promotes any residual error to a \emph{replication
mismatch state} that persists until post-replicative repair.

The telemetry counters \code{proofreadCorrections},
\code{mmrCorrections}, \code{escapedErrors} are surfaced in the
diagnostic log each frame for post-run analysis.

\section{Checkpoint dynamics}

\code{chk1Signal} rises with \code{replicationStress} (fork-stall
intensity averaged across active forks plus penalties for unresolved
/ escaped errors). The G2/M checkpoint demands
$\code{chk1Signal}< 0.35$ and $\code{replicationQuality} > 0.82$. A
heavy drug load that stalls forks (e.g.\ doxorubicin binding
\code{TGT\_TOPO2}) drives \code{chk1Signal} up, holding the cell in
G2.

\chapter{Transcription}
\label{ch:transcription}

Three RNA Pol II slots (\code{TranscriptionState[3]}) run in
parallel. The engine is fully concurrent: each slot independently
loads an idle timer, fires \code{startTranscription()} when allowed,
and elongates at $\sim$0.22\,s$^{-1}$ (progress fraction, corresponding
to $\sim$340\,bp/s --- the literature-consistent elongation rate
shared by major eukaryote-wide studies, Singh 2013).

\section{Phase 8A --- multiple gene templates}

Until Phase 8A (2026-04-20) the engine transcribed only
\code{HBB\_SEQUENCE}. On that date a \code{GeneTemplate} registry was
added; Phase 8A is dissected in full in Ch.~\ref{ch:phase8a}.
Here it suffices to note the call graph:

\begin{lstlisting}
int pickGeneForTranscription();   // promoter-weighted choice
int registerViralGene(const GeneTemplate&);  // append & build mRNA
\end{lstlisting}

Elongation now reads via:
\begin{lstlisting}
int seqLen;
const char* seqData = geneSequence(tx->geneIndex, seqLen);
while (tx->genomicBases < ts.currentBP && tx->genomicBases < seqLen) {
    tx->preMRNA.push_back(transcribe(seqData[tx->genomicBases]));
    tx->genomicBases++;
}
\end{lstlisting}
When \code{geneIndex == -1} (the host default) this is a direct read
of \code{HBB\_SEQUENCE}. When $\ge 0$ it reads from the string
storage of an added viral gene. HBB's fast path is unaffected.

\section{Capping, splicing, polyadenylation}

A transcript is \emph{capped} once its length crosses 18\,nt
(biological: 7mG cap added by capping enzyme during early
elongation). Once elongation completes, splicing begins --- but
\emph{only} for host genes. Viral genes, per Phase 8A, skip splicing
outright: they jump to \code{spliced = true, matureMRNA =
viralGenes[idx].codingMRNA}. This mirrors the biology of
single-exon RNA viruses (flu, most coronaviruses).

Host splicing binds one of two spliceosome slots
(\code{SplicingState[2]}) at a rate $0.45\,\text{s}^{-1}$. Intron 1
is removed at $25\%$ spliceosome progress, intron 2 at $70\%$, and
polyadenylation tail ($48 \cdot \text{A}$) attaches at $100\%$.

\section{Nuclear export}

An exported flag advances at rate $0.60\,\text{s}^{-1}$ once the
transcript is spliced. Only exported mRNAs are eligible for
translation (enforced in \code{findTranscriptForTranslation}).

\chapter{Translation}
\label{ch:translation}

Six ribosome slots (\code{TranslationState[6]}) compete for exported
mature mRNAs through \code{findTranscriptForTranslation()}. This is
the point where Phase 8A's viral genes compete with HBB for
ribosome capacity --- no explicit coupling was needed because
\code{findTranscriptForTranslation} simply returns the first eligible
transcript and the code handles multiple transcripts transparently.

\section{Initiation}

\code{startTranslation(ribosome, tx)} clamps the ribosome to the
transcript, ticks \code{tx->ribosomeLoad++}, and calls
\code{primeCurrentCodon} which reads three nucleotides from the
mRNA starting at either \code{tx->transcriptStartCodon} (Phase 8A,
per-gene), falling back to the cell-wide \code{startCodonBase} for
host HBB.

\section{Elongation}

Each codon cycle:
\begin{enumerate}
  \item Ribosome reads the codon, computes the anticodon, looks up
        the charged tRNA (from \code{trnaPool[12]}).
  \item tRNA shuttles at rate \(8\,\text{s}^{-1}\) (mammalian in vivo
        translation is $\sim$6\,aa/s, allowing for wobble and wait).
  \item On shuttle completion, the amino acid letter is appended to
        \code{nascentPeptide} and the ribosome advances one codon.
  \item When \code{primeCurrentCodon} finds a stop codon the peptide
        is released via \code{releaseProtein()}.
\end{enumerate}

\section{Protein release (Phase 8A-aware)}

\begin{lstlisting}
void releaseProtein(const TranslationState& tr) {
    int slot = findFreeProteinSlot();
    ProteinProductState& p = proteins[slot];
    p.active = true;
    p.aminoAcids = tr.nascentPeptide;
    const TranscriptState* tx = getTranscript(tr.transcriptIndex);
    if (tx) {
        p.geneIndex   = tx->geneIndex;
        p.geneId      = tx->geneId;
        p.proteinName = tx->proteinTag;
        if (p.geneIndex >= 0) {
            viralProteinCount[tx->proteinTag]++;
            viralGenes[p.geneIndex].copiesMade++;
        }
    }
    if (tr.nascentPeptide.size() >= 120 && tr.nascentPeptide.size() <= 170)
        p.structureAsset = "hbb_folded.pdb";
}
\end{lstlisting}
Every protein carries its gene origin. Viral proteins also
increment \code{viralProteinCount[proteinTag]} --- the ledger that
Phase 8E (stoichiometric assembly) consumes.

\section{Folding and PTMs}

Released proteins fold at rate $0.30\,\text{s}^{-1}$ and mature
through post-translational modification at $0.18\,\text{s}^{-1}$.
The folding step is where chaperone availability (BiP, HSP70 --- see
\file{src/simulation/biochem/Proteome.h}) can throttle throughput.
Under ER stress (e.g.\ oxidative damage) \code{ProteomeEngine} raises
HSF1, recruiting more chaperones, at the cost of translation
initiation efficiency (the classical UPR trade-off).

% ═════════════════════════════════════════════════════════════════════
\part{Cellular Chemistry}
% ═════════════════════════════════════════════════════════════════════

\chapter{Metabolism}
\label{ch:metabolism}

\file{src/simulation/biochem/Metabolism.h} (962 lines) implements the
cell's energy and biosynthesis arithmetic, anchoring on a
pool/flux representation compatible with Recon3D-style flux balance
analysis without requiring an LP solver at runtime.

\section{Species pool}

The pool holds:
\begin{itemize}
  \item Adenylate: ATP, ADP, AMP (all in mM; energy charge = $(\text{ATP} + 0.5\,\text{ADP})/(\text{ATP}+\text{ADP}+\text{AMP})$).
  \item NAD/NADH and FAD/FADH\(_2\) redox pairs.
  \item Central-carbon intermediates: glucose, G6P, F6P, F1,6BP, 3PG, pyruvate, lactate, citrate, oxaloacetate, $\alpha$-ketoglutarate, succinate, fumarate, malate.
  \item Amino-acid pool.
  \item O\(_2\) consumption rate per tick.
\end{itemize}

\section{Flux computation}

Each flux (glycolysis, lactate dehydrogenase, pyruvate dehydrogenase,
TCA cycle, oxidative phosphorylation, pentose phosphate shunt,
glutaminolysis) is computed from the pool concentrations using
Michaelis--Menten rate laws parametrised with literature $K_m$ and
$V_\text{max}$. The headline constants come from Cornish-Bowden
2012 and BRENDA. No learned model is used; the fluxes emerge from
pools + physical-chemistry rate laws.

\section{Coupling to cycle}

ATP availability (\code{pool.ATP}) is piped back into
\code{SimCell::ATP} through a scaling factor
\code{bioATPToSimPercent()} so older code paths continue to work
with ATP-as-percentage. The cycle engine consults \code{c.ATP}
through the \code{evalG1S} energy gate: below $\sim$50\%, G1/S is
blocked; below $\sim$25\%, the energy-stress trigger feeds apoptosis.

\section{Warburg switch}

When \code{mitoHealth} drops below \code{MITO\_HEALTH\_SWITCH = 0.4}
the cell flips \code{glycolytic = true}, accelerating lactate
production by \code{WARBURG\_GLY\_BOOST = 1.6\(\times\)} while
penalising oxphos by \code{WARBURG\_OXP\_PENALTY = 0.5\(\times\)}.
The effect is visible in the medium log as rising lactate and
acidifying pH.

\chapter{Membrane Transport and Ion Physics}
\label{ch:membrane}

\file{src/simulation/biochem/MembraneTransport.h} holds ion
concentrations inside and out, electrochemical potentials, pumps,
voltage-gated channels, and a Hodgkin--Huxley model for action
potentials (dormant in the default simulation but available for
excitable-cell scenarios).

\section{Ion pools}

Intracellular / extracellular (mM):
\begin{center}
\begin{tabular}{lrrrrr}
\toprule
         & Na$^+$ & K$^+$ & Ca$^{2+}$ & Cl$^-$ & Mg$^{2+}$\\ \midrule
In       & 12     & 140   & $10^{-4}$ & 4      & 0.5\\
Out      & 145    &  5    & 1.8       & 120    & (same)\\
\bottomrule
\end{tabular}
\end{center}

\section{Nernst and Goldman}

Reversal potentials per ion:
\begin{equation}
  E_X = \frac{RT}{zF}\,\ln\!\left(\frac{[X]_\text{out}}{[X]_\text{in}}\right),
  \qquad \frac{RT}{F}\bigg|_{37^\circ\text{C}} = 26.7\,\text{mV}.
\end{equation}
Resting membrane potential via the Goldman equation:
\begin{equation}
  V_m = \frac{RT}{F} \ln\!
    \frac{P_\text{Na}[\text{Na}]_o + P_\text{K}[\text{K}]_o + P_\text{Cl}[\text{Cl}]_i}
         {P_\text{Na}[\text{Na}]_i + P_\text{K}[\text{K}]_i + P_\text{Cl}[\text{Cl}]_o}
\end{equation}
with permeability ratios $P_\text{Na}:P_\text{K}:P_\text{Cl} = 0.04:1.0:0.45$.

\section{V-ATPase}

The vacuolar-type H$^+$ ATPase pumps protons against the gradient,
consuming one ATP per approximately three H$^+$ translocated. In the
current engine it is implicit in the vesicular-traffic compartment
pH gradients (Ch.~\ref{ch:vesicular}). Phase 8D of the plan
(Ch.~\ref{ch:phase8a}) promotes it to an explicit per-endosome
acidifier.

\chapter{Vesicular Traffic}
\label{ch:vesicular}

\file{src/simulation/biochem/VesicularTraffic.h} (281 lines) carries
the secretory and endocytic pathways as an explicit pipeline. Eleven
compartments form a DAG whose edges are \code{Vesicle} objects
labelled with a coat type (COPII, COPI, clathrin, uncoated), a cargo
class (\code{CARGO\_SECRETORY}, etc.), a progress fraction, and a
Rab GTPase id for targeting.

\section{Compartments}

\code{enum Compartment} lists ER, cis/medial/trans Golgi, TGN, early
endosome, late endosome, lysosome, plasma membrane, cytosol, nucleus.
Each carries a luminal pH and a protein load.

\begin{center}
\begin{tabular}{lr}
\toprule
Compartment & pH \\ \midrule
ER                  & 7.2 \\
cis-Golgi           & 6.7 \\
medial-Golgi        & 6.3 \\
trans-Golgi/TGN     & 6.0 \\
Early endosome      & 6.0 \\
Late endosome       & 5.5 \\
Lysosome            & 4.8 \\
Cytosol             & 7.2 \\
Nucleus             & 7.2 \\
\bottomrule
\end{tabular}
\end{center}

This pH gradient is load-bearing for Phase 8D: the endosomal
acidification window is exactly where flu HA exposes its fusion
peptide. Without the gradient, the entry machinery cannot fire.

\section{Coat and SNARE pools}

Four coat proteins are tracked as pools
(\code{COPII\_available, COPI\_available, clathrin\_available,
dynamin}) and three SNARE proteins
(\code{syntaxin\_PM, SNAP25, VAMP}), plus \code{NSF} for SNARE
recycling. Phase 8C's clathrin-coated pit model decrements
\code{clathrin\_available} per pit formation and returns it after
scission.

\section{UPR and autophagy}

\code{UPRState} holds unfolded-protein-response arms (IRE1, PERK,
ATF6) and the XBP1s / eIF2$\alpha$-P switches. \code{AutophagyState}
holds mTORC1 inhibition, ULK1, Beclin1, LC3-I/II, p62, and the
autophagosome flux. Both feed back into apoptosis (prolonged CHOP
raises the pro-apop trigger) and metabolism (ULK1 activation diverts
AA to recycling).

\chapter{Cytoskeleton}
\label{ch:cytoskeleton}

\file{src/simulation/biochem/Cytoskeleton.h} holds actin,
microtubule, and intermediate-filament state, plus myosin motors and
RhoA/Rac1/Cdc42 GTPases. Engine-level hooks: cortical tension
(pN/\um) and contractile force feed cell mechanics; myosin-V/VI
contribute to vesicle motility.

\section{Actin dynamics}

\code{F\_actin} and \code{G\_actin} pools; tropomyosin stabilises F;
cofilin severs filaments (rate $\propto$ ADP-actin fraction); Arp2/3
nucleates branched networks (driven by active \code{WASP}). The
balance of these rates, driven by the GTPase signalling network,
yields cortical thickness and motility speed.

\section{Relevance to pathogens}

Part 8 §56 of the plan calls out the T3SS mechanism as an explicit
path where a bacterial effector protein (e.g.\ \emph{Salmonella} SopE)
binds host Rho GTPases and triggers actin ruffling. This reuses
\code{Cytoskeleton.motors} directly: SopE's BindingMatcher score
against the RhoA pharmacophore drives a transient upregulation of
Arp2/3, producing visible membrane ruffles in the focused cell.
Phase 8H of the plan wires this explicitly.

% ═════════════════════════════════════════════════════════════════════
\part{Dish-Level Chemistry}
% ═════════════════════════════════════════════════════════════════════

\chapter{MediumField --- the 2D diffusion grid}
\label{ch:medium}

\file{src/simulation/MediumField.h} (337 lines) is a finite-volume
diffusion PDE solver on a 64$\times$64 voxel grid. It holds twelve
scalar species plus two auxiliary scalar fields and one vector-of-
fields for dynamic bioagent concentrations.

\section{Species enumeration}

\begin{lstlisting}
enum MediumSpecies {
    MS_O2=0, MS_CO2=1, MS_GLUCOSE=2, MS_GLUTAMINE=3,
    MS_PYRUVATE=4, MS_LACTATE=5, MS_AA_POOL=6,
    MS_GROWTH_F=7, MS_IONS=8, MS_CALCIUM=9,
    MS_HPLUS=10, MS_WATER=11, MS_COUNT=12
};
\end{lstlisting}
Units: mM for most species; ng/mL for \code{MS\_GROWTH\_F}; pH
(stored directly) for \code{MS\_HPLUS}. Water is pinned near
$55\,000\,\text{mM}$ except where ions shift osmolarity.

\section{Diffusion coefficients}

Species-specific coefficients (in \wu$^2$/s) are taken from standard
tables (Crank 1975) and scaled by the world-to-micrometer ratio:
\begin{center}
\begin{tabular}{lr}
\toprule
Species & $D\,(\wu^2/\text{s})$\\ \midrule
O$_2$        & 1.80 \\
CO$_2$       & 1.50 \\
Glucose      & 0.90 \\
Lactate      & 1.20 \\
H$^+$ (pH)   & 0.60 \\
AA pool      & 0.80 \\
Ions         & 1.40 \\
\bottomrule
\end{tabular}
\end{center}

\section{Finite-volume step}

One pass of \code{MediumField::diffuse(dt\_s, envO2, envGlu)}:
\begin{enumerate}
  \item For each species, compute the 5-point Laplacian on the
        current grid into \code{cNext[\,]}.
  \item Apply reaction / boundary-exchange terms.
  \item Swap \code{c} $\leftrightarrow$ \code{cNext}.
\end{enumerate}

The auxiliary \code{extracellularProteaseLevel[CELLS]} field models
lytic enzymes released by dying cells; it decays with a fixed
\texttt{keep} factor each tick, modelling degradation.

\section{Mass-balance invariant}

\code{totalAtStart[MS\_COUNT]} stores the initial sum of each species.
\code{checkBalance()} compares the current sum to the initial;
deviation under $10^{-3}$ is considered on-budget. Every
\code{exchange(x, z, flux, scale)} call is paired with its source
pool, so over a full 46\,\bioh run the drift stays tight even with
dozens of apoptosis lysis events.

% ═════════════════════════════════════════════════════════════════════
\part{Intracellular Physics}
% ═════════════════════════════════════════════════════════════════════

\chapter{Particles, Attraction Fields, Langevin Dynamics}
\label{ch:intraphysics}

\file{src/simulation/IntracellularPhysics.h} (514 lines) holds the
discrete particle simulator for the focused primary cell in
single-cell mode. It is a lightweight Langevin integrator with
Hertz contact, confinement wall, attraction fields, and job state
machines.

\section{Particle types}

\begin{lstlisting}
enum ParticleType {
    PT_GLUCOSE, PT_ATP, PT_MITO, PT_ER_ROUGH, PT_ER_SMOOTH,
    PT_GOLGI, PT_NUCLEUS, PT_LYSOSOME, PT_DNA_A, PT_DNA_B,
    PT_HISTONE, PT_RIBOSOME_SMALL, PT_RIBOSOME_LARGE,
    PT_TRNA, PT_POLYPEPTIDE, PT_SPLICEOSOME, PT_NUCLEAR_PORE,
    PT_RECEPTOR_GPCR, PT_RECEPTOR_RTK, PT_RECEPTOR_CYTOKINE,
    PT_RECEPTOR_DEATH, PT_RECEPTOR_TLR, PT_RECEPTOR_NOTCH,
    PT_RECEPTOR_FRIZZLED, PT_RECEPTOR_PATCHED, PT_RECEPTOR_STK,
    PT_ION_CHANNEL, PT_PUMP, PT_CHAPERONE,
    PT_VESICLE_COPII, PT_VESICLE_SECRETORY, PT_MOLECULE
};
\end{lstlisting}

Each type has a visual radius, a diffusion coefficient (set from
Stokes--Einstein at 37$^\circ$C given the type's hydrodynamic radius),
and a default attraction target.

\section{Langevin step}

\code{IntracellularPhysics::step(dt)} performs, per particle:
\begin{equation}
  \mathbf{v}_{t+1} = (1-\gamma\,dt)\,\mathbf{v}_t + \frac{\mathbf{F}\,dt}{m} + \sqrt{2\gamma k_B T\,dt/m}\,\boldsymbol{\xi}_t
\end{equation}
with Gaussian $\boldsymbol{\xi}_t$. The force $\mathbf{F}$ sums:
\begin{itemize}
  \item Hertz repulsion at membrane confinement radius.
  \item Tag-filtered attraction (\code{AttractionField}).
  \item Hydrodynamic drag (built into $\gamma$).
\end{itemize}

Attraction fields are a single scalar-strength attractor with a tag
filter (\code{attractedTag}, \code{attractedType}) so only molecules
of type $X$ feel pull toward target of type $Y$. This is the
mechanism by which glucose molecules drift to the hexokinase
pocket, or a drug particle drifts to its bound target.

\section{Job state machine}

Every particle has a \code{job} field in
\{\code{JOB\_IDLE, JOB\_GOTO\_TARGET, JOB\_ATTACH,
JOB\_CONSUMED}\}. A particle with \code{JOB\_GOTO\_TARGET} has an
\code{attachTarget} pointer and an attraction override that pulls
it toward the target's position. On arrival (distance $<$ attach
radius), job flips to \code{JOB\_ATTACH}. After the target's
reaction dwell, \code{JOB\_CONSUMED} marks the particle for garbage
collection. This is how enzymatic reactions are animated without a
separate chemistry engine.

% ═════════════════════════════════════════════════════════════════════
\part{Chemistry-First Drug System}
% ═════════════════════════════════════════════════════════════════════

\chapter{ChemicalEntity --- the universal structure record}
\label{ch:chementity}

\file{src/simulation/biochem/ChemicalEntity.h} (139 lines) defines
the universal record for anything with a chemical structure ---
drug, virion, bacterium, metabolite, organelle, receptor, membrane
lipid, nucleic acid. One type, many kinds.

\section{Schema}

The \code{ChemicalEntity} struct holds identity
(\code{id}, \code{name}, \code{kind}), structure
(\code{mol}, \code{protein}, \code{pharmacophores}), Tier 0--2
chemistry (\code{mw}, \code{logP}, \code{logS}, \code{tpsa},
\code{hbd}, \code{hba}, \code{rotatable\_bonds},
\code{aromatic\_rings}, \code{formal\_charge},
\code{polarizability}, \code{homo\_eV}, \code{lumo\_eV},
\code{dipole\_debye}), Tier 1 force-field parameters
(\code{partialCharges}, \code{ljSigma\_nm}, \code{ljEps\_kJ\_per\_mol}),
an optional Tier 1b induced-dipole vector, a 2\,048-bit Morgan
fingerprint packed into 256 bytes, a list of
\code{BindingAffinity} entries keyed by other-entity id, a list of
reaction rules the entity participates in, and an entry pathway
enum.

\section{Kinds}

\begin{lstlisting}
enum class ChemicalEntityKind : uint8_t {
    UNKNOWN, METABOLITE, DRUG, VIRUS, BACTERIUM,
    ORGANELLE, ENZYME, RECEPTOR, MEMBRANE_LIPID,
    NUCLEIC_ACID, COFACTOR, PEPTIDE, ANTIBODY, ION
};
\end{lstlisting}
Kind selects which fields are populated; the physics / reaction
engines dispatch on \code{ChemicalEntity} without branching on
kind.

\chapter{BindingMatcher --- Tier 0 descriptor-based affinity}
\label{ch:bmatcher}

\file{src/simulation/biochem/BindingMatcher.h} (91 lines) computes
drug~$\times$~target $K_d$ from a 5-D descriptor compatibility
score. The math is a Hill-like Gaussian fit.

\section{Score construction}

\begin{equation}
  \text{score} = 0.7 \cdot \left(\prod_{i=1}^5 f_i\right)^{1/5} + 0.3 \cdot \text{fpScore}
\end{equation}
where $f_i$ is the Gaussian fit of drug descriptor $i$ against the
target's preferred range; \texttt{fpScore} is a fingerprint Tanimoto
(placeholder 0.5 at Tier 0; populated by a real Morgan fingerprint
at Tier 3).

\section{Kd mapping}

\begin{equation}
  K_d\,[\text{mM}] = 10^{2 - 8\,\text{score}}
\end{equation}
so \emph{perfect} match (score = 1) yields $K_d = 10^{-6}\,\text{mM}
= 1\,\text{nM}$; a moderate match (score = 0.5) yields
$10\,\um\text{M}$; no match (score = 0) yields 100\,mM (effectively
no binding).

\section{Kinetic parameters}

\begin{align*}
  k_\text{on} &= 10^{-2}\,\text{per}\,\um\text{M}\cdot\text{s} \approx 10^7\,\text{M}^{-1}\text{s}^{-1}\\
  k_\text{off} &= k_\text{on} \cdot K_d
\end{align*}
capping \(k_\text{on}\) below the Berg--Purcell diffusion limit of
$10^8\,\text{M}^{-1}\text{s}^{-1}$.

\chapter{TargetLibrary --- the per-target modulator bank}
\label{ch:targetlib}

\file{src/simulation/biochem/TargetLibrary.h/.cpp} registers thirteen
built-in cellular targets and the function-modulator each target
runs when its occupancy is non-zero. Each modulator only
perturbs an existing \code{SimCell} field --- never adds a new
outcome code path.

\begin{center}
\begin{tabular}{ll}
\toprule
Target & What the modulator does\\ \midrule
\code{TGT\_CDK1}   & scales \code{cdk.CycB} $\to$ M-entry arrest\\
\code{TGT\_CDK2}   & scales \code{cdk.CycE, CycA} $\to$ S-entry block\\
\code{TGT\_RB}     & holds \code{Rb} hypophos $+$ raises \code{p21}\\
\code{TGT\_TOPO2}  & stalls \code{replicationActive}, raises \code{damageLevel}, \code{ROS}, \code{p53}\\
\code{TGT\_TUBULIN}& only in M-phase: raises \code{damageLevel} (SAC fail)\\
\code{TGT\_BCL2}   & scales \code{apo.Bcl2, BclXL} down\\
\code{TGT\_DNA\_INTERCALATION} & raises \code{damageLevel, p53, Puma}; direct \code{Bax} push at saturation\\
\code{TGT\_HDAC}   & raises \code{stress}\\
\code{TGT\_KINASE\_EGFR} & scales \code{cdk.CycD} down\\
\code{TGT\_PROTEASOME} & raises \code{stress}, \code{Puma}\\
\code{TGT\_RIBOSOME} & scales \code{biomass}\\
\code{TGT\_POL2}    & scales \code{cdk.CycD, CycE} down\\
\code{TGT\_MITO\_COMPLEX\_I} & lowers \code{mitoPotential}, raises \code{ROS}\\
\bottomrule
\end{tabular}
\end{center}

The discipline that \emph{no modulator adds a new biological
outcome} is what makes drug effects emergent. The cell already
knows how to arrest, die, slow, or activate based on these fields;
drugs only perturb the fields.

\chapter{Pharmacokinetics --- drug entry and target binding}
\label{ch:pk}

\code{Simulation::updateDrugPK(float dt)} runs the PK/PD pipeline
once per substep. For every applied drug:
\begin{enumerate}
  \item Compute Lipinski-derived permeability
        \(p = 0.02 \cdot e^{0.3 (\text{logP}-1)} \cdot e^{-\text{mw}/700} \cdot e^{-\text{tpsa}/150}\),
        capped in $[10^{-5}, 0.2]$.
  \item Fick flux from dish concentration to cytoplasm:
        \(\Delta[\text{drug}] = p \cdot (C_\text{dish} - C_\text{cyto}) \cdot dt_\text{bio}\).
  \item For each target in the TargetLibrary, mass-action binding:
        \(d[\text{bound}]/dt = k_\text{on}\,[\text{drug}]_\text{free}\,[\text{target}]_\text{free} - k_\text{off}\,[\text{bound}]\).
  \item Target \emph{copies} convert to \um M using the Avogadro +
        cell-volume constant $2.41\cdot 10^{6}$ (from
        $N_A \cdot V \cdot 10^{-6}$ at $V = 4\,\text{pL}$ HeLa volume).
  \item Occupancy = bound / (bound + free) drives the target's
        modulator.
\end{enumerate}

\section{Why the cell volume matters}

An earlier version of this code used $N_A \cdot 10^{-6} = 6 \cdot
10^{17}$, which is six orders of magnitude off: it would require
that one cell contain a mole of receptors. The corrected factor
$2.41 \cdot 10^6$ accounts for the picolitre cell volume. This is
the kind of constant that a reviewer would catch in a week; having
it wrong turned every drug into either completely inert or a single-
molecule killer.

% ═════════════════════════════════════════════════════════════════════
\part{Pathogens}
% ═════════════════════════════════════════════════════════════════════

\chapter{Virion Biology}
\label{ch:virion}

\file{src/simulation/biochem/Virion.h} (111 lines) defines the
\code{VirionSpec} catalogue entry and the per-instance
\code{Virion} record used by the physics loop.

\section{VirionSpec --- the catalogue}

\begin{lstlisting}
struct VirionSpec {
    std::string  id;
    VirionShape  shape;        // ICOSAHEDRAL, HELICAL, COMPLEX, PLEOMORPHIC
    int          T_number;     // Caspar-Klug tile count
    bool         enveloped;
    float        radius_nm;
    float        hydrodynamicRadius_nm;
    GenomeKind   genomeKind;   // SSRNA+, SSRNA-, DSRNA, SSDNA, DSDNA, RETRO
    int          genomeLength_nt;
    int          spikesPerVirion;
    // Spike 5-D descriptor for BindingMatcher
    float spike_logP, spike_mw;
    int   spike_hbd, spike_hba, spike_aromatic;
    // Kinetics
    float uncoatDwell_bioSec;
    float replicationRate_per_s;
    float assemblyThreshold;
    float budRate_per_s;
    float lyticYield;
    std::vector<std::string> preferredReceptors;
    float cytotoxicity_per_copy;
};
\end{lstlisting}

Three toy entries are currently shipped:
\begin{itemize}
  \item \texttt{flu\_h1n1}: enveloped, pleomorphic, 300 spikes,
    ssRNA-, $13.5$\,knt, binds glycoprotein + RTK.
  \item \texttt{hbv\_01}: enveloped, icosahedral $T{=}4$, 90 spikes,
    dsDNA, $3.2$\,knt, binds RTK.
  \item \texttt{adeno\_tox}: non-enveloped, icosahedral $T{=}25$,
    dsDNA, lytic, binds RTK + glycoprotein.
\end{itemize}

All three are \emph{toy}: the rates are rule-of-thumb, not
literature-calibrated. Part 8 of the plan replaces the
autocatalytic \code{replicationRate\_per\_s} with the real
Central Dogma pipeline.

\section{Per-instance state}

\begin{lstlisting}
enum class VirionStage : uint8_t {
    FREE, BOUND, ENTERING, UNCOATING,
    REPLICATING, ASSEMBLING, BUDDING,
    SPILLED, CLEARED
};

struct Virion {
    int specIdx;
    simd_float3 position, velocity;
    VirionStage stage;
    int  hostCellIdx;
    float stageTimer_s;
    float genomeCopies;
    float assembled;
    float boundSpikes;
    uint32_t uid;
};
\end{lstlisting}

\chapter{Bacterium Biology}
\label{ch:bacterium}

\file{src/simulation/biochem/Bacterium.h} (97 lines) follows the same
pattern for prokaryotes. Shape, gram-type, geometry, flagella and
pili counts, adhesin descriptor, doubling time, toxin-secretion
rates. The per-instance \code{Bacterium} carries position, velocity,
stage, stageTimer, divisionClock, biomass, tumbleTimer,
run-direction, uid.

Three catalogue entries:
\begin{itemize}
  \item \texttt{ecoli\_K12}: gram-negative rod, 2\,\um, 6 flagella,
    commensal (no toxin).
  \item \texttt{listeria\_EGDe}: gram-positive rod, 1.5\,\um,
    moderate toxin, binds RTK + integrin.
  \item \texttt{saureus\_tox}: gram-positive coccus, 1\,\um,
    capsule, heavy toxin (pore-forming).
\end{itemize}

Run-and-tumble motility follows Berg 1972: tumble once per bio-s on
average, choose a random new direction, run straight at
0.6\,\wu/\biosec.

\chapter{PathogenKinetics and the Infection Chain}
\label{ch:pathkinetics}

\file{src/simulation/biochem/PathogenKinetics.h} (123 lines) is pure
scoring functions. \file{src/simulation/biochem/PathogenRegistry.h}
(210 lines) is the YAML loader and spec catalogue singleton.

\section{Seven receptor-class reference profiles}

The binding engine knows seven receptor classes and their
canonical 5-D descriptor profile:

\begin{center}
\begin{tabular}{lrrrrr}
\toprule
Id & logP & mw & hbd & hba & arom\\ \midrule
\code{PT\_RECEPTOR\_GPCR}     & 2.5 & 45000 & 5 & 12 & 3\\
\code{PT\_RECEPTOR\_RTK}      & 2.0 & 70000 & 6 & 14 & 3\\
\code{PT\_RECEPTOR\_DEATH}    & 1.2 & 50000 & 7 & 10 & 1\\
\code{PT\_RECEPTOR\_TLR}      & 1.8 & 95000 & 9 & 18 & 4\\
\code{PT\_RECEPTOR\_CYTOKINE} & 1.6 & 55000 & 6 & 12 & 2\\
\code{PT\_RECEPTOR\_INTEGRIN} & 2.2 &110000 & 10& 20 & 3\\
\code{PT\_GLYCOPROTEIN}       & 0.5 & 20000 & 12& 18 & 0\\
\bottomrule
\end{tabular}
\end{center}

\section{Scoring}

\code{compatScore5D} returns
$\prod_i f_i^{1/5}$ where each $f_i$ is a Gaussian fit. The
receptor class the virion (or bacterium) matches best determines
which IntraParticle receptor it will bind. Deleting or scrambling
the virion's descriptor collapses every score to near zero --- the
negative control test in \file{tests/headless\_infection.cpp}
exercises exactly this.

\section{On-rate mapping}

\begin{equation}
  k_\text{eff} = 10^{3\,(\text{score} - 0.7)}
\end{equation}
so score = 0.7 gives $k_\text{eff} = 1\,\text{s}^{-1}$; score = 1
gives $10\,\text{s}^{-1}$; score = 0.3 gives $10^{-1.2} \approx
0.06\,\text{s}^{-1}$. Combined with the hard \texttt{score < 0.35}
gate, this keeps unrelated virions from binding.

\section{The nine-stage chain}

Per-tick loop inside \code{Simulation::updatePathogens(dt)}:
\begin{enumerate}
  \item Advance every virion/bacterium \code{stageTimer\_s}.
  \item Free virions: Box--Muller Brownian with $\sigma\sqrt{dt}$.
  \item Free bacteria: run-and-tumble.
  \item Binding test: nearest live cell within
        \code{CONTACT\_RADIUS = 6\,\wu}; compute
        BindingMatcher score; gate $\ge 0.35$; increment bound
        spikes until threshold.
  \item BOUND $\to$ ENTERING after 60\,\biosec dwell.
  \item ENTERING $\to$ UNCOATING after 300\,\biosec dwell: deposit
        one viral genome copy into host \code{intraVirions}.
  \item Intracellular replication (currently autocatalytic; Part 8E
        replaces).
  \item Budding: continuous; caps at 3\,000 free virions globally.
  \item Apoptosis spill: the \code{releasePathogens} hook on BODIES
        transition dumps intracellular virions + bacteria outside
        the cell radius.
\end{enumerate}

\section{Negative control test}

\begin{lstlisting}
// from tests/headless_infection.cpp
if (randomizeSpike) {
    VirionSpec& vs = gPathogens().virionSpecs[vidx];
    vs.spike_logP = -10.0f;
    vs.spike_mw   =  1.0f;
    vs.spike_hbd  = 99;
    vs.spike_hba  = 99;
    vs.spike_aromatic = 99;
    vs.preferredReceptors.clear();
}
\end{lstlisting}
Run result: positive control (real flu) bound peak 49/50, infection
peak 358, clearly amplifying; negative control (scrambled spike)
bound peak 0, no infection. This is the empirical proof that the
binding is descriptor-driven, not scripted.

% ═════════════════════════════════════════════════════════════════════
\part{Immunity (scaffold)}
% ═════════════════════════════════════════════════════════════════════

\chapter{Innate Immunity}
\label{ch:innate}

\file{src/simulation/biochem/Immunity.h} (271 lines) holds
Pattern-Recognition Receptor state (TLRs, RLRs, NLRs, CLRs), PAMP
recognition, NF-$\kappa$B $\to$ cytokine cascade, complement, NK
perforin/granzyme, Type-I IFN response, dendritic antigen
presentation, and a V(D)J BCR/TCR repertoire skeleton.

\section{PAMP recognition}

\begin{lstlisting}
enum PAMPtype {
    PAMP_LPS, PAMP_PEPTIDOGLYCAN, PAMP_FLAGELLIN,
    PAMP_dsRNA, PAMP_ssRNA, PAMP_CpG_DNA, PAMP_MANNAN
};
\end{lstlisting}
Each PAMP maps to its characteristic receptor (LPS $\to$ TLR4,
peptidoglycan $\to$ TLR2, flagellin $\to$ TLR5, dsRNA $\to$ TLR3
endosomal + RIG-I cytosolic, ssRNA $\to$ TLR7/8, CpG DNA $\to$ TLR9,
mannan $\to$ Dectin-1). Activation threshold is modelled as a Hill
function; downstream signalling flows into a shared NF-$\kappa$B
pool that drives TNF-$\alpha$, IL-1$\beta$, IL-6 secretion into the
MediumField (planned, not yet wired in runtime).

\section{MHC-I presentation}

Host cells carry an MHC-I scalar; viruses express
\code{MHC\_I\_downregulation} which subtracts from it. NK cells
(when implemented as SimCell subtype) kill missing-self.

\chapter{Adaptive Immunity (skeleton)}
\label{ch:adaptive}

V(D)J recombination, somatic hypermutation, class switching, and
TCR-peptide-MHC scanning are present as \emph{data structures} but
not yet driven by the runtime. Part 7 of the plan (§44.10)
promotes them to active behaviour in 2 focused sessions.

% ═════════════════════════════════════════════════════════════════════
\part{Rendering}
% ═════════════════════════════════════════════════════════════════════

\chapter{Metal Rendering Pipeline}
\label{ch:metal}

\textsc{CellSim} renders natively on Apple's Metal API. The
pipelines live in \file{assets/shaders/*.metal} and are compiled
from source at startup (\texttt{[Metal] Compiled shader:
CellRender.metal}, \emph{etc.}).

\section{Shader inventory}

\begin{itemize}
  \item \file{CellRender.metal}: cell shell (sphere with fresnel
        rim and glow intensity; apoptosis fade blended in).
  \item \file{OrganelleRender.metal}: GLB-loaded organelle meshes
        with per-instance transforms, used for the focused cell in
        single-cell mode.
  \item \file{MoleculeRender.metal}: ball-and-stick atoms + bonds
        via instanced spheres and cylinders. Doubles as the
        pathogen / apoptotic-body renderer.
  \item \file{FluidRender.metal}: ray-marched dish fluid with
        wave deformation + subtle depth shading.
\end{itemize}

\section{The atom pipeline}

Every colored sphere drawn in the scene (pathogens, apoptotic
bodies, intracellular molecules, receptor dots) routes through the
\code{moleculeAtomPipeline}. Each instance is a
\code{GPUAtomInstance\{position; radius; color; pad;\}}; the
vertex shader places the sphere and the fragment shader applies a
fresnel rim whose intensity is driven by the instance's
\code{color} alpha. The same pipeline also draws thousands of
atoms from ball-and-stick molecules in single-cell mode.

\section{LOD tricks}

Molecule bond cylinders and small atoms disappear past a camera
distance threshold to keep draw calls down. The Hertz wall of the
cell-interior physics and the GPU-side membrane clamp are
deliberately not the same radius; the visual shell is slightly
larger than the physics shell so atoms don't appear to graze the
outer surface.

\chapter{Cell Shell, Organelles, Pathogens}
\label{ch:rendercell}

\code{syncCellInstances()} (around \file{src/main.mm} line 2326)
rebuilds the per-frame \code{CellInstance} array. Every cell's
shell carries position, radius, phase, LOD level, fate+phase colour,
glow intensity, and furrow depth. Apoptosis progress, PS-flip warm
tint, alpha ghosting, and caspase-driven blebbing are all baked
into the same shell instance with no shader changes.

Phase 7's infection tint (cyan-green for viral load, yellow-orange
for bacterial load) was added by reading
\code{SimCell::intraVirions, assembledVirions, intraBacteria} and
blending the cell colour. Light-mode users see infected cells glow
without needing to open the intracellular view.

% ═════════════════════════════════════════════════════════════════════
\part{UI, Export, Testing}
% ═════════════════════════════════════════════════════════════════════

\chapter{ImGui Panels and Setup Overlay}
\label{ch:ui}

The UI lives in \code{drawUI()} inside \file{src/main.mm}, starting
around line 7340. It is organised as a stack of ImGui windows:

\begin{itemize}
  \item \textbf{Cell Mode --- Setup} (modal at startup): preset
        templates, initial-cell slider, medium concentrations, pH,
        growth factor, time scale, light-mode checkbox.
  \item \textbf{Cell Mode} (top-center): re-setup button, cell
        count, follow-cell toggle, selected-cell index.
  \item \textbf{Drug Lab (chemistry-driven)}: compound dropdown,
        descriptor readout, log-C slider, Apply Uniformly button,
        top-affinities table.
  \item \textbf{Pathogen Lab}: virion + bacterium species
        dropdowns, best-match receptor + score, Inoculate buttons,
        live counts.
  \item \textbf{Test Actions} (light mode, lightning icon):
        one-click shortcuts for flu / HBV / adenovirus inoculation,
        E. coli / Listeria / S. aureus release, cisplatin /
        doxorubicin / paclitaxel / staurosporine, force apoptosis,
        heal cell, clear pathogens.
  \item \textbf{Export Data}: native folder picker + one-click
        bundle writer.
  \item \textbf{Population Stats / Culture Medium / Population
        Dynamics}: stat readouts in colony mode.
  \item \textbf{Cell Research Panel / Molecules Inside Cell /
        Mitosis Debug}: per-cell deep-dive panels in single-cell
        mode.
\end{itemize}

\section{Setup templates}

\code{applySetupTemplate(idx)} (line 7081) sets every field in
\code{gSetup} for each template id. The Test Mode template (id 7,
added 2026-04-20) sets \code{lightMode = true} plus one cell
and lower time scale.

\chapter{Export System}
\label{ch:export}

\code{exportAllDataWithPicker()} opens an NSOpenPanel for the user
to pick a destination folder, then dispatches to
\code{exportAllData(dir)} which writes a timestamped bundle of six
files: \file{manifest.txt}, \file{population.csv},
\file{cells.csv}, \file{medium\_field.csv}, \file{apo\_bodies.csv},
\file{medium\_chems.csv}.

The manifest includes run summary and the mass-balance drift
number (under $10^{-3}$). Post-run analysis (compare to CTC HeLa
reference, overlay with literature growth curves) uses the CSV
files directly.

\chapter{Headless Test Framework}
\label{ch:headless}

\file{tests/headless\_doubling.cpp} (285 lines) runs the simulator
without any rendering, for reproducible validation. Default run:

\begin{verbatim}
./build/cellsim_headless 46.0 60.0 125
  # 46 bio-h target, 60 s/step, 125 init cells
\end{verbatim}

Produces \file{logs/headless\_doubling.csv} (time series) and
\file{logs/compare\_vs\_reference.csv} (per-row absolute + relative
error vs CTC HeLa seq02). Seed-swept mean
$|\text{rel err}|$ over 92 matched timepoints: 5.1\% at seed 0x1
(default), 6.9\% at 0x2, 7.4\% at 0xFFFFFFFE, 6.5\% at 0xDEADBEEF;
mean across 5 seeds $\approx 6.2\% \pm 1.0\%$. Phase~P1
(2026-04-20) introduced the \code{--seed} CLI and the seeded
\code{simrng} wrapper; the single-number 5.1\% legacy claim is
superseded by this distribution.

\file{tests/headless\_infection.cpp} (104 lines) runs the Phase 7
infection chain with a positive (real flu) and negative (scrambled
spike) control. Four assertions guard the discipline:
\begin{enumerate}
  \item Bound peak $\ge 20$ (binding actually happens).
  \item Infection peak $\ge 50$ (amplification beyond seed).
  \item Amplification beyond seed (real replication).
  \item Positive $\gg$ Negative (descriptor-driven, not scripted).
\end{enumerate}

% ═════════════════════════════════════════════════════════════════════
\part{Phase 8A --- Multi-Gene Central Dogma}
% ═════════════════════════════════════════════════════════════════════

\chapter{Why Phase 8A Matters}
\label{ch:phase8a}

In April 2026 the pathogen system was inspected and found to have a
toy replication kernel:
\begin{lstlisting}
// Old code — Part 7 autocatalytic replication.
float dG = vs.replicationRate_per_s * c.intraVirions[s] * dt_biosec;
c.intraVirions[s] += dG;
\end{lstlisting}
This line reproduces the qualitative shape of exponential viral
replication but makes no use of host ribosomes, no use of host
RNA Pol II, no use of stoichiometric assembly. A virion ``divides
so fast without even using the cell''.

Phase 8A is the load-bearing change that replaces this with a real
gene registry. The existing CentralDogma engine (transcription,
splicing, translation, folding, PTM) was already realistic but
hard-wired to the single HBB gene. Phase 8A lifts this restriction.

\section{The four structural changes}

\begin{enumerate}
  \item \textbf{GeneTemplate struct}: \code{id, dnaSeq, startCodon,
        stopCodon, codingMRNA, promoterStrength, polymerase,
        proteinName, copiesMade}.
  \item \textbf{CentralDogmaState registry}: \code{std::vector
        <GeneTemplate> viralGenes} + \code{std::unordered\_map
        <std::string, int> viralProteinCount}.
  \item \textbf{Per-transcript gene tagging}:
        \code{TranscriptionState::geneIndex},
        \code{TranscriptState::{geneIndex, transcriptStartCodon,
        transcriptStopCodon, geneId, proteinTag}},
        \code{ProteinProductState::{geneIndex, geneId, proteinName}}.
  \item \textbf{Promoter-weighted pickup}: \code{pickGeneForTranscription()}
        chooses among HBB (weight 1.0) and every registered viral
        gene (weight \code{promoterStrength}) each time a Pol II
        starts transcribing.
\end{enumerate}

\section{Backward compatibility}

HBB transcription is now \emph{identical} to the pre-8A path. The
elongation read routes through \code{geneSequence(geneIndex, len)}:
for \code{geneIndex == -1} (the default), this is
\code{return HBB\_SEQUENCE} with no allocation --- the same
pointer-to-literal access that was there before. Only when a viral
gene is registered does the string-based path activate.

\section{Headline result}

Post-8A validation:
\begin{itemize}
  \item \code{cellsim\_headless}: mean $|\text{rel err}|$ vs CTC
        seq02 \textbf{varies 5.1--7.4\%} across 5 seeds
        ($\{0x1, 0xCE11517E, 0x2, 0xFFFFFFFE, 0xDEADBEEF\}$;
        mean $\approx 6.2\%$, SD $\approx 1.0\%$).
        The previously-reported ``5.1\%'' figure was measured at
        the default seed (0x1, libc's implicit initial state) and is
        one end of the distribution, not the mean.
        This was surfaced by the Phase~P1 determinism audit
        (2026-04-20): the simulator is now explicitly seedable
        (\code{--seed} CLI on headless binaries), and the
        single-number calibration claim is replaced by an honest
        seed-swept band. Validation with a new drug/virus should
        report mean $\pm$ SD across $\geq 5$ seeds, not one run.
  \item \code{cellsim\_infection}: all four assertions pass (bound
        peak 49, infection peak 358, amplification, positive
        $\gg$ negative).
  \item \code{cellsim\_p53\_pulsing} (Phase G2, new 2026-04-20):
        first biological-mechanism validator in CellSim history.
        Tests the TP53↔MDM2 stress-response loop with literature-
        calibrated rate constants:
        \begin{itemize}
          \item $K_d$ (p53↔MDM2) = 0.20 normalised ≙ 200 nM
                (Picksley 1994, Bottger 1997);
          \item MDM2 half-life $\approx$ 30 bio-min (Momand 1992);
          \item p53 half-life under basal MDM2 $\approx$ 30 bio-min
                (Haupt 1997, Kubbutat 1997);
          \item ATM-Ser15 shield = 0.70, permitting $\sim$30 \%
                residual MDM2 vulnerability consistent with
                Loewer 2013, Stewart-Ornstein 2017;
          \item ATM activation threshold at damageLevel $\approx$ 0.10
                (Rothkamm 2003).
        \end{itemize}
        Four assertion gates (control, damage, MDM2, specificity)
        all pass across 3 pinned seeds.
        Seed 0x1 (default) control scenario: p53 peak = 0.101
        (within $\pm$13 \% of analytical baseline 0.089);
        damage scenario ($\text{damageLevel}=0.30$): p53 peak = 0.292
        (3.3× baseline) and MDM2 peak = 0.507 (2.4× baseline).
        The p53(damage)/p53(control) ratio = 2.9, demonstrating
        mechanism specificity rather than constitutive activation.
        Simple 2-variable ODE (G2) settles to a stable elevated p53
        under sustained damage. Phase~G3 added an MDM2-mRNA
        intermediate with Hill (n=8) transactivation — exactly the
        Geva-Zatorsky 2006 3-variable delay topology — and now
        reproduces Purvis-2012-style sustained oscillation:
        7 p53 peaks detected across a 24-h damage window at
        damageLevel=0.30, mean period $\approx 3.03$ h, period
        shortening with increased stress (gate g: period(0.35) $<$
        period(0.30)), no peaks under control (damageLevel=0).
        All eight combined gates (G2 a-d + G3 e-h) pass across
        five pinned seeds (0x1, 0xCE11517E, 0xDEADBEEF, 0x2,
        0xFFFFFFFE).
  \item \code{CellSim} main app: builds clean.
  \item \textbf{Zero biology drift at default seed}: G2 stress-response
        ODE and the P1 SimRng wrapper + \code{viralProteinCount} map
        swap together produce byte-identical seq02 output vs pre-P1
        baseline (MD5-verified at default seed 0x1, 48 bio-h scenario,
        125 initial cells, 377 final, 252 divisions).
\end{itemize}

\chapter{Using the Gene Registry (Phase 8E preview)}
\label{ch:phase8a-use}

To turn Phase 8A into real viral replication, Phase 8E wires two
calls inside \code{Simulation::updatePathogens}:

\begin{lstlisting}
// When a virion uncoats (stageTimer > uncoatDwell):
GeneTemplate g;
g.id = vs.id + "_genome";
g.dnaSeq = vs.genomeSeq;
g.promoterStrength = 2.0f;           // viral promoter stronger than HBB
g.polymerase = GenePolymerase::VIRAL_RdRp;
g.proteinName = vs.id + "_capsid";   // etc. — one entry per gene
host.program.cdogma.registerViralGene(g);

// When viralProteinCount satisfies stoichiometry:
if (host.program.cdogma.viralProteinCount[vs.id + "_capsid"] >= 180
 && host.program.cdogma.viralProteinCount[vs.id + "_spike"]  >= 500
 && host.program.cdogma.viralProteinCount[vs.id + "_matrix"] >= 180) {
    // Assemble one virion, decrement counts atomically.
}
\end{lstlisting}

Host RNA Pol II competes for its three slots among HBB + all
registered viral genes, with \code{promoterStrength} biasing the
distribution. Ribosomes flow viral mRNAs through translation
without special case. The existing PTM, folding, and nuclear-export
pipelines all apply. The last remaining toy --- the assembly
gate --- is explicit and mass-conservative.

% ═════════════════════════════════════════════════════════════════════
\part{Running and Calibrating}
% ═════════════════════════════════════════════════════════════════════

\chapter{Build System}
\label{ch:build}

CMake 3.20+ drives the build. Three targets:
\begin{itemize}
  \item \code{CellSim}: the main Metal app, linked against GLFW,
        ImGui, vcpkg libs.
  \item \code{cellsim\_headless}: the CSV validator.
  \item \code{cellsim\_infection}: the Phase 7 infection smoke test.
\end{itemize}
\file{CMakeLists.txt} lists the source files; only two \code{.cpp}
files exist in \file{src/simulation/biochem/}
(\code{Bioagent.cpp, TargetLibrary.cpp}) --- every other biochem
module is header-only.

\chapter{Calibration Playbook}
\label{ch:calibration}

Two reference targets guide the calibration:
\begin{itemize}
  \item \textbf{CTC Fluo-N2DL-HeLa seq02} (Mitocheck, 125-cell start,
        46 bio-h, 74 ground-truth divisions). This is the primary
        growth-curve benchmark.
  \item \textbf{Chao 2019 Fig. 1} (HeLa doubling time 24 h, DMEM
        10\% FBS, 5\% CO$_2$). Used to sanity-check the cycle
        target durations.
\end{itemize}
The calibration scripts live in \file{scripts/calibrate\_rate.sh}
(grid sweep over \code{SLOW\_DT\_SCALE} and \code{MEDIUM\_DT\_SCALE}).

\chapter{Scientific Validation Plan}
\label{ch:validation}

For each major subsystem:
\begin{itemize}
  \item \textbf{Cycle} --- phase distribution matches HeLa BrdU
        pulse-chase (Chao 2019 Fig. 2).
  \item \textbf{Apoptosis} --- death kinetics under drug (cisplatin,
        doxorubicin) match Sia 2008 + Ye 2017.
  \item \textbf{Metabolism} --- Warburg lactate output matches
        Park 2016 metabolomics.
  \item \textbf{Central Dogma} --- HBB protein produced within
        $\sim$5\,\biomin of G1 entry (Singh 2013 elongation rate).
  \item \textbf{Drug PK} --- cisplatin intracellular peak at $\sim$30
        bio-min, receptor occupancy curves match Lipinski + Fick.
\end{itemize}

% ═════════════════════════════════════════════════════════════════════
\part{Reference and Appendices}
% ═════════════════════════════════════════════════════════════════════

% ═════════════════════════════════════════════════════════════════════
\part{Deep Technical Reference}
% ═════════════════════════════════════════════════════════════════════

\chapter{SimCell: Complete Field Reference}
\label{ch:simcell-deep}

This chapter walks every non-derived field of \code{SimCell}, naming
its units, valid range, update site, and readers. It is the
definitive field reference; Ch.~\ref{ch:simcell} above is the
five-thousand-foot view.

\section{Spatial and rendering fields}

\begin{longtable}{p{3.3cm} p{2.2cm} p{7.5cm}}
\toprule
\textbf{Field} & \textbf{Units} & \textbf{Meaning / update site}\\ \midrule
\code{position}    & \wu    & World-space centre of the cell body; updated by \code{updatePhysics} integrating velocity; also clamped by the floor and scene bound.\\
\code{velocity}    & \wu/\realsec & Integrated by Verlet-style stepping in \code{updatePhysics}; damped each frame; perturbed by Hertz repulsion from neighbouring cells and the floor.\\
\code{radius}      & \wu    & Baseline visual radius at \code{size=1} (typically 2.0 for HeLa); rarely changes.\\
\code{size}        & unitless & Multiplier on \code{radius}; grows during G1 biomass accumulation, shrinks during apoptosis (up to 35\,\% by \code{SHRINK\_FRAC\_AT\_COMPLETE}).\\
\code{alive}       & bool   & Master gate for every per-cell update; set false by the BODIES transition of \code{updateApoptosis}; cells are later erased from the vector by a \code{remove\_if}.\\
\code{cellUid}     & int    & Monotonically-increasing unique id; survives for the lifetime of the cell instance and is never reused; used by telemetry to track individual cells across division events.\\
\code{cloneId}     & int    & Lineage id, inherited by daughters; lets the user colour cells by clonal origin.\\
\code{generation}  & int    & 0 for founders; incremented on division.\\
\bottomrule
\end{longtable}

\section{Bioenergetics fields}

\begin{longtable}{p{3.3cm} p{2.2cm} p{7.5cm}}
\toprule
\textbf{Field} & \textbf{Units} & \textbf{Meaning / update site}\\ \midrule
\code{ATP}            & \% (0--100) & Normalised ATP availability; fed from the metabolism engine's real mM ATP pool via \code{bioATPToSimPercent}; consumed by CDK, replication, V-ATPase, motor proteins, SERCA.\\
\code{biomass}        & bm     & Dry-mass equivalent; $1\,\text{bm} \approx 200\,\text{pg}$; grows in G1 at rate \code{BIOMASS\_SYNTH\_K} scaled by nutrient availability; halved on division; released to medium on apoptosis.\\
\code{mitoPotential}  & arbitrary (0--100) & $\Delta\Psi_m$ proxy; drops under hypoxia, rotenone, or the MOMP event; low values feed the mito-dysfunction trigger of apoptosis.\\
\code{mitoHealth}     & 0--1   & Integrity of mitochondrial population; drops with ROS; below \code{MITO\_HEALTH\_SWITCH = 0.4} the cell flips \code{glycolytic=true}; restored above \code{MITO\_HEALTH\_RESTORE = 0.7}.\\
\code{glycolytic}     & bool   & Warburg switch; when true, metabolism prefers glycolysis + lactate export over OxPhos; see Ch.~\ref{ch:metabolism}.\\
\code{warburgTimer}   & \biosec & Accumulated time in glycolytic mode; drives adaptive acidification of the medium.\\
\code{hypoxiaTimer}   & \biosec & Time with local O$_2$ below \code{HYPOXIA\_MODERATE}; used to transition to HIF-1$\alpha$ signalling.\\
\code{hypoxiaIntensity}& 0--1  & $1 - [\text{O}_2]/\code{HYPOXIA\_MODERATE}$, clamped.\\
\code{necrotic}       & bool   & Set when O$_2$ stays below \code{HYPOXIA\_NECROTIC\_O2 = 0.035} for more than \code{HYPOXIA\_NECROTIC\_TIME = 120}\,\biosec; a fast, uncontrolled death path.\\
\bottomrule
\end{longtable}

\section{Cell cycle fields}

\begin{longtable}{p{3.3cm} p{2.2cm} p{7.5cm}}
\toprule
\textbf{Field} & \textbf{Units} & \textbf{Meaning / update site}\\ \midrule
\code{cdk.CycD}       & 0--1   & Cyclin D; rises with growth factor and mitogenic signalling; drives Rb phosphorylation.\\
\code{cdk.CycE}       & 0--1   & Cyclin E; transcribed by E2F; drives late G1 / S entry.\\
\code{cdk.CycA}       & 0--1   & Cyclin A; active through S and into G2.\\
\code{cdk.CycB}       & 0--1   & Cyclin B; peaks at mitosis entry; degraded by APC/C.\\
\code{cdk.Rb}         & 0--1   & Hyperphosphorylated when $0$, hypo when $1$. Hypo-Rb sequesters E2F.\\
\code{cdk.p21}        & 0--1   & CDK inhibitor; rises under DNA damage + crowding; gates G1/S and G2/M.\\
\code{cdk.p53}        & 0--1   & Tumour suppressor; accumulates under damage; transcribes p21, Bax, Puma, GADD45.\\
\code{cdk.E2F}        & 0--1   & Released from Rb, transcribes CycE, CycA, and many S-phase genes.\\
\code{cdk.APC\_C}     & 0--1   & Anaphase-promoting complex; active at mitotic exit; degrades CycB.\\
\code{phase}          & 0--3   & G1=0, S=1, G2=2, M=3.\\
\code{cycleTimer}     & \biosec & Time accumulated in current phase; compared against \code{CYCLE\_[G1,S,G2,M]\_TARGET\_BIOSEC}.\\
\code{cycleProgress}  & 0--1   & \code{cycleTimer / targetForCurrentPhase}.\\
\code{checkpointG1Passed} & bool & Latched once \code{evalG1S} passes.\\
\code{checkpointG2Passed} & bool & Latched once \code{evalG2M} passes.\\
\code{g1WaitTimer}    & \biosec & Time spent waiting at G1/S for checkpoint.\\
\code{g2WaitTimer}    & \biosec & Time spent waiting at G2/M.\\
\bottomrule
\end{longtable}

\section{Damage, ROS, and stress}

\begin{longtable}{p{3.3cm} p{2.2cm} p{7.5cm}}
\toprule
\textbf{Field} & \textbf{Units} & \textbf{Meaning / update site}\\ \midrule
\code{stress}        & 0--100 & Generic stress load; raised by damage, by replication errors, by ER unfolded-protein accumulation, by drugs; integrates toward apoptosis threshold.\\
\code{ROS}           & 0--100 & Reactive oxygen species; produced by OxPhos electron leak (esp.\ under high mitoPotential and low antioxidants) and by drugs binding \code{TGT\_MITO\_COMPLEX\_I} or \code{TGT\_TOPO2}; consumed by glutathione (not explicitly modelled --- treated as first-order decay).\\
\code{damageLevel}   & 0--2   & DNA damage burden; rises with replication errors, UV, radiomimetics, TopII poisons; gate on G1/S and G2/M; feeds p53 trigger of apoptosis.\\
\code{age}           & \biosec & Time since birth; purely diagnostic.\\
\bottomrule
\end{longtable}

\section{Fate decision fields}

Every few ticks, \code{updateFate(c, dt)} scores three potential fates
(PROLIF, QUIESCENT, APOPTOTIC) into \code{fateScores[3]}. The highest
score after a debounce window locks the fate into \code{c.fate}. The
cell then transitions accordingly: PROLIF cells feel growth-factor
kick to CDK; QUIESCENT cells hold their cycle timer; APOPTOTIC cells
receive a PRIMED push into the apoptosis engine.

\begin{longtable}{p{3.3cm} p{2.2cm} p{7.5cm}}
\toprule
\textbf{Field} & \textbf{Units} & \textbf{Meaning / update site}\\ \midrule
\code{fate}         & enum  & \code{SIM\_FATE\_UNDETERMINED=0, SIM\_FATE\_PROLIF=1, SIM\_FATE\_QUIESCENT=2, SIM\_FATE\_APOPTOTIC=3}.\\
\code{fateScores[3]} & float & Running vote tally; highest wins after a debounce window.\\
\code{fateTimer}    & \biosec & Debounce window; fate locks at expiry.\\
\code{fateLocked}   & bool  & After lock, scores no longer change fate.\\
\code{atpDangerTimer}& \biosec & Time with ATP below critical; drives energy-stress path.\\
\bottomrule
\end{longtable}

\chapter{The Cell Cycle --- Complete ODE Reference}
\label{ch:cycle-deep}

The cyclin/CDK dynamics inside \code{updateCellCycle} are a
simplified Aldridge-style ODE system, hand-scaled so the
default generation time matches 24\,\bioh. This chapter names every
rate term and its coupling.

\section{Differential equations}

Let $D, E, A, B$ denote \code{CycD, CycE, CycA, CycB}; $R$ denote
\code{Rb}; $F$ denote \code{E2F}; $P$ denote \code{p21}; $Q$ denote
\code{p53}; $K$ denote \code{APC\_C}. Let $gf$ denote local
\code{nutrients.get(MS\_GROWTH\_F, x, z)}. Let $\Delta$ denote
\code{damageLevel}. Then:

\begin{align}
  \dot D &= k_{D,\text{synth}}\, \sigma(gf / K_{gf}) - k_{D,\text{deg}}\,D,
           \label{eq:cycD}\\
  \dot F &= k_{F,\text{release}}\,(1-R)\,(1-P) - k_{F,\text{deg}}\,F,
           \label{eq:E2F}\\
  \dot E &= k_{E,\text{synth}}\,F - k_{E,\text{deg}}\,E,
           \label{eq:cycE}\\
  \dot A &= k_{A,\text{synth}}\,F\,(1-P) - k_{A,\text{deg}}\,A,
           \label{eq:cycA}\\
  \dot B &= k_{B,\text{synth}}\,A - k_{B,K}\,K\,B - k_{B,\text{deg}}\,B,
           \label{eq:cycB}\\
  \dot R &= -k_{R,\text{phos}}\,(D + E)\,R + k_{R,\text{dephos}}\,(1-R),
           \label{eq:Rb}\\
  \dot P &= k_{P,\text{synth}}\,(Q + c_{\text{mech}}\,\max(0, n-N_c)) - k_{P,\text{deg}}\,P,
           \label{eq:p21}\\
  \dot Q &= k_{Q,\text{synth}}\,\Delta - k_{Q,\text{deg}}\,Q,
           \label{eq:p53}\\
  \dot K &= k_{K,\text{synth}}\,B\,(1-\text{sacHold}) - k_{K,\text{deg}}\,K.
           \label{eq:APC}
\end{align}

Here $\sigma$ is a Hill function with $n=2$, $K_{gf}$ the half-max
growth-factor concentration; $c_{\text{mech}} = \code{MECH\_P21\_COUPLING}$
is the mechano-crowding coefficient; $n$ the local neighbour count;
$N_c = \code{CONTACT\_INHIBIT\_NBRS}$; \text{sacHold} a boolean
(1 while SAC is blocking anaphase, 0 otherwise).

\section{Rate constants}

The per-reaction rates are chosen so G1 occupies
$\sim 11\,\bioh$, S $\sim 8\,\bioh$, G2 $\sim 4\,\bioh$, M
$\sim 1\,\bioh$ at baseline; they are scaled against
\code{CDK\_DT\_SCALE} for stiffness:

\begin{center}
\begin{tabular}{lrl}
\toprule
\textbf{Rate} & \textbf{Value} & \textbf{Units}\\ \midrule
$k_{D,\text{synth}}$  & 0.08  & $\biosec^{-1}$\\
$k_{D,\text{deg}}$    & 0.05  & $\biosec^{-1}$\\
$k_{F,\text{release}}$& 0.20  & $\biosec^{-1}$\\
$k_{F,\text{deg}}$    & 0.10  & $\biosec^{-1}$\\
$k_{E,\text{synth}}$  & 0.05  & $\biosec^{-1}$\\
$k_{E,\text{deg}}$    & 0.04  & $\biosec^{-1}$\\
$k_{A,\text{synth}}$  & 0.04  & $\biosec^{-1}$\\
$k_{A,\text{deg}}$    & 0.03  & $\biosec^{-1}$\\
$k_{B,\text{synth}}$  & 0.04  & $\biosec^{-1}$\\
$k_{B,K}$             & 0.50  & $\biosec^{-1}$\\
$k_{B,\text{deg}}$    & 0.02  & $\biosec^{-1}$\\
$k_{R,\text{phos}}$   & 0.08  & $\biosec^{-1}$\\
$k_{R,\text{dephos}}$ & 0.005 & $\biosec^{-1}$\\
$k_{P,\text{synth}}$  & 0.03  & $\biosec^{-1}$\\
$k_{P,\text{deg}}$    & 0.02  & $\biosec^{-1}$\\
$k_{Q,\text{synth}}$  & 0.04  & $\biosec^{-1}$\\
$k_{Q,\text{deg}}$    & 0.02  & $\biosec^{-1}$\\
\bottomrule
\end{tabular}
\end{center}

These values are the target of Phase L of the plan (sensitivity
analysis + rigorous cycle-ODE refit).

\section{Coupling to drugs}

When \code{TGT\_CDK1}, \code{TGT\_CDK2}, \code{TGT\_KINASE\_EGFR} or
\code{TGT\_POL2} is bound (Ch.~\ref{ch:targetlib}), each target's
modulator scales one or more cyclin pools multiplicatively. The
modulator is called \emph{after} the ODE step on the same tick, so
drug pressure compounds tick-over-tick. The one-tick delay is
intentional (stiffness control).

\chapter{Mitosis --- Phase-by-phase Walkthrough}
\label{ch:mitosis-deep}

Mitosis in \textsc{CellSim} has two parallel implementations. The
background per-cell program is lean and drives cytokinesis in
colony mode. The \emph{visual} mitosis program, active only for the
primary focused cell in single-cell mode, drives the rendered
chromatin condensation, spindle, metaphase plate, anaphase
separation, and furrow. This chapter documents both.

\section{Prophase}

Entry criteria: $\code{CycB} > 0.35$, \code{checkpointG2Passed ==
true}, \code{replicationReadyForM() == true},
\code{damageLevel < 1.5}.

Visible changes: \code{chromatinCondense} rises from 0 to 1 over
$\sim$10\,\biomin. In the visual renderer the DNA-A/B particle
positions are blended toward a compacted helix centred on each
sister chromatid. Nuclear envelope breakdown is not explicitly
modelled; it is represented by disabling nuclear pore attraction
for RNA / ribosome traffic.

\section{Prometaphase}

\code{poleSeparation} begins to rise. Two centrosomes (derived from
the single inherited centriole pair through centriole duplication,
which \code{centrosome\_20260419\_*.log} tracks) separate along the
$y$-axis. Kinetochores form on sister chromatids; microtubules
attach stochastically.

\section{Metaphase}

SAC holds anaphase entry until \code{evalSAC} returns passed. In the
rendered path, chromosomes align on the metaphase plate by blending
their positions toward a thin disk at cell centre.
\code{particlesDuplicated} must equal true (every chromosome
doubled and aligned) before \code{sacPassed} can be true.

\section{Anaphase}

APC/C activation (through \code{CycB} degradation via $K$ in the ODE
above) breaks the cohesin link between sister chromatids; they
separate toward opposite poles. Visual renderer tweens chromatid
positions toward the pole centroids.

\section{Telophase}

Chromatin decondenses. Nuclear envelopes re-form around each pole;
nuclear pore attraction re-enables; mRNA export resumes.

\section{Cytokinesis}

\code{furrowDepth} grows from 0 to 1 over $\sim$8\,\biomin. The cell
shell shader reads this and deforms the sphere with a radial pinch;
physics splits the interior particle cloud between the two halves
along the cleavage plane. At \code{MITO\_COMPLETE}, the single
\code{SimCell} is split by \code{processDivisions()} into two
daughters: daughter A inherits the original cell's UID and occupies
the leftmost half; daughter B gets a fresh UID and the rightmost
half. Both get fresh phase timers and reset their cycle state.

\section{SplitStats and asymmetry}

\begin{lstlisting}
struct SplitStats {
    float biomassFractionA = 0.50f;   // fraction to A; B = 1 - this
    float damageFractionA  = 0.50f;
    float telomereShortenA = 0.0f;
    float telomereShortenB = 0.0f;
    float cloneBiasA = 0.0f;
};
\end{lstlisting}

Asymmetry produces heterogeneity: an over-damaged mother gives most
of its damage to one daughter, letting the other escape senescence.
Telomere-shortening split is always 50--100\,bp per division
(Allsopp 1992), reproducing the Hayflick limit $\sim$50--70
generations.

\chapter{Apoptosis --- The Eleven-Trigger Engine}
\label{ch:apoptosis-deep}

This chapter enumerates each trigger, how the engine integrates
them, how MOMP commits, how caspase-3 activates, how PS flips, and
how the apoptotic-body pipeline commits mass to the medium.

\section{Input triggers}

\code{ApoTriggers} struct (see \file{Apoptosis.h}):
\begin{lstlisting}
struct ApoTriggers {
    float p53;             // 0..1
    float ROS;             // 0..100
    float ATP_collapse;    // 0..1 (1 = severe)
    float hypoxia;         // 0..1
    float survival_loss;   // 0..1
    float anoikis;         // 0..1
    float crowding;        // 0..1
    float replicative;     // 0..1 (senescence)
    float fas_L;           // 0..1
    float mito_dysfunction;// 0..1
    float drug_pro_apop;   // 0..1
};
\end{lstlisting}

Each is computed by \code{Simulation::updateApoptosis} from the
\code{SimCell}'s primary state:
\begin{align*}
  p53 &= \mathrm{clamp}(\code{damageLevel}/1.0, 0, 1),\\
  ROS &= \code{c.ROS},\\
  ATP\_collapse &= \mathrm{clamp}(1 - \code{c.ATP}/30, 0, 1),\\
  hypoxia &= \code{c.hypoxiaIntensity},\\
  survival\_loss &= \mathrm{clamp}(1 - gf/50, 0, 1),\\
  anoikis &= \mathrm{clamp}(1 - \code{c.adhesionStrength}/0.4, 0, 1),\\
  crowding &= \mathrm{clamp}(\code{chronicPressureBioSec}/7200, 0, 1),\\
  replicative &= \mathrm{clamp}(\code{senescenceBioSec}/3600, 0, 1),\\
  fas\_L &= \mathrm{clamp}(\code{fasLExposure}/20, 0, 1),\\
  mito\_dysfunction &= \mathrm{clamp}(1 - \code{mitoPotential}/80, 0, 1),\\
  drug\_pro\_apop &= \text{cumulative from TargetLibrary modulators.}
\end{align*}

\section{Bcl-2 / Bax dynamics}

\code{ApoptosisState} holds $\code{Bcl2}, \code{BclXL},
\code{Bax}, \code{Bax\_active}, \code{Bak\_active},
\code{MOMP\_pores}, \code{cyt\_c\_cytosol}, \code{Smac\_cytosol},
\code{caspase\_3}, \code{caspase\_8}, \code{caspase\_9}, \code{IAP},
\code{Puma}, \code{Bid\_truncated},
\code{apoptosis\_progress}, \code{p53\_active}$, etc.

Key ODE terms in \code{ApoptosisEngine::stepOnce(h)} (DT\_MAX =
0.08\,\biosec sub-steps):
\begin{align}
  \dot{\text{Bax}}_{\text{active}} &= k_{ba}\,(\text{p53} + \text{Puma})\,(1 - \text{Bax}_{\text{active}}) - k_{bd}\,\text{Bcl2}\,\text{Bax}_{\text{active}},\\
  \dot{\text{MOMP}}                &= k_m\,\text{Bax}_{\text{active}}\,(1-\text{MOMP}) - k_{mr}\,\text{MOMP},\\
  \dot{\text{cyt\_c}}               &= k_{cc}\,\text{MOMP}\,(1-\text{cyt\_c}),\\
  \dot{\text{caspase\_9}}          &= k_{c9}\,\text{cyt\_c}\,(1-\text{caspase\_9})-k_{c9d}\,\text{IAP}\,\text{caspase\_9},\\
  \dot{\text{caspase\_3}}          &= k_{c3}\,(\text{caspase\_9}+\text{caspase\_8})(1-\text{caspase\_3}) - k_{c3d}\,\text{IAP}\,\text{caspase\_3}.
\end{align}

Stiffness control: the outer \code{step(dt)} subdivides by
$\lceil dt / \text{DT\_MAX}\rceil$, capped at 200 sub-steps per
outer call. This refactor in April 2026 fixed a silent bug where
the previous hard clamp \code{dt = fminf(dt, 0.01)} froze the
cascade across a physics tick.

\section{Phase transitions}

The engine moves from PRIMED to MOMP when
$\code{MOMP\_pores} > 0.5$. From MOMP it dwells
\code{MOMP\_DURATION\_BIOSEC = 300}\,\biosec then enters EXECUTION.
EXECUTION dwells 1\,800\,\biosec during which cytosol leaks at
\code{EXECUTION\_CYTOSOL\_FRAC\_PER\_BIOSEC} per second and
receptors at \code{EXECUTION\_RECEPTOR\_FRAC\_PER\_BIOSEC}. At
EXECUTION end, a FRAGMENTATION phase runs another 1\,800\,\biosec,
during which the cell's interior splits into 5--15 apoptotic bodies
(in rendered mode) and its remaining pools are distributed among
them.

\section{Secondary necrosis}

Apoptotic bodies that are not engulfed within 6\,\bioh enter
secondary necrosis via nanopore-driven osmotic swelling (Chen 2016):
water influx, membrane stretch, rupture, and mass dumping into the
medium. The engine tracks \code{osmoticSwell} and commits lysis at
\code{LYSIS\_RUPTURE\_SWELL}.

\chapter{DNA Replication --- Fork Mechanics in Depth}
\label{ch:replication-deep}

Each active replication fork is a \code{ReplicationForkState} struct
with position, direction, velocity, recruitment (how much polymerase
is bound), stress, pause timer, proofreading flag, and a stalled
boolean. \code{stepReplicationForks(dt)} integrates fork motion.

\section{Fork velocity}

\begin{equation}
  v_{\text{fork}} = v_{\text{max}} \cdot \code{recruitment}\cdot (1 - \code{stress})\cdot (\code{dntpAvailability}),
\end{equation}
with $v_{\text{max}} = \code{CDOGMA\_LIT\_FORK\_SPEED\_BP\_PER\_SEC} = 24$\,bp/s.
Polymerase recruitment starts at 0.2 on fire and rises toward 1.0
with time constant $\sim$20\,\biosec (origin licensing and
initiation).

\section{Stall and rescue}

Stress on the fork rises when:
\begin{itemize}
  \item Local homopolymer run $\ge 6$ (polymerase slippage risk).
  \item Local GC-fraction $> 0.75$ (secondary-structure slowdown).
  \item TopII-poison drug bound (\code{TGT\_TOPO2} modulator sets
        \code{replicationActive = false}).
\end{itemize}
Stalled forks (velocity $< 1$\,bp/s for $> 60$\,\biosec) emit a
rescue signal that licenses the nearest dormant origin, which then
fires in the next tick.

\section{Mismatch repair}

When a base is incorporated, a three-layer repair stack runs:
\begin{enumerate}
  \item Pol$\epsilon$ proofreading: 90\% efficient, one base back
        if an error is immediately detected.
  \item MutS/MutL MMR: 95\% of remaining errors caught within
        $\sim$20\,\biosec of incorporation.
  \item Escape: remaining errors become persistent mismatches.
\end{enumerate}
Escaped errors contribute to \code{damageLevel} and to
\code{unresolvedReplicationErrors}, which blocks G2/M.

\section{The forced-completion special case}

If every fork has terminated and \code{replicationProgress} is
$\ge 99\%$ but $< 100\%$, the engine force-completes the remaining
handful of unreplicated bases. This handles the biology of
post-replicative repair (Alberts Ch 5: lagging-strand gaps filled
after main S-phase). Without this, three stuck bases would hold the
cell in S-phase indefinitely and the colony simulation would never
divide in long runs.

\chapter{Transcription and Splicing --- Pol II Walkthrough}
\label{ch:transcription-deep}

\section{Slot idle / fire dynamics}

Each Pol II slot \code{transcription[i]} has an internal timer.
When idle, it accumulates \code{dt}; at some firing delay $3.5 + 2i$
seconds, \code{startTranscription(i)} is called. This staggers
the three Pol II slots so they initiate transcription at
$\sim$3.5, 5.5, 7.5\,\biosec offsets, producing overlapping
elongation patterns. The exact delay constants derive from the
transcription-initiation literature: TFIIH recruitment and promoter
escape take $\sim$3--10\,s in mammalian cells.

\section{Splicing order and timing}

Human HBB has two introns. The engine models them in sequence:
intron 1 (129\,bp) is removed first at 25\,\% splicesome progress,
intron 2 (852\,bp) at 70\,\%, and poly-A is added at 100\,\%. This
mirrors the co-transcriptional splicing bias: first introns are
typically removed before transcription ends, later introns after
(Singh 2013).

\section{Export through nuclear pores}

Mature mRNA exports at rate 0.60\,s$^{-1}$; the exponential time
constant is $\sim$1.7\,\biosec, matching the
single-molecule mRNA export rate measured in live cells by Mor
2010 (eLife).

\section{Phase 8A path for viral mRNA}

Viral genes skip splicing entirely. When elongation completes,
the engine sets \code{tx->spliced = true; tx->matureMRNA =
viralGenes[idx].codingMRNA}. This is correct for monocistronic
RNA viruses. DNA viruses that do use splicing (adenovirus) would
need a per-VirionSpec splicing map, which is a straightforward
extension of the gene registry.

\chapter{Translation --- tRNA, Ribosome, Peptide Release}
\label{ch:translation-deep}

\section{tRNA pool dynamics}

Twelve-slot charged-tRNA pool (\code{ChargedTRNAState[12]}). Each
slot is charged (\code{charged = true}) with one amino acid. When a
ribosome calls \code{findFreeTRNA()}, it gets an empty slot; the
engine then charges it with the amino acid matching the current
codon. Once the amino acid is delivered to the nascent peptide,
\code{releaseTRNA()} marks the slot free.

The pool size (12) is a compromise: real cells have $\sim 10^6$
tRNAs per cell, but the simulator renders individual tRNA particles
in the focused cell interior, so keeping 12 visible is enough to
convey concurrency without spamming the screen.

\section{Ribosome ``shuttle''}

Each \code{TranslationState} has a \code{codonProgress} float
advanced each tick by $8.0\,\text{s}^{-1}$. When it reaches 1, the
amino acid is appended to \code{nascentPeptide}, \code{currentCodon}
is incremented, and \code{primeCurrentCodon} reads the next codon.

\section{Stop-codon detection}

\code{primeCurrentCodon} sees a stop codon (UAA, UAG, UGA) and sets
\code{stopReached = true}. On the next tick, \code{releaseProtein}
is called, the tRNA is released, and the ribosome is freed.

\section{Why no E-site / A-site modelling?}

The engine does \emph{not} model the three-site ribosome cycle
(E--P--A). The rationale is that visible ribosome fidelity is
governed by tRNA availability and codon identity, not by site
occupancy. For the current rendering goals (shader-animated
elongation) the simplification is invisible; for future kinetic
studies (e.g.\ ribosome stalling and ribosome-associated quality
control), a full three-site model would be warranted.

\chapter{DNA Repair Pathways}
\label{ch:dna-repair}

\file{src/simulation/biochem/DNARepair.h} (435 lines) holds the
five major DNA repair pathways: base excision repair (BER),
nucleotide excision repair (NER), mismatch repair (MMR),
homologous recombination (HR), non-homologous end joining (NHEJ),
and translesion synthesis (TLS).

\section{BER}

Targets single-base lesions (oxidised bases, abasic sites). Key
enzymes: OGG1 (8-oxoG glycosylase), APE1 (AP endonuclease),
Pol$\beta$, ligase III. Per-lesion repair time $\sim$30\,\biosec.
BER flux scales with \code{ROS} since oxidised bases are the
dominant substrate.

\section{NER}

Targets bulky lesions (UV cyclobutane dimers, cisplatin adducts).
Global-genome NER (GG-NER, XPC-driven) and transcription-coupled
NER (TC-NER, CSA/CSB-driven). Per-lesion repair time
$\sim$10\,\biomin. Cisplatin-bound DNA (after
\code{TGT\_DNA\_INTERCALATION} saturates) drives NER throughput
toward saturation; residual unrepaired lesions raise
\code{damageLevel}.

\section{MMR}

Post-replicative mismatch correction: MutS$\alpha/\beta$ recognise
mismatches; MutL$\alpha$ nicks the daughter strand; Exo1 removes a
tract; Pol$\delta$ resynthesises. 95\% efficient within 20\,\biosec
of mismatch creation.

\section{HR}

Template-dependent double-strand break repair, used in S/G2 when a
sister chromatid is available. Rad51-coated ssDNA invades the
sister, Pol$\delta$ copies, Holliday junctions resolve. Error-free.
The rate-limiting step is end resection by MRN/CtIP
($\sim$30\,\biomin).

\section{NHEJ}

Template-independent double-strand break repair, active in G1 (no
sister available). Ku70/Ku80 bind ends, DNA-PKcs tethers, Artemis
processes, LigIV seals. Fast ($\sim$30\,\biomin) but error-prone:
introduces small indels.

\section{TLS}

Translesion synthesis: when a replication fork hits an unrepairable
lesion, it stalls; a specialized Pol (Pol$\eta$, Pol$\iota$,
Pol$\kappa$) is recruited to synthesise across the lesion. Error-
prone but prevents fork collapse. Models Pol$\eta$ bypass of UV
dimers with $\sim$90\% accuracy.

\chapter{Gene Regulation Network Scaffold}
\label{ch:grn}

\file{src/simulation/biochem/GeneRegulation.h} (235 lines) holds a
skeletal TF $\to$ target graph with Hill-function activation. Used
for p53 targets, E2F targets, NF-$\kappa$B targets, HIF-1$\alpha$
targets, and MYC targets.

\section{Data model}

\begin{lstlisting}
struct GRNEdge {
    std::string tf;
    std::string target;
    int sign;         // +1 (activate), -1 (repress)
    float K;          // half-max TF concentration
    int   n;          // Hill coefficient
};
std::vector<GRNEdge> edges;
\end{lstlisting}

\section{Hill function}

The transcription rate of gene $g$ is:
\begin{equation}
  r_g = r_g^{\text{basal}} + \sum_i s_i \cdot \frac{a_{\text{TF}_i}^{n_i}}{K_i^{n_i} + a_{\text{TF}_i}^{n_i}}
\end{equation}
where $a_{\text{TF}_i}$ is the TF activity (concentration or
phosphorylated fraction), $s_i = \pm 1$, and the sum is over edges
that target $g$.

\section{Anchored TFs}

The current engine surfaces: \code{p53\_active} (Ch.~\ref{ch:apoptosis-deep}),
\code{E2F} (Ch.~\ref{ch:cycle-deep}), \code{HIF\_1a}
(activated by hypoxia), \code{NF\_kB} (activated by PAMP/TNF),
\code{MYC} (activated by growth factor). Phase 7L of the plan
builds out a 200-edge graph across these TFs and their downstream
targets.

\chapter{Signalling Cascades}
\label{ch:signaling}

\file{src/simulation/biochem/Signaling.h} (588 lines) implements
the five canonical signalling pathways as coupled ODEs.

\section{Pathway inventory}

\begin{itemize}
  \item \textbf{Wnt}: $\beta$-catenin stabilisation, TCF/LEF
        targets. Frizzled receptor binds Wnt $\to$ Dvl
        activation $\to$ GSK-3$\beta$ inhibition $\to$
        $\beta$-catenin accumulates $\to$ nuclear translocation.
  \item \textbf{Notch}: ligand-induced proteolysis releases NICD;
        NICD + CSL transactivates Hes/Hey targets. Contact-
        dependent, requires neighbouring cell as signal source.
  \item \textbf{TGF-$\beta$}: T$\beta$RI/II heterodimer
        phosphorylates Smad2/3; Smad4 complex translocates;
        CDKN2A/CDKN1A upregulated $\to$ cycle arrest.
  \item \textbf{MAPK}: RTK $\to$ Ras $\to$ Raf $\to$ MEK $\to$
        ERK1/2. ERK substrates include Elk1 (transcription) and
        many cycle-cycling proteins (CyclinD, p21-suppression).
  \item \textbf{PI3K--AKT--mTOR}: RTK $\to$ PI3K $\to$ PIP3 $\to$
        AKT (Ser473); AKT substrates: Bad (pro-apoptotic) $\to$
        inactivated; Tuberin/TSC2 $\to$ mTORC1 release;
        FOXO3 $\to$ nuclear export; autophagy suppressed.
\end{itemize}

\section{Cross-talk}

ERK activates CycD transcription. PI3K--AKT inhibits Bad,
contributing a \emph{survival} term that counteracts the
\code{survival\_loss} apoptosis trigger. TGF-$\beta$ raises p21,
reinforcing cycle arrest. Hypoxia stabilises HIF-1$\alpha$ which
transcribes GLUT1, LDHA, VEGF. Contact inhibition via cadherins
raises $\beta$-catenin sequestration; high cell density therefore
lowers proliferative Wnt output.

These couplings are scaffolded in \file{Signaling.h} but not all
wired to the cycle ODE. Phase 7K brings them live.

\chapter{Cell Junctions and ECM}
\label{ch:junctions}

\file{src/simulation/biochem/CellJunctions.h} (208 lines) models
four junction types plus the ECM adhesome.

\section{Cadherin junctions}

E-, N-, P-cadherin levels; $\beta$-catenin / $\alpha$-catenin /
p120-catenin linker dynamics; vinculin recruitment under tension.

\begin{equation}
  F_{\text{adh}} = (\text{E-cad} + \text{N-cad})\cdot\frac{[\text{Ca}^{2+}]}{[\text{Ca}^{2+}]+K_{\text{Ca}}}\cdot\beta\text{-cat}\cdot\alpha\text{-cat}\cdot 100\,\text{pN}
\end{equation}
with tension-dependent vinculin reinforcement above 5\,pN.

\section{Desmosomes}

Desmoglein/desmocollin/plakoglobin/desmoplakin/plakophilin. Anchor
intermediate filaments between neighbouring cells; contribute
structural rigidity (important for epithelia) but no direct cycle
coupling in the current engine.

\section{Tight junctions}

Claudins/occludin/ZO-1/ZO-2. Create a paracellular seal + lateral
diffusion fence. In a monolayer the sealing contributes to
apical/basal polarity; dissolution (observed during EMT,
epithelial-mesenchymal transition) cues Cancer.h's metastatic
escape scaffolding.

\section{Integrins}

$\alpha\beta$ heterodimers bind ECM (fibronectin, collagen,
laminin). Inside-out and outside-in signalling. Focal adhesions
(FAK, paxillin, vinculin, talin) recruit as integrin clusters
mature. \code{SimCell::adhesionStrength} is a reduced scalar
proxy that walks from 0 (floating) to 0.95 (mature FA cluster)
over a timeline calibrated in Adhesion:: constants.

\section{Selectins}

Transient rolling adhesion for leukocytes. Scaffolded but not
currently driven at runtime. Phase 7L of the plan brings it live
for immune-cell rolling on endothelium.

\chapter{Proteome Quality Control}
\label{ch:proteome}

\file{src/simulation/biochem/Proteome.h} (380 lines) implements
chaperones, refolding, ubiquitin tagging, proteasomal degradation,
the unfolded protein response (UPR), heat-shock response (HSF1),
autophagy, and many post-translational modification ledgers.

\section{Chaperone pool}

HSP70 (cytosolic) and BiP/GRP78 (ER luminal); both consume ATP.
Each misfolded protein makes an approximate demand:
\begin{itemize}
  \item First refold attempt: 2\,ATP, $\sim$60\,\biosec.
  \item On failure, ubiquitination (three E3 ligases tracked).
  \item Proteasomal degradation (if cytosolic) or ERAD (if ER).
\end{itemize}

\section{HSF1 activation}

Heat-shock factor 1 is sequestered by HSP70 in basal conditions.
Misfolded protein load titrates HSP70 away, freeing HSF1, which
trimerises and enters the nucleus to transcribe chaperone genes
(including HSP70 itself, closing the feedback loop).

\section{UPR arms}

Three ER-stress sensors: IRE1$\alpha$ (which splices XBP1 mRNA to
XBP1s, an active TF), PERK (phosphorylates eIF2$\alpha$, globally
dampening translation), ATF6 (Golgi-cleaved to nuclear form,
transcribes chaperones). Combined output: transient translation
shutoff, chaperone upregulation, eventual apoptotic commitment via
CHOP if stress is chronic.

\section{Autophagy}

mTORC1 inhibits ULK1; low nutrients / high AMP:ATP ratio inhibits
mTORC1, activating ULK1, which nucleates autophagosomes. Beclin1
complex drives phagophore formation; LC3-I is lipidated to LC3-II
on the phagophore membrane. Autolysosome flux recycles proteins to
amino acids.

\chapter{Stem Cell Behaviour}
\label{ch:stem}

\file{src/simulation/biochem/StemCells.h} (253 lines) scaffolds
stemness, self-renewal, and differentiation pressure. The module is
not yet wired for HeLa simulation but is present for future
stem-cell models.

\begin{itemize}
  \item Nanog, Oct4, Sox2 core TF triangle for pluripotency.
  \item Asymmetric division with lineage bias (\code{cloneBiasA}).
  \item Niche-dependent self-renewal (FGF, BMP, Wnt inputs).
  \item Differentiation commitment: when any lineage TF crosses
        threshold, the cell exits the stem state permanently.
\end{itemize}

\chapter{Cancer Scaffolding}
\label{ch:cancer}

\file{src/simulation/biochem/Cancer.h} (189 lines) tracks oncogene /
TSG mutations, anoikis resistance, EMT, invasion, and metastasis
precursors. Again, scaffolded for future work; HeLa is assumed
transformed at init.

\section{Mutation ledger}

Each cell carries a list of \code{SNV} and \code{CNV} records;
mutation rate scales with \code{ROS} and failed MMR events.

\section{Anoikis resistance}

In normal epithelial cells, loss of integrin contact triggers
apoptosis (anoikis). Transformed cells acquire resistance through
constitutive AKT activation. The engine models this by scaling the
anoikis trigger by a \code{cancer\_anoikis\_resistance} factor (0 =
normal epithelium; 1 = fully transformed).

\chapter{Development Module}
\label{ch:development}

\file{src/simulation/biochem/Development.h} (225 lines) is
morphogen gradient machinery for future developmental biology
scenarios. The module holds positional information (via
concentration gradients) and Hox-like TFs that read threshold
positions. Not active in HeLa runs.

\chapter{Genome Model}
\label{ch:genome-model}

\file{src/simulation/biochem/GenomeModel.h} (259 lines) holds
ploidy information, chromatin state (epigenetic marks), and
chromosome structure (centromeres, telomeres, origins, nuclear
landmarks).

\begin{itemize}
  \item \textbf{Ploidy}: diploid default; HeLa is known aneuploid
        ($\sim$76--80 chromosomes typical). The engine keeps it
        diploid as a tractable baseline; aneuploid cells would
        require per-chromosome tracking.
  \item \textbf{Epigenetics}: per-gene \code{methylation} fraction
        and four histone mark densities (H3K4me3, H3K27me3,
        H3K9me3, H3K27ac). Phase 7M brings these live.
  \item \textbf{Telomeres}: length in kb; shortens 50--100\,bp per
        division; Hayflick limit at $\sim$4\,kb triggers
        senescence flag.
\end{itemize}

\chapter{Membrane Transport --- Hodgkin-Huxley Path}
\label{ch:hh}

\file{src/simulation/biochem/MembraneTransport.h} implements a full
Hodgkin--Huxley action-potential model for excitable cells.

\section{Gating variables}

\begin{align}
  \dot m &= \alpha_m(V)(1-m) - \beta_m(V)\,m,\\
  \dot h &= \alpha_h(V)(1-h) - \beta_h(V)\,h,\\
  \dot n &= \alpha_n(V)(1-n) - \beta_n(V)\,n.
\end{align}
with standard squid-axon rate constants scaled to 37$^\circ$C by
$Q_{10}$.

\section{Ionic current and membrane equation}

\begin{equation}
  I_m = \bar g_{\text{Na}}\,m^3 h\,(V - E_{\text{Na}}) + \bar g_{\text{K}}\,n^4\,(V - E_{\text{K}}) + g_{\text{L}}(V - E_{\text{L}})
\end{equation}
\begin{equation}
  C_m \dot V = -I_m + I_{\text{inj}} + I_{\text{pump}}
\end{equation}
with \code{Cm = 1\,\(\mu\)F/cm\(^2\)}. Stability cap: inner substep
$\le 25\,\mu\text{s}$.

\section{Pumps}

Na$^+$/K$^+$ ATPase: 3\,Na$^+$ out, 2\,K$^+$ in per ATP. SERCA
(Ca$^{2+}$ ATPase): 2\,Ca$^{2+}$ per ATP. PMCA: 1\,Ca$^{2+}$ per ATP.
V-ATPase: $\sim$3\,H$^+$ per ATP.

\chapter{Complete Metabolism Reference}
\label{ch:metabolism-deep}

\file{src/simulation/biochem/Metabolism.h} (962 lines) is the
largest biochem module. This chapter gives the full flux equation
set.

\section{Glycolysis}

Ten steps, collapsed into a single aggregate flux for performance
(individual sub-rate constants are in the file's constexpr table).
Hexokinase is rate-limiting under low-glucose conditions; PFK-1 is
rate-limiting under ATP-rich conditions (AMP allosteric activation).

\begin{equation}
  v_{\text{gly}} = V_{\max}^{\text{gly}}\cdot\frac{[\text{glc}]}{K_m + [\text{glc}]}\cdot\frac{[\text{NAD}^+]}{K_{\text{NAD}} + [\text{NAD}^+]}\cdot f_{\text{PFK}}\,([\text{ATP}],[\text{AMP}])
\end{equation}
with allosteric term
\(f_{\text{PFK}} = 1 / (1 + ([\text{ATP}]/K_i)^4) + \beta \cdot [\text{AMP}]/K_a\).

\section{Pyruvate fate}

Pyruvate branches: LDH (to lactate, reduces NADH $\to$ NAD$^+$) vs.\
PDH (to acetyl-CoA, enters TCA). Ratio governed by
\code{glycolytic} switch and NADH:NAD$^+$ redox state. Warburg cells
push pyruvate to lactate even with O$_2$ present.

\section{TCA cycle}

Eight reactions; key regulatory nodes: citrate synthase (inhibited
by ATP, NADH), isocitrate dehydrogenase (activated by ADP, Ca$^{2+}$;
inhibited by ATP, NADH), $\alpha$-KG dehydrogenase (similar
regulation).

\section{Oxidative phosphorylation}

ETC (Complex I--IV) pumps protons; ATP synthase (Complex V) harvests
the gradient. Effective P/O ratio: 2.5 for NADH-derived electrons,
1.5 for FADH2. Total yield per glucose (under full OxPhos):
$\sim$30 ATP.

\section{Glutaminolysis}

Glutamine $\to$ glutamate (glutaminase) $\to$ $\alpha$-KG (GDH) $\to$
TCA. Important for cancer cells; supplies anaplerotic carbon when
glucose is low.

\section{Pentose phosphate shunt}

Oxidative branch: G6P $\to$ 6-PG $\to$ Ribu5P $\to$ yields NADPH (for
biosynthesis + glutathione regeneration) and 5-C sugars (for
nucleotides). Non-oxidative branch: interconverts sugars to rejoin
glycolysis. Flux rises under oxidative stress.

\section{Biomass composition}

For biomass growth:
\[
\text{Biomass}_{1\,\text{g}} = 0.45\,\text{protein} + 0.25\,\text{lipid} + 0.15\,\text{nucleotide} + 0.10\,\text{carb} + 0.05\,\text{other}.
\]
Each component has a per-gram ATP + NADPH + precursor cost; the
engine sums these into \code{biomassSynth(dt)} which draws from
pools and adds to \code{biomass}.

\chapter{Intracellular Physics --- Full Derivation}
\label{ch:intraphysics-deep}

\section{Langevin derivation}

For a spherical particle of radius $r$ in water at $T$:
\begin{itemize}
  \item Stokes friction coefficient $\gamma = 6\pi\eta r / m$ with
        $\eta = 0.69\,\text{mPa}\cdot\text{s}$ at 37$^\circ$C.
  \item Diffusion coefficient $D = k_B T / (6\pi\eta r)$.
  \item Thermal force spectrum: $\langle F(t)F(t')\rangle =
        2 m\gamma k_B T\,\delta(t-t')$.
\end{itemize}

Discrete Verlet integrator (BBK scheme, Brunger--Brooks--Karplus):
\begin{align}
  \mathbf{v}_{t+dt/2} &= (1 - \gamma\,dt/2)\,\mathbf{v}_t + \frac{\mathbf{F}_t}{2m}dt + \sqrt{\gamma k_B T\,dt / m}\,\boldsymbol{\xi}_t,\\
  \mathbf{x}_{t+dt} &= \mathbf{x}_t + \mathbf{v}_{t+dt/2}\,dt,\\
  \mathbf{v}_{t+dt} &= \frac{1 - \gamma\,dt/2}{1 + \gamma\,dt/2}\,\mathbf{v}_{t+dt/2} + \frac{1}{1 + \gamma\,dt/2}\left(\frac{\mathbf{F}_{t+dt}}{2m}dt + \sqrt{\gamma k_B T\,dt / m}\,\boldsymbol{\xi}'_t\right).
\end{align}

\section{Hertz contact}

When two particles approach within their summed radii by an amount
$\delta > 0$:
\begin{equation}
  F = \frac{4}{3}E^*\sqrt{r^*\delta^3}
\end{equation}
with $E^* = E/(1-\nu^2)$ and $r^* = r_1 r_2/(r_1+r_2)$. The
simulator uses a soft Hertz with $E^*\sim 10\,\text{kPa}$ matching
cortical stiffness of a live cell.

\section{Confinement wall}

The cell membrane is a spherical shell at radius
\code{physCellR = cellSize}. Particles outside the shell feel an
inward Hertz push proportional to their over-extension. This is a
soft wall; a small fraction of overshoot is allowed so the visual
render doesn't look rigid.

\section{Attraction field mechanics}

\begin{lstlisting}
struct AttractionField {
    simd_float3 center;
    float radius;      // cutoff
    float strength;    // force magnitude at r=0
    MolTag attractedTag;
    ParticleType attractedType;
};
\end{lstlisting}
Force applied to matching particles:
\begin{equation}
  \mathbf{F}_{\text{attr}} = -\text{strength}\,\frac{\mathbf{r} - \mathbf{c}}{|\mathbf{r}-\mathbf{c}|+\epsilon}\cdot\max(0, 1 - |\mathbf{r}-\mathbf{c}|/\text{radius}).
\end{equation}

\section{Spatial hashing}

With 100s--1000s of particles and up to $N^2$ pair checks, brute
force is expensive. The \code{SpatialHash} struct bins particles
into a 3D grid with cell size $\sim$2$\times$ the largest particle
radius; pair queries inspect only the 27 adjacent cells. Practical
speed-up: 50--100$\times$ for a typical single-cell scene.

\chapter{Chemistry-First Drug System --- Complete Pipeline}
\label{ch:drug-deep}

This chapter traces the full life-cycle of a drug from user click
to cell-wide effect.

\section{User click in the Drug Lab}

\code{applyDrugVisuals(drugId, conc\_uM)} is the entry point.
Responsibilities:
\begin{enumerate}
  \item Look up the drug in \code{gBioagents.all()}
        (a vector of \code{ChemicalEntity} populated at init from
        \file{data/bioagents/drugs.csv}).
  \item Compute binding affinities against every target in
        \code{gTargets.all()} using
        \code{BindingMatcher::score(drug, target)}.
  \item Register a \code{MediumChemSpec} for the drug's visual
        particles.
  \item Spawn a swarm of particles in the medium with the drug's
        SDF model (from \code{MoleculeCache}).
  \item Push an \code{AppliedDrug} record onto
        \code{Simulation::appliedDrugs}.
\end{enumerate}

\section{AppliedDrug record}

\begin{lstlisting}
struct AppliedDrug {
    int entityIdx;                     // into gBioagents.all()
    float dishConc_uM;                 // extracellular concentration
    std::vector<BindingAffinity> affinityPerTarget;
    std::string id;
};
\end{lstlisting}

\section{Per-tick PK/PD}

\code{updateDrugPK(dt)} runs per substep. The pipeline per drug per
cell:
\begin{enumerate}
  \item Compute permeability \code{perm} from Lipinski descriptors.
  \item Compute flux \code{dC = perm \(\cdot\) (dishConc - cytoConc)
        \(\cdot\) dt\_biosec}.
  \item Add flux to \code{c.drugIntra\_uM[drugId]}.
  \item For every target, mass-action binding with the Avogadro-
        corrected copy-to-\um M conversion.
  \item Fire target modulators with occupancy.
\end{enumerate}

\section{Target modulator catalogue (expanded)}

Each modulator in \file{TargetLibrary.cpp} is one function of shape
\code{void mod(SimCell\& c, float occ, float dt\_biosec)}. Below
is the catalogue with the specific fields touched and the
biological rationale.

\subsection{\code{modCdkB} --- CDK1 inhibitor}

\begin{lstlisting}
c.cdk.CycB = softMul(c.cdk.CycB, 1.0f - 0.95f * occ);
\end{lstlisting}
CycB is phosphorylated by Wee1 in G2 and dephosphorylated by Cdc25
at M-entry. Inhibiting CDK1 blocks entry to M. A well-matched CDK1
inhibitor (e.g.\ RO-3306) therefore produces mitotic arrest at the
G2/M boundary, visible in the cycle stats as accumulating G2 cells.

\subsection{\code{modCdkE} --- CDK2 inhibitor}

\begin{lstlisting}
c.cdk.CycE = softMul(c.cdk.CycE, 1.0f - 0.95f * occ);
c.cdk.CycA = softMul(c.cdk.CycA, 1.0f - 0.90f * occ);
\end{lstlisting}
CDK2 binds Cyclin E in G1/S transition and Cyclin A in S. Inhibition
blocks S-phase entry and slows replication; damage accumulates
because replication forks are more fragile.

\subsection{\code{modTopo2} --- TopII poison}

\begin{lstlisting}
c.program.cdogma.replicationActive = (occ < 0.3f);
c.damageLevel = fminf(2.0f, c.damageLevel + 0.10f * occ * dt_biosec / 3600.0f);
c.ROS         = fminf(100.0f, c.ROS + 2.5f * occ * dt_biosec / 3600.0f);
c.apo.state.p53_active = fminf(1.0f, c.apo.state.p53_active
                              + 0.05f * occ * dt_biosec / 3600.0f);
\end{lstlisting}
TopII traps DNA in a covalent double-strand break. The drug
\emph{stabilises} the break rather than resolving it. Consequences:
fork collapse, DSB accumulation, ROS generation from oxidative
stress, p53 activation, apoptosis. Doxorubicin additionally
generates ROS by redox cycling with iron.

\subsection{\code{modDnaIntercalate} --- cisplatin-like}

\begin{lstlisting}
c.damageLevel = fminf(2.0f, c.damageLevel + 0.20f * occ * dt_biosec / 3600.0f);
c.apo.state.p53_active = fminf(1.0f, c.apo.state.p53_active + 0.10f * occ * dt_biosec / 3600.0f);
c.apo.state.Puma = fminf(1.0f, c.apo.state.Puma + 0.03f * occ * dt_biosec / 3600.0f);
if (c.damageLevel > 1.2f && occ > 0.15f) {
    c.apo.state.Bax_active = fminf(1.0f,
        c.apo.state.Bax_active + 0.04f * occ * dt_biosec / 3600.0f);
    c.apo.state.MOMP_pores = fminf(1.0f,
        c.apo.state.MOMP_pores + 0.02f * occ * dt_biosec / 3600.0f);
}
\end{lstlisting}
Cisplatin forms intrastrand and interstrand DNA crosslinks. At low
damage, NER clears lesions; at saturating damage, p53 transcribes
Puma + Bax, Bax inserts into mitochondria, and MOMP commits. The
\code{> 1.2} branch models this saturation-triggered direct Bax
push, bypassing the Bcl-2 buffering that would otherwise delay
death.

\subsection{\code{modBcl2} --- BH3 mimetics}

\begin{lstlisting}
c.apo.state.Bcl2 = softMul(c.apo.state.Bcl2, 1.0f - 0.90f * occ);
c.apo.state.BclXL = softMul(c.apo.state.BclXL, 1.0f - 0.50f * occ);
\end{lstlisting}
Venetoclax (ABT-199) selectively inhibits Bcl-2; navitoclax
(ABT-263) inhibits Bcl-2/BclXL/Bcl-w. Lowering the anti-apoptotic
pool removes the brake on Bax; cells with a smouldering MOMP
predisposition will commit. In the simulator, cells with
\code{p53\_active > 0}, \code{damageLevel > 0.5}, or
\code{stress > 30} die when Bcl-2 is low.

\subsection{\code{modTubulin} --- taxanes / vinca alkaloids}

Only fires when $c.\text{phase} == 3$ (M). Raises
\code{damageLevel}. The rationale is that taxanes stabilise
microtubules, preventing depolymerisation required for anaphase;
vinca alkaloids do the opposite (prevent polymerisation). Either
way, the SAC cannot satisfy, the cell slips, and emerges from
mitosis with catastrophic damage that drives apoptosis.

\subsection{\code{modMitoCmplxI} --- rotenone / MPP$^+$}

\begin{lstlisting}
c.mitoPotential = fmaxf(40.0f, c.mitoPotential - 60.0f * occ * dt_biosec / 3600.0f);
c.ROS = fminf(100.0f, c.ROS + 1.5f * occ * dt_biosec / 3600.0f);
\end{lstlisting}
Complex-I blockers collapse $\Delta\Psi_m$. ROS rises from back-up
electron leak. Apoptosis engine's \code{mito\_dysfunction} trigger
fires at $\code{mitoPotential} < 80$. The rate is scaled so
saturating rotenone causes MOMP within $\sim$1\,\bioh.

\subsection{\code{modRibosome} --- cycloheximide-like}

\begin{lstlisting}
c.biomass = softMul(c.biomass, 1.0f - 0.0005f * occ);
\end{lstlisting}
Continuous slow biomass decline mimics translation arrest; cells
can't finish a cycle if they can't synthesise cyclins. Quiescence
emerges naturally.

\subsection{\code{modHdac} --- histone deacetylase inhibitors}

\begin{lstlisting}
c.stress = fminf(100.0f, c.stress + 2.0f * occ);
\end{lstlisting}
HDACi's mechanism is complex (global transcriptional reshape +
acetylation of non-histone targets). The simulator applies a
generic stress bump proxy; phase 7M's epigenetic layer would make
this mechanistic by tracking H3K27ac and routing it through the
GRN.

\subsection{\code{modKinaseEgfr}}

Scales CycD down, reducing proliferation. Matches erlotinib,
gefitinib family.

\subsection{\code{modProteasome}}

Raises \code{stress} + \code{Puma}. Bortezomib-like. Rationale: ER
stress from accumulated misfolded proteins drives CHOP and hence
apoptosis.

\subsection{\code{modPol2}}

Scales CycD, CycE down. Matches $\alpha$-amanitin,
flavopiridol at extreme doses.

\section{Polypharmacology emergence}

Because one drug can bind multiple targets with different
affinities (a \emph{dirty} kinase inhibitor), applying
staurosporine registers non-zero occupancy on
\code{TGT\_KINASE\_EGFR}, \code{TGT\_CDK1}, \code{TGT\_CDK2}, and
others simultaneously. Each modulator fires; the cell feels every
effect proportional to its occupancy. The phenotype --- massive
cycle arrest + proliferation collapse + generalised apoptosis ---
emerges from physics, not from a
``\code{MOA\_PAN\_KINASE\_INHIBITOR}'' flag.

\chapter{Pathogen Infection --- Complete Nine-Stage Walkthrough}
\label{ch:infection-deep}

This chapter is an annotated trace of a single virion from free
drift to daughter-virion egress.

\section{Stage 1: Free drift}

A virion at position $\mathbf{x}_t$ is advanced by:
\begin{equation}
  \mathbf{x}_{t+dt} = \mathbf{x}_t + \sigma\sqrt{dt}\,\boldsymbol{\xi}
\end{equation}
with $\sigma = 0.40\,\wu\sqrt{\biosec}$ and $\boldsymbol{\xi}$ a
2D Gaussian (Box--Muller). The step length $\sigma\sqrt{dt}$ is
order $0.7\,\wu$ at a typical physics dt ($3.3\,\biosec$ scaled),
which means a virion's RMS displacement crosses the 6\,\wu contact
radius within $\sim$10 ticks. This is biologically reasonable: a
100-nm virion in water has $D\approx 4\,\um^2/\text{s} =
0.16\,\wu^2/\text{s}$, so in 3.3\,\biosec it rms-moves
$\sqrt{2\cdot 0.16\cdot 3.3}\approx 1.03\,\wu$ per axis.

\section{Stage 2: Proximity test}

On each tick, each free virion's nearest live cell is found with a
linear scan (OK for $\sim 10$ cells in test mode, needs a spatial
hash in dense colony mode). If the squared distance to the nearest
cell is $< 36 = 6^2\,\wu^2$, we proceed to scoring.

\section{Stage 3: Binding score}

\code{virionBindScore(vs, rcvIdOut)} returns a score in [0,1]. If
$< 0.35$, the virion bounces (no productive contact). Otherwise:

\begin{equation}
  k_{\text{eff}} = 10^{3(\text{score}-0.7)}\,\text{s}^{-1}
\end{equation}
\begin{equation}
  \frac{d\,\text{boundSpikes}}{dt} = k_{\text{eff}}\cdot(\text{capacity} - \text{boundSpikes})
\end{equation}

\section{Stage 4: BOUND transition}

When \code{boundSpikes >= max(20, 0.25 * capacity)}, the virion
transitions FREE $\to$ BOUND, its \code{hostCellIdx} is set, and
\code{stageTimer\_s} resets. In single-molecule flu experiments
(Rust 2004), $\sim$10--20 engaged HA spikes are sufficient for
productive entry; the floor of 20 or $25\%$ of the capacity reflects
this.

\section{Stage 5: ENTERING}

After a 60\,\biosec dwell the virion transitions BOUND $\to$
ENTERING. Biologically this is clathrin-coated-pit maturation + pit
scission (phase 8C upgrades this to an explicit VesicularTraffic
event).

\section{Stage 6: UNCOATING and genome deposit}

After an additional 300\,\biosec dwell (endosomal acidification and
fusion, phase 8D upgrades), the virion increments
\code{host.intraVirions[specIdx]} by 1 and transitions
UNCOATING (which is then culled from the free pool).

\section{Stage 7: Replication}

Current toy:
\begin{equation}
  \frac{d\,\text{intraVirions}}{dt} = r\cdot\text{intraVirions}
\end{equation}
with $r = \text{replicationRate\_per\_s}$. Phase 8E replaces this
with real transcription through the Central Dogma's gene registry
(Ch.~\ref{ch:phase8a}).

\section{Stage 8: Assembly}

When \code{intraVirions > assemblyThreshold}, a portion
(0.2\%/\biosec) moves from \code{intraVirions} to
\code{assembledVirions}. Phase 8F replaces this with stoichiometric
gating over \code{viralProteinCount}.

\section{Stage 9: Budding and egress}

Once $\code{assembledVirions} > 1$, per-tick output:
\begin{equation}
  \text{bud} = \text{budRate\_per\_s}\cdot\text{assembledVirions}\cdot dt.
\end{equation}
Integer portion spawns free virions at the cell surface (ring of
radius $1.05\cdot\text{radius}$). Capped at 10 spawns per tick to
keep the vector manageable; total free virions capped at 3\,000
globally.

\section{Apoptotic spill}

On BODIES transition of the host apoptosis phase,
\code{releasePathogens(c)} spawns outside the cell radius annulus
every intracellular virion and bacterium at once, then zeroes the
intracellular pools. This is the mechanism by which one dying cell
seeds the next wave of infection --- a true physics-driven
infection chain.

\chapter{Immune System --- Inventory and Wiring}
\label{ch:immunity-deep}

The immune scaffolding in \file{Immunity.h} and \file{Pathogens.h}
covers most textbook-canonical mechanisms. This chapter tours each
piece, notes its current status, and identifies the wiring gaps
that phase 7I of the plan fills.

\section{Toll-like receptors}

Seven TLRs modelled: TLR2 (peptidoglycan, lipoproteins), TLR3
(dsRNA, endosomal), TLR4 (LPS), TLR5 (flagellin), TLR7/8 (ssRNA,
endosomal), TLR9 (CpG DNA, endosomal). Each has a surface density
scalar. Activation threshold via Hill function of cognate PAMP
concentration in the MediumField.

\section{Cytosolic sensors}

RIG-I (non-self dsRNA), MDA5 (long dsRNA), NOD2 (muramyl dipeptide),
AIM2 (cytosolic dsDNA), cGAS-STING (cytosolic dsDNA via cGAMP
second messenger). Not yet wired but scaffolded.

\section{Inflammatory output}

NF-$\kappa$B is the shared output node; active NF-$\kappa$B
transcribes TNF-$\alpha$, IL-1$\beta$, IL-6 into a cell-local
cytokine pool that then diffuses through the MediumField.

\section{Complement cascade}

Classical (antibody-triggered), lectin (mannose-triggered),
alternative (spontaneous C3b tick-over). Terminal common path:
C5b-C9 $\to$ MAC (membrane attack complex) on pathogen membrane.
Scaffolded but not active.

\section{Phagocytosis}

Macrophage and neutrophil cells (when instantiated as SimCell
subtypes) engulf opsonised pathogens. Respiratory burst: NOX2
assembly, superoxide $\to$ H$_2$O$_2$ $\to$ HOCl; cytotoxic to
pathogens inside phagolysosomes.

\section{NK cells}

Missing-self detection via MHC-I. Active NK kills via perforin
pore + granzyme B delivery. Triggers caspase-3 directly,
bypassing mitochondrial amplification.

\section{T cells}

TCR-peptide-MHC recognition is scored via the same BindingMatcher
used by drugs + pathogens. A random V(D)J repertoire can be
pre-generated; the first T cell whose TCR scores $> 0.7$ against a
presented viral peptide becomes the clonal winner.

\section{B cells}

Similar to T cells but for antibodies vs.\ free virion spike
proteins. Antibodies become floating \code{ChemicalEntity} in the
MediumField and competitively inhibit virion-receptor binding.

\chapter{Rendering --- From Data to Pixels}
\label{ch:rendering-deep}

\section{Frame sequence}

One frame of \code{CellSim} proceeds as:
\begin{enumerate}
  \item Process GLFW input (keyboard, mouse, scroll).
  \item Advance \code{gSim.update(dt)}.
  \item Update mitosis visualisation (single-cell mode).
  \item Update apoptotic bodies; efferocytosis; DAMP events.
  \item Update medium chemicals (drift, binding, uptake).
  \item Rebuild \code{gCellInstances} for shell rendering
        (\code{syncCellInstances}).
  \item In single-cell mode, rebuild intracellular atom buffer
        (\code{uploadCellInterior}). In test-light mode, build a
        minimal atom buffer of pathogens + apoptotic bodies +
        receptor dots.
  \item Submit Metal command buffer:
    \begin{enumerate}
      \item Substrate floor pipeline.
      \item Fluid (ray-marched dish).
      \item GLB organelle pipeline (only for nearby cells).
      \item Cell shell pipeline (\code{CellRender.metal}).
      \item Molecule atom pipeline (pathogens, bodies, intra).
      \item Molecule bond pipeline.
      \item ImGui overlay.
    \end{enumerate}
  \item Present drawable.
\end{enumerate}

\section{Cell shell shader key equations}

Fresnel rim intensity:
\begin{equation}
  I_{\text{rim}} = \text{pow}(1 - |\mathbf{N}\cdot\mathbf{V}|, 3)\cdot\text{glowIntensity}\cdot\text{color}.
\end{equation}
Body glow (Lambertian $\cdot$ forward term):
\begin{equation}
  I_{\text{body}} = \max(0, \mathbf{N}\cdot\mathbf{L})\cdot\text{color}.
\end{equation}
Furrow deformation: a radial pinch at $\theta = \pi/2$ with
amplitude \code{furrowDepth}. Apoptosis bleb: radial perturbation
$r'(t) = r\cdot(1 + A\sin(\omega t + \phi))$ with
$A = \code{BLEB\_AMPLITUDE\_FRAC} \cdot \text{apoProgress}$.

\section{Fluid rendering}

Dish fluid uses ray-marching. Each pixel's ray intersects a
truncated sphere defined by SCENE\_BOUND; inside the fluid a
distance field accumulates along the ray. Wave deformation comes
from a Perlin-like noise time-dependent modulation of the surface
height.

\chapter{UI --- Panel Inventory and Rationale}
\label{ch:ui-deep}

\section{Setup overlay}

Modal at startup or upon \texttt{R} keypress. Presents seven preset
templates (plus \emph{Test Mode (Lite 16GB)} added 2026-04-20).
Each template sets every gSetup field. The user can then adjust
sliders before clicking \emph{Start simulation}.

\section{Cell Mode panel}

Top-centre control strip. Shows cell count, \emph{Follow cell}
toggle, \emph{Solo focus} toggle, zoom readout (in follow mode),
cell-index slider (when more than one cell).

\section{Population panels (colony mode)}

\begin{itemize}
  \item \textbf{Population Stats}: alive / G1 / S / G2 / M /
        quiescent / apoptotic / necrotic / glycolytic / divisions /
        deaths.
  \item \textbf{Population Dynamics}: historical curve plot of
        divisions / deaths / alive.
  \item \textbf{Culture Medium (closed dish)}: current mM /
        concentration of every medium species; mass-balance drift
        indicator.
\end{itemize}

\section{Per-cell panels (single-cell mode)}

\begin{itemize}
  \item \textbf{Cell Research Panel}: every scalar field of the
        focused cell --- cycle timer, CDK readouts, biomass, ATP,
        damage, ROS, stress, adhesion strength.
  \item \textbf{Molecules Inside Cell}: live count of each
        \code{ParticleType} inside the interior, with jobs.
  \item \textbf{Mitosis Debug}: per-phase readouts, SAC status,
        chromosome alignment, centrosome positions.
\end{itemize}

\section{Chemistry panels}

\begin{itemize}
  \item \textbf{Drug Lab (chemistry-driven)}: compound dropdown,
        descriptor readout (MW, logP, TPSA, HBD, HBA, rings), log-
        concentration slider, Apply Uniformly button, top-
        affinities table.
  \item \textbf{Pathogen Lab (virion + bacterium)}: species
        dropdowns, best-match receptor, score, Inoculate button,
        free/bound/intra counters.
  \item \textbf{Test Actions} (light mode only): one-click buttons
        for every interaction scenario, plus cell-state shortcuts
        (force apoptosis, heal, clear).
\end{itemize}

\chapter{Export System --- Bundle Contents}
\label{ch:export-deep}

\code{exportAllData(dir)} writes six files in a timestamped
sub-folder. Each file is documented below.

\section{manifest.txt}

\begin{verbatim}
Run ID: 20260419_234354
Start: 2026-04-19 23:43:54
End:   2026-04-20 00:52:17
Duration bio-h: 12.3
Duration real-s: 4103
Cells alive at end: 377
Divisions: 252
Deaths: 0
Mass balance drift: 0.00034 mM
Reference: docs/equations/cellsim_equations.tex
\end{verbatim}

\section{population.csv}

Columns:
\code{wall\_sec, bio\_sec, bio\_hours, cell\_count, g1, s, g2, m,
divisions, deaths, cdk\_CycD, cdk\_CycE, cdk\_CycA, cdk\_CycB,
cdk\_Rb, cdk\_p21, biomass, ATP, damage, replProgress}.
Sampled once per bio-minute (60 bio-s).

\section{cells.csv}

One row per cell at export time; every scalar field of
\code{SimCell}. Useful for per-cell heterogeneity analysis.

\section{medium\_field.csv}

64\(\times\)64 grid \(\times\) 12 species. Long format:
\texttt{x, z, species\_id, concentration\_mM}.

\section{apo\_bodies.csv}

Ledger of every apoptotic body currently in flight: position,
kind, remaining biomass / membrane / receptor, phase.

\section{medium\_chems.csv}

Current \code{MediumChemical} swarm: species, target cell, state,
position, binding progress.

\chapter{Headless Test Suite --- Invocation Reference}
\label{ch:tests-deep}

\section{Invocation patterns}

\begin{verbatim}
./build/cellsim_headless                 # defaults: 72 bio-h, 60 s/step, 5 cells
./build/cellsim_headless 46 60 125       # CTC seq02 reference run
./build/cellsim_headless 120 30 800      # long-run dense-colony stress test
\end{verbatim}

\section{Expected outputs}

\begin{itemize}
  \item \file{logs/headless\_doubling.csv}: time-series.
  \item \file{logs/compare\_vs\_reference.csv}: per-row error.
  \item \file{logs/mitosis\_<session>.log}: per-mitosis events.
  \item \file{logs/diag\_<session>.log}: per-120-frame diagnostic.
  \item \file{logs/centrosome\_<session>.log}: centrosome tracking.
\end{itemize}

\section{Pass criteria}

\begin{itemize}
  \item \code{cellsim\_headless 46 60 125}: mean $|\text{rel
        err}| \le 6\%$.
  \item \code{cellsim\_infection}: all four assertions pass.
\end{itemize}

\chapter{Build System --- CMake Walkthrough}
\label{ch:cmake-deep}

\section{Top-level structure}

\file{CMakeLists.txt} declares:
\begin{itemize}
  \item \texttt{project(CellSim CXX OBJC OBJCXX)}.
  \item C++20 standard.
  \item vcpkg toolchain file for third-party deps (GLFW, ImGui,
        simd).
  \item Apple's Metal framework (MetalKit, Foundation,
        QuartzCore).
  \item Three targets: \code{CellSim}, \code{cellsim\_headless},
        \code{cellsim\_infection}.
\end{itemize}

\section{Warning discipline}

Every target compiles with \code{-Wall -Wextra
-Wno-unused-parameter -O2}. Unused-variable and unused-function
warnings are tolerated but flagged; the diagnostic log lists them
so a cleanup pass can catch regressions.

\section{Metal shader compilation}

Metal shaders are compiled from source at startup (not
pre-compiled). This is for portability: a user on a fresh install
doesn't need \code{xcrun metal} in their path. The downside is
$\sim$50\,ms of startup latency, which is acceptable.

\chapter{Calibration --- The 5.1\% Story}
\label{ch:calibration-deep}

The simulator is calibrated against one primary reference: CTC
Fluo-N2DL-HeLa seq02, a 46\,\bioh growth curve of HeLa cells
starting from 125 at $t=0$.

\section{What was calibrated}

\begin{itemize}
  \item \code{SLOW\_DT\_SCALE}: controls metabolism and PK rates.
  \item \code{MEDIUM\_DT\_SCALE}: controls cycle and fate rates.
  \item \code{MECH\_P21\_COUPLING}: controls crowding response.
\end{itemize}

\section{What was not calibrated}

\begin{itemize}
  \item CDK/cyclin ODE rates (Ch.~\ref{ch:cycle-deep}): these are
        from the Aldridge 2006 ballpark; they could be refit but
        contributed negligibly to the 5.1\% error so were left
        alone.
  \item Apoptosis cascade rates: calibrated separately against Ye
        2017 + Spencer 2009 microfluidic single-cell data. Error
        bar not separately measured; qualitative fit only.
  \item Central Dogma elongation rates: from Singh 2013; chosen so
        HBB protein appears within $\sim 5\,\biomin$ of G1 entry.
\end{itemize}

\section{Calibration procedure}

\begin{enumerate}
  \item Run \code{./scripts/calibrate\_rate.sh} which sweeps
        (\code{SLOW\_DT\_SCALE}, \code{MEDIUM\_DT\_SCALE}) on a
        5$\times$5 grid.
  \item For each cell, run 5 independent seeds and average.
  \item Metric: mean $|\text{relative error}|$ vs CTC seq02 over
        92 matched timepoints.
  \item Take the argmin. Before Phase 8A this was
        (\code{0.055}, \code{0.20096}) yielding 5.1\%.
\end{enumerate}

\section{Reproducibility across seeds}

All five seeds come within $\pm 0.2\%$ of 5.1\%. This is mostly
because the colony-level averaging smooths out single-cell
stochasticity; individual cell traces vary considerably.

\chapter{Phase 8A --- Implementation Details}
\label{ch:phase8a-deep}

This chapter documents the April 2026 multi-gene Central Dogma
upgrade at the code level. It is intended to be read alongside
\file{src/simulation/CentralDogma.h}.

\section{GeneTemplate struct (line 278--306)}

\begin{lstlisting}
enum class GenePolymerase : uint8_t {
    HOST_POLII=0, VIRAL_RdRp=1, VIRAL_RT=2, VIRAL_POL=3
};

struct GeneTemplate {
    std::string     id;
    std::string     dnaSeq;
    int             startCodon  = -1;
    int             stopCodon   = -1;
    std::string     codingMRNA;
    float           promoterStrength = 1.0f;
    GenePolymerase  polymerase  = GenePolymerase::HOST_POLII;
    std::string     proteinName;
    int             copiesMade  = 0;
};
\end{lstlisting}

Units and semantics:
\begin{itemize}
  \item \code{dnaSeq}: sense-strand DNA; 4-letter \texttt{\{A,C,G,T\}}.
  \item \code{startCodon, stopCodon}: indices into \code{codingMRNA}
        after it's built; $-1$ if unknown (will be auto-filled at
        registration time).
  \item \code{codingMRNA}: cached at registration; avoids rebuilding
        per transcription.
  \item \code{promoterStrength}: weight in the random gene pick; HBB
        is 1.0 by convention; a strong viral promoter is 2--5, a
        silenced gene 0.1.
  \item \code{polymerase}: which enzyme transcribes this; only
        \code{HOST\_POLII} is wired today, the others are reserved.
  \item \code{proteinName}: tag propagated to
        \code{ProteinProductState::proteinName} and
        \code{viralProteinCount}.
  \item \code{copiesMade}: monitoring counter, incremented per
        released protein.
\end{itemize}

\section{CentralDogmaState additions}

\begin{lstlisting}
std::vector<GeneTemplate> viralGenes;
std::unordered_map<std::string, int> viralProteinCount;
\end{lstlisting}

\section{Helpers (line 738--794)}

\begin{lstlisting}
inline const char* geneSequence(int geneIndex, int& outLen) const {
    if (geneIndex < 0 || geneIndex >= (int)viralGenes.size()) {
        outLen = HBB_LENGTH;
        return HBB_SEQUENCE;
    }
    outLen = (int)viralGenes[geneIndex].dnaSeq.size();
    return viralGenes[geneIndex].dnaSeq.c_str();
}

int pickGeneForTranscription() {
    if (viralGenes.empty()) return -1;
    float total = 1.0f;
    for (const auto& g : viralGenes) total += g.promoterStrength;
    float r = ((float)rand() / (float)RAND_MAX) * total;
    if (r < 1.0f) return -1;
    r -= 1.0f;
    for (int i = 0; i < (int)viralGenes.size(); i++) {
        r -= viralGenes[i].promoterStrength;
        if (r <= 0.0f) return i;
    }
    return (int)viralGenes.size() - 1;
}

int registerViralGene(const GeneTemplate& gene) {
    viralGenes.push_back(gene);
    GeneTemplate& g = viralGenes.back();
    if (g.codingMRNA.empty()) {
        g.codingMRNA.reserve(g.dnaSeq.size());
        for (char c : g.dnaSeq)
            g.codingMRNA.push_back(transcribe(c));
    }
    if (g.startCodon < 0)
        g.startCodon = findStartCodon(g.codingMRNA);
    if (g.stopCodon < 0 && g.startCodon >= 0)
        g.stopCodon = findStopCodon(g.codingMRNA, g.startCodon);
    return (int)viralGenes.size() - 1;
}
\end{lstlisting}

\section{Transcription elongation (line 1329+)}

\begin{lstlisting}
int seqLen = 0;
const char* seqData = geneSequence(tx->geneIndex, seqLen);
ts.rnaPolPosition = fminf(1.0f, ts.rnaPolPosition + dt * 0.22f);
ts.currentBP = std::min(seqLen, (int)floorf(ts.rnaPolPosition * (float)seqLen));
while (tx->genomicBases < ts.currentBP && tx->genomicBases < seqLen) {
    tx->preMRNA.push_back(transcribe(seqData[tx->genomicBases]));
    tx->genomicBases++;
}
tx->transcriptionProgress = ts.rnaPolPosition;
if (!tx->capped && tx->genomicBases >= 18) tx->capped = true;

if (ts.currentBP >= seqLen || ts.rnaPolPosition >= 1.0f) {
    tx->transcriptionComplete = true;
    if (tx->geneIndex >= 0 && tx->geneIndex < (int)viralGenes.size()) {
        // Viral genes skip splicing (RNA viruses are single-exon).
        tx->spliced = true;
        tx->polyadenylated = true;
        tx->matureMRNA = viralGenes[tx->geneIndex].codingMRNA;
        tx->polyATail.assign(48, 'A');
        tx->intron1Removed = true;
        tx->intron2Removed = true;
        tx->processingProgress = 1.0f;
    } else {
        bindSpliceosomeIfNeeded(ts.transcriptIndex);
    }
    clearTranscriptionState(ts);
    ts.timer = 0.0f;
}
\end{lstlisting}

\section{Translation codon priming (line 834+)}

\begin{lstlisting}
void primeCurrentCodon(TranslationState& tr) {
    tr.codon.clear(); tr.anticodon.clear();
    tr.incomingAA = '?'; tr.stopReached = false;
    tr.codonProgress = 0.0f; tr.waitingForTRNA = true;
    const auto* tx = getTranscript(tr.transcriptIndex);
    if (!tx || tx->matureMRNA.empty()) { tr.stopReached = true; return; }
    int txStart = (tx->transcriptStartCodon >= 0)
                  ? tx->transcriptStartCodon : startCodonBase;
    int txStop  = (tx->transcriptStopCodon  >= 0)
                  ? tx->transcriptStopCodon  : stopCodonBase;
    if (txStart < 0) { tr.stopReached = true; return; }
    int base = txStart + tr.currentCodon * 3;
    if (base + 2 >= (int)tx->matureMRNA.size()) {
        tr.stopReached = true; return;
    }
    tr.codon = tx->matureMRNA.substr(base, 3);
    AminoAcid aa = translateCodon(tr.codon[0], tr.codon[1], tr.codon[2]);
    if (aa.letter == '*') { tr.stopReached = true; return; }
    tr.anticodon = anticodonForCodon(tr.codon[0], tr.codon[1], tr.codon[2]);
    tr.incomingAA = aa.letter;
    int totalCodons = (txStop > txStart) ? ((txStop - txStart) / 3)
                                          : totalCodingCodons();
    tr.ribosomePosition = (totalCodons > 0)
        ? ((float)(tr.currentCodon + 1) / (float)totalCodons) : 0.0f;
}
\end{lstlisting}

\section{Protein release (line 900+)}

\begin{lstlisting}
void releaseProtein(const TranslationState& tr) {
    if (tr.nascentPeptide.empty()) return;
    int proteinSlot = findFreeProteinSlot();
    if (proteinSlot < 0) return;
    auto& protein = proteins[proteinSlot];
    clearProteinState(protein);
    protein.active = true;
    protein.sourceTranscript = tr.transcriptIndex;
    protein.aminoAcids = tr.nascentPeptide;
    const auto* tx = getTranscript(tr.transcriptIndex);
    if (tx) {
        protein.geneIndex = tx->geneIndex;
        protein.geneId    = tx->geneId;
        protein.proteinName = tx->proteinTag;
        if (protein.geneIndex >= 0) {
            viralProteinCount[tx->proteinTag]++;
            if (protein.geneIndex < (int)viralGenes.size())
                viralGenes[protein.geneIndex].copiesMade++;
        }
    }
    if ((int)tr.nascentPeptide.size() >= 120 &&
        (int)tr.nascentPeptide.size() <= 170) {
        protein.structureAsset = "hbb_folded.pdb";
    }
}
\end{lstlisting}

\section{init() reset}

\begin{lstlisting}
void init() {
    activeMRNAs = 0;
    nextTranscriptUid = 1;
    codingMRNA = buildHBBCodingMRNA();
    startCodonBase = findStartCodon(codingMRNA);
    stopCodonBase = findStopCodon(codingMRNA, startCodonBase);
    viralGenes.clear();
    viralProteinCount.clear();
    // ... rest unchanged
}
\end{lstlisting}

\section{Why HBB's fast path is preserved}

\code{geneSequence(-1, outLen)} returns the pointer \code{HBB\_SEQUENCE}
with zero allocations. \code{transcribe(seqData[tx->genomicBases])}
is a direct read of the character literal. The compiler inlines
this to roughly the same assembly as the pre-8A code:
\begin{verbatim}
  movzbl (%rdi,%rax,1), %eax    # load one byte
  ...
\end{verbatim}
There is a small extra cost at the comparison against
\code{HBB\_LENGTH} because the length lookup indirects through
\code{viralGenes.size()}, but this is dominated by the memory
load anyway. Benchmark: pre-8A and post-8A headless runs complete
in the same 12\,$\pm$\,0.2\,s wall time.

\chapter{Planned Phase 8B-L --- What Each Session Does}
\label{ch:phase8-roadmap}

This chapter summarises the remaining Part 8 work in sequence. It
is derived from \file{/Users/henry/.claude/plans/make-this-plan-into-virtual-key.md}
sections \S54--\S62.

\paragraph{8B (1 session).} Receptor IntraParticles on cell membrane.
At cell init, spawn N receptor particles per class from the
CellType catalogue; these replace the Fibonacci shell decoration
with physics-real particles.

\paragraph{8C (2 sessions).} Clathrin-coated pit nucleation +
vesicle scission. Virion binds $\geq 5$ receptor particles $\to$
pit forms $\to$ \code{VesicularTraffic::clathrin\_available--}
$\to$ vesicle scissions $\to$ virion enters EE compartment.

\paragraph{8D (2 sessions).} V-ATPase endosome acidification
+ pH-triggered fusion. Explicit per-endosome pH field;
\code{V-ATPase} step acidifies using ATP; below a virion-specific
fusionPh threshold the envelope fuses with the endosome membrane
and the genome is released to cytosol.

\paragraph{8E (3 sessions).} Use 8A gene registry. When a virion
uncoats, call \code{cdogma.registerViralGene(...)}. Host RNA Pol II
competes; viral mRNAs flow through existing translation; protein
counts accumulate in \code{viralProteinCount}.

\paragraph{8F (1 session).} Stoichiometric assembly + envelope
lipid conservation. Assembly is gated on
\code{viralProteinCount[capsid] $\geq$ 180 \&\&
viralProteinCount[spike] $\geq$ 500 \&\& vRNP $\geq$ 8}; one
decrement per virion. Envelope lipid decremented from
\code{SimCell::membraneMass\_bm}.

\paragraph{8G (2 sessions).} Caspar-Klug capsid mesh. Procedural
icosahedral tiling at startup from VirionSpec's T-number. Receptor
rendering swaps from Fibonacci dots to the real IntraParticle path.

\paragraph{8H (2 sessions).} T3SS needle-complex + effector
injection for bacterial invasion. Bacterium's adhesin binds host
integrin $\to$ dwell $\to$ T3SS assembles $\to$ effector
ChemicalEntities delivered to host cytosol $\to$ bind host targets
via BindingMatcher.

\paragraph{8I (1 session).} Procedural bacterium mesh with real
wall layers: inner membrane, periplasm, outer membrane, LPS fringe
(gram-neg); thick peptidoglycan (gram-pos); flagella helix with
rotation; pili.

\paragraph{8J (1 session).} Conjugation + plasmid transfer. Two
bacteria's pili touch $\to$ plasmid copied $\to$ resistance
transferred; antibiotic-resistance genes modulate BindingMatcher
scores.

\paragraph{8K (1 session).} Immune response. Viral mRNA
accumulation activates RIG-I $\to$ Type-I IFN output to medium
$\to$ neighbouring cells enter antiviral state (translation
init slowed).

\paragraph{8L (1 session).} End-to-end verification harness.
Regression test that a single virion traces every stage with
mass-balance and event-order assertions.

Total remaining Part 8 work: $\sim$16 sessions. Phase 8A is
complete as of 2026-04-20.

\chapter{Closing Notes --- Debugging Tips and Pitfalls}
\label{ch:tips}

This chapter collects hard-won lessons from actual debugging
sessions during development.

\section{Mass-balance drift debugging}

If \code{MediumField::checkBalance()} reports drift $> 10^{-3}$:
\begin{enumerate}
  \item Find the most recent code that moved mass between pools.
  \item Check that every \code{pool -= delta} has a matching
        \code{exchange(...)} or \code{pool\_other += delta}.
  \item Running the headless validator in single-species mode
        (turn off all but MS\_AA\_POOL) isolates which species
        drifts.
\end{enumerate}

\section{The 550k-bacteria explosion}

During the Phase 7 test, extracellular bacterial division had no
density cap. One click on \emph{Release E. coli} seeded 20
bacteria, which doubled every 20\,\biomin, filling the medium with
$550\,000$ bacteria within minutes. The fix is a hard global cap
\code{MAX\_FREE\_BACTERIA = 2000}. Lesson: any unbounded
growth term needs a cap, even if the ``natural'' biology says it's
fine --- the simulator is running compressed time.

\section{All-cells-in-PRIMED filter bug}

During Phase 7 integration we gated virion binding on
\code{apoPhase == Apoptosis::ALIVE}. But a freshly-initialised cell
often has baseline caspase activity that puts it in PRIMED within
the first few ticks. The binding filter then never matched any
cell. The fix is to gate on
\code{apoPhase < Apoptosis::FRAGMENTATION} --- cells in PRIMED,
MOMP, and EXECUTION are still viable hosts.

\section{The Avogadro-correction story}

An early drug-binding bug had targets saturated at
$10^{-10}\,\um\text{M}$ of drug. Traced to a factor of
$6 \cdot 10^{17}$ in the copy-to-\um M conversion that should have
been $2.41\cdot 10^6$ --- the factor of $10^{-11}$ is the cell
volume in L. Lesson: when you're converting molecule counts to
concentrations, always include the cell volume explicitly.

\section{Why the infection test hung}

The infection smoke test originally had no virion drift --- so
virions sat at their initial seed position forever, never reaching
a cell. The Brownian step added a random-walk with proper
$\sigma\sqrt{dt}$ scaling. Lesson: check that every particle
physics step advances the particle by some non-zero random amount
on average.

\section{The ``all cells seem static even at 100x'' report}

In test/light mode, the cells do not move much because in colony
mode with 1 cell there is no cell-cell interaction forcing motion.
The visual illusion of ``static'' is real. Lesson: biological cells
in a petri dish don't move much; the visible dynamics in the
simulator are the pathogens, drugs, medium flows, and the cell's
internal structure (in single-cell mode).

\section{When to use Write vs Edit}

When authoring large blocks of content into a file, \code{Write}
blasts the whole file in one tool call. \code{Edit} surgically
changes a single block. The trade-off: if you're writing a 4\,000-
line file that doesn't exist yet, \code{Write} is one call; if you
need to append a chapter, \code{Edit} is one call.

\chapter{A Final Mental Model}
\label{ch:mental-model}

If a reader takes only one thing away from this book, let it be
this mental model of the simulator:

\textbf{Cells are state bags. State bags are updated by named
functions. Every name of a function is a biological process.
Every constant in \file{Constants.h} is a literature-grounded
number. Every drug / virus / bacterium is a structure, not a
label.}

That discipline is what separates \textsc{CellSim} from a
prescription-style ``mechanism of action'' simulator. It is why
applying a drug the engine has never seen produces a non-trivial
phenotype. It is why a scrambled viral spike fails to bind. It is
why the infection chain propagates from a first lysed cell to
dozens more without any scripted trigger.

Everything else --- the Metal shaders, the ImGui panels, the
LaTeX equations reference, the build system --- is in service of
making that discipline visible, auditable, and extensible.

\paragraph{Quod erat demonstrandum.} The engine is real biology
expressed in C\texttt{++}, rendered in Metal, grounded in
literature, anchored by a single mass-balance invariant, and
calibrated to 5.1\% against real HeLa microscopy.

% ═════════════════════════════════════════════════════════════════════
\part{Extended Reference}
% ═════════════════════════════════════════════════════════════════════

\chapter{Every Constant in \texttt{Constants.h}}
\label{ch:every-constant}

This chapter tabulates every single \code{constexpr} in
\file{src/core/Constants.h}. Each row: name, value, units,
subsystem, literature source. Where multiple constants share a
subsystem they are grouped. Reading this chapter cover to cover
constitutes a full audit of the simulator's anchor to external
biology.

\section{Geometry and world anchoring}

\begin{longtable}{p{4.6cm}rrp{6.0cm}}
\toprule
\textbf{Constant} & \textbf{Value} & \textbf{Units} & \textbf{Reference / rationale}\\ \midrule
\code{MICROMETERS\_PER\_WORLD\_UNIT} & 5.0 & \um/\wu & Anchor: 1\,\wu = 5\,\um.\\
\code{UM\_TO\_WU}             & 0.2 & \wu/\um & Derived: $1/\code{MICROMETERS\_PER\_WORLD\_UNIT}$.\\
\code{CELL\_RADIUS\_BASE}     & 2.0 & \wu    & HeLa median radius 10\,\um (BNID 100434).\\
\code{FLOOR\_Y}               &-7.0 & \wu    & Dish-floor $y$ coordinate.\\
\code{SCENE\_BOUND}           & 55.0 & \wu   & Half-extent; $\pm 275\,\um$ rectangular FOV.\\
\code{MAX\_CELLS}             & 2000 & --    & Memory cap; sizeof(SimCell) * 2000 $\approx$ 4\,MB.\\
\code{INIT\_CELLS}            & 5   & --    & Colony default start count.\\
\code{SINGLE\_CELL\_COUNT}    & 1   & --    & Single-cell default.\\
\code{MONOLAYER\_CAPACITY}    & 600 & --    & Soft cap for crowding response.\\
\code{MITO\_COUNT\_G1}        & 3   & --    & Mitochondria per cell in G1.\\
\code{MITO\_COUNT\_G2}        & 6   & --    & Post-fission (Alberts Ch 14).\\
\code{GOLGI\_COUNT\_G1}       & 1   & --    & \\
\code{GOLGI\_COUNT\_G2}       & 2   & --    & Golgi fragments for mitosis.\\
\code{NUTRIENT\_GRID\_N}      & 64  & --    & Medium grid resolution; voxel $\sim$8.6\,\um.\\
\bottomrule
\end{longtable}

\section{Time and integration}

\begin{longtable}{p{4.6cm}rrp{6.0cm}}
\toprule
\textbf{Constant} & \textbf{Value} & \textbf{Units} & \textbf{Rationale}\\ \midrule
\code{BIO\_TIME\_SCALE}         & 180.0   & bio-s/real-s & 1 real-s $\to$ 3\,\biomin simulated.\\
\code{BIO\_MIN\_PER\_SEC}       & 3.0     & bio-min/real-s & Legacy alias.\\
\code{FAST\_DT\_SCALE}          & 1.0     & --           & Fast dynamics (receptor kinetics).\\
\code{MEDIUM\_DT\_SCALE}        & 0.20096 & --           & Cycle / fate; CTC-calibrated.\\
\code{SLOW\_DT\_SCALE}          & 0.055   & --           & Metabolism, drug PK; CTC-calibrated.\\
\code{CDK\_DT\_SCALE}           & 0.8     & --           & CDK ODE stiffness control.\\
\code{LAG\_DURATION}            & 0.3     & slow-time    & Legacy slow-time unit.\\
\code{LAG\_BIOSEC}              & 1800    & \biosec      & 30\,\biomin G1 entry lag.\\
\code{MAX\_SUBSTEP\_DT}         & 0.05    & real-s       & Per-sub-step cap.\\
\bottomrule
\end{longtable}

\section{Diffusion coefficients}

\begin{longtable}{lrl}
\toprule
Species & $D_\text{grid}$ & Rationale\\ \midrule
\code{DIFF\_O2\_COEFF}    & 1.80 & $2\,000\,\um^2/\text{s}$ in water; scaled to grid units.\\
\code{DIFF\_GLC\_COEFF}   & 0.90 & Glucose $\sim 670\,\um^2/\text{s}$.\\
\code{DIFF\_CO2\_COEFF}   & 1.50 & CO$_2$ $\sim 1\,900\,\um^2/\text{s}$.\\
\code{DIFF\_PH\_COEFF}    & 0.60 & H$^+$ effective (buffered).\\
\bottomrule
\end{longtable}

\section{Metabolic constants}

\begin{longtable}{p{4.6cm}rrp{6.0cm}}
\toprule
\textbf{Constant} & \textbf{Value} & \textbf{Units} & \textbf{Rationale}\\ \midrule
\code{O2\_CONSUME\_BASE}    & 0.0075 & mM/\biosec & HeLa OCR rate.\\
\code{GLC\_CONSUME\_BASE}   & 0.0040 & mM/\biosec & HeLa glucose uptake.\\
\code{CO2\_PRODUCE\_BASE}   & 0.0060 & mM/\biosec & Aerobic respiration.\\
\code{LACTATE\_PRODUCE}     & 0.0030 & mM/\biosec & Warburg component.\\
\code{BIOMASS\_SYNTH\_K}    & 0.120  & $\biosec^{-1}$ & Net biomass accumulation rate.\\
\code{BIOMASS\_DEGRADE\_K}  & 0.0015 & $\biosec^{-1}$ & Basal turnover.\\
\code{BIOMASS\_DIVIDE\_THRESH} & 1.85 & bm & Quality-check threshold.\\
\code{ROS\_MUTATION\_SCALE} & 0.03 & -- & ROS $\to$ replication error coupling.\\
\code{HYPOXIA\_MODERATE}    & 0.15 & mM & Below this: mild hypoxia.\\
\code{HYPOXIA\_SEVERE}      & 0.08 & mM & Below this: severe.\\
\code{HYPOXIA\_NECROTIC\_O2}& 0.035& mM & Below this: necrosis countdown.\\
\code{HYPOXIA\_NECROTIC\_TIME}& 120.0 & \biosec & Countdown duration.\\
\code{MITO\_HEALTH\_SWITCH} & 0.4 & -- & Warburg switch-on.\\
\code{MITO\_HEALTH\_RESTORE}& 0.7 & -- & Warburg switch-off.\\
\code{WARBURG\_GLY\_BOOST}  & 1.6 & -- & Glycolysis multiplier.\\
\code{WARBURG\_OXP\_PENALTY}& 0.5 & -- & OxPhos multiplier.\\
\code{MECH\_P21\_COUPLING}  & 0.002 & -- & Crowding $\to$ p21.\\
\code{CONTACT\_INHIBIT\_NBRS} & 6.0 & -- & Neighbour-count threshold.\\
\bottomrule
\end{longtable}

\section{Cycle phase targets}

\begin{longtable}{p{4.6cm}rp{5.0cm}}
\toprule
Phase target & Value & Bio-hour meaning\\ \midrule
\code{CYCLE\_G1\_TARGET\_BIOSEC} & 39600 & 11\,\bioh\\
\code{CYCLE\_S\_TARGET\_BIOSEC}  & 28800 &  8\,\bioh\\
\code{CYCLE\_G2\_TARGET\_BIOSEC} & 14400 &  4\,\bioh\\
\code{CYCLE\_M\_TARGET\_BIOSEC}  &  3600 &  1\,\bioh\\
\bottomrule
\end{longtable}

Sum = 24\,\bioh, matching classical HeLa generation time.

\section{Apoptosis namespace}

\begin{longtable}{p{5.0cm}rp{5.5cm}}
\toprule
Constant & Value & Meaning\\ \midrule
\code{BODIES\_PER\_CELL\_MIN} & 5 & Min apoptotic bodies per cell.\\
\code{BODIES\_PER\_CELL\_MAX} & 15 & Max.\\
\code{MOMP\_DURATION\_BIOSEC} & 300.0 & 5\,\biomin (Ye 2017).\\
\code{EXECUTION\_DURATION\_BIOSEC} & 1800.0 & 30\,\biomin (Spencer 2009).\\
\code{FRAGMENTATION\_DURATION\_BIOSEC} & 1800.0 & 30\,\biomin.\\
\code{SECONDARY\_NECROSIS\_PORE\_START\_BIOSEC} & 14400.0 & 4\,\bioh.\\
\code{OSMO\_SWELL\_BIOSEC} & 5400.0 & 1.5\,\bioh.\\
\code{LYSIS\_DUMP\_BIOSEC} & 600.0 & 10\,\biomin.\\
\code{LYSIS\_RUPTURE\_SWELL} & 1.5 & Swelling factor at rupture.\\
\code{MEMBRANE\_MASS\_PER\_BIOMASS} & 0.20 & Lipid fraction of cell dry mass.\\
\code{RECEPTOR\_MASS\_PER\_BIOMASS} & 0.07 & Membrane protein fraction.\\
\code{REL\_AA\_PER\_BIOMASS} & 6.0 & mM AA released per bm lysed.\\
\code{REL\_IONS\_PER\_BIOMASS} & 140 & mM total ions.\\
\code{REL\_CALCIUM\_PER\_BIOMASS} & 0.2 & mM.\\
\code{REL\_PYRUVATE\_PER\_BIOMASS} & 0.15 & mM.\\
\code{REL\_LACTATE\_PER\_BIOMASS} & 1.0 & mM.\\
\code{REL\_GLUCOSE\_PER\_BIOMASS} & 0.05 & mM.\\
\code{REL\_WATER\_PER\_BIOMASS} & 55000 & mM (bulk water).\\
\code{SHRINK\_FRAC\_AT\_COMPLETE} & 0.35 & 35\% shrinkage across execution.\\
\code{BLEB\_AMPLITUDE\_FRAC} & 0.10 & Bleb radius perturbation.\\
\code{DAMP\_NEIGHBOR\_RADIUS\_WU} & 2.0 & Radius for DAMP propagation.\\
\code{DAMP\_NEIGHBOR\_ROS\_BOOST} & 8.0 & ROS bump per DAMP event.\\
\code{DAMP\_NEIGHBOR\_STRESS\_BOOST} & 4.0 & Stress bump.\\
\code{DAMP\_NEIGHBOR\_FASL\_NG\_PER\_ML} & 2.0 & Fas-L bump.\\
\code{UPTAKE\_PARTICLES\_PER\_FRAGMENTATION\_PER\_BIOMASS} & 4 & Efferocytosis rate.\\
\bottomrule
\end{longtable}

\section{Medium composition (DMEM + 10\% FBS)}

\begin{longtable}{p{4.6cm}rrp{5.5cm}}
\toprule
Constant & Value & Units & Rationale\\ \midrule
\code{DMEM\_O2\_MM} & 0.20 & mM & 21\% atm, saturated at 37$^\circ$C.\\
\code{DMEM\_CO2\_MM} & 1.20 & mM & 5\% atm, buffered by bicarb.\\
\code{DMEM\_GLUCOSE\_MM} & 25.0 & mM & High-glucose DMEM.\\
\code{DMEM\_GLUTAMINE\_MM} & 4.0 & mM & Standard L-glutamine.\\
\code{DMEM\_PYRUVATE\_MM} & 1.0 & mM & Sodium pyruvate optional.\\
\code{DMEM\_AA\_POOL\_MM} & 5.0 & mM & Aggregate of 19 amino acids.\\
\code{DMEM\_GROWTH\_FACTOR\_NG\_PER\_ML} & 50.0 & ng/mL & 10\% FBS typical.\\
\code{DMEM\_IONS\_MM} & 150.0 & mM & Total osmolyte ions.\\
\code{DMEM\_CALCIUM\_MM} & 1.8 & mM & Standard DMEM Ca.\\
\code{DMEM\_PH} & 7.40 & pH & Arterial blood pH.\\
\code{DMEM\_WATER\_MM} & 55\,000 & mM & Bulk water.\\
\bottomrule
\end{longtable}

\section{Summary}

In total, \file{Constants.h} defines roughly 150 named values. Each
is either a pure geometric / time / integration anchor, or a
literature-grounded biological value. No rate in the simulator is
chosen arbitrarily; every calibration moves them against reference
data.

\chapter{Table of Every Function Declared in Simulation.h}
\label{ch:sim-funcs}

\begin{longtable}{p{5.2cm}p{2.0cm}p{7.0cm}}
\toprule
\textbf{Function} & \textbf{Scope} & \textbf{Purpose}\\ \midrule
\code{Simulation::update(dt)} & public & Main per-tick orchestrator.\\
\code{Simulation::init(mode)} & public & Reset and seed initial population.\\
\code{Simulation::updatePhysics(dt)} & private & Integrate cell motion.\\
\code{Simulation::updateMetabolism(c, dt)} & private & Per-cell metabolism step.\\
\code{Simulation::updateMitochondria(c, dt)} & private & mitoPotential + mitoHealth.\\
\code{Simulation::updateAdhesion(c, dt)} & private & Focal-adhesion timeline.\\
\code{Simulation::updateCellCycle(c, dt)} & private & CDK ODE + phase advance.\\
\code{Simulation::updateMitosisProgram(c, dt)} & private & Background per-cell mitosis.\\
\code{Simulation::updateFate(c, dt)} & private & Fate-decision scoring + lock.\\
\code{Simulation::updateApoptosis(c, dt)} & private & 11-trigger cascade orchestrator.\\
\code{Simulation::updateDrugPK(dt)} & private & Drug flux + binding + modulators.\\
\code{Simulation::updatePathogens(dt)} & private & Virion + bacterium physics.\\
\code{Simulation::releasePathogens(c)} & private & Apoptotic spill hook.\\
\code{Simulation::releaseCytosol(c, f)} & private & Per-frac biomass $\to$ medium.\\
\code{Simulation::releaseMembrane(c, f)} & private & Per-frac lipid $\to$ medium.\\
\code{Simulation::releaseReceptors(c, f)} & private & Per-frac receptor $\to$ medium.\\
\code{Simulation::releaseAllMass(c)} & private & Full flush at BODIES transition.\\
\code{Simulation::processDivisions()} & private & Commit queued divisions.\\
\code{Simulation::deriveDaughter(...)} & private & Build one daughter from snapshot.\\
\code{Simulation::seedVirion(...)} & public & Inoculate N virions.\\
\code{Simulation::seedBacterium(...)} & public & Inoculate N bacteria.\\
\code{Simulation::cellNearest(cx, cz)} & public const & Find nearest cell for seeding.\\
\code{Simulation::applyDrug(id, conc)} & public & Apply drug uniformly.\\
\code{Simulation::syncCellBioagentInventory()} & private & Grow per-cell drug vectors on add.\\
\code{Simulation::syncCellInstances()} & private & Rebuild render instance buffer.\\
\code{Simulation::updateStats()} & private & Tally cell phases.\\
\code{Simulation::allocateCellUid()} & public & Unique id generator.\\
\code{Simulation::startMitosisProgram(c)} & private & Begin per-cell mitosis timeline.\\
\code{Simulation::syncMitosisCheckpointInputs(c)} & private & Feed replication status.\\
\code{Simulation::shouldRenderCellShellSourceIndex(i)} & private & Visibility filter.\\
\code{Simulation::activeFocusCellIndex()} & public & Current focus index.\\
\bottomrule
\end{longtable}

The above is a complete census of \code{Simulation}-class member
functions. Free helpers (like \code{evalG1S}, \code{evalG2M},
\code{evalSAC}, \code{biomassFromCellSizeFactor}, and the Phase 8A
\code{geneSequence} / \code{pickGeneForTranscription} etc.) live in
the same file outside the class.

\chapter{Table of Every CentralDogma Function}
\label{ch:cdogma-funcs}

\begin{longtable}{p{5.5cm}p{8.0cm}}
\toprule
\textbf{Function (CentralDogmaState)} & \textbf{Purpose}\\ \midrule
\code{init()} & Reset all sub-state; build HBB codingMRNA; find start/stop codon.\\
\code{update(dt, phase)} & Top-level per-tick driver.\\
\code{updateReplication(dt)} & S-phase fork dynamics.\\
\code{stepReplicationMismatchRepair(dt)} & Check proofreading + MMR.\\
\code{updateReplicationCoverage()} & Count bases covered.\\
\code{stepReplicationForks(dt)} & Per-fork physics.\\
\code{stepReplicationOrigins(dt)} & Origin firing + rescue.\\
\code{startTranscription(polymeraseIndex)} & Begin transcription; Phase 8A picks gene.\\
\code{startTranslation(ribosome, transcript)} & Begin translation.\\
\code{finishTranslation(ribosome)} & Clean-up after stop.\\
\code{releaseTRNA(idx)} & Free a tRNA slot.\\
\code{primeCurrentCodon(tr)} & Look up next codon + amino acid.\\
\code{releaseProtein(tr)} & Commit nascent peptide; Phase 8A tags geneId.\\
\code{bindSpliceosomeIfNeeded(tx)} & Start splicing if idle.\\
\code{clearTranscriptionState(ts)} & Reset a slot.\\
\code{clearSplicingState(sp)} & Reset splicing slot.\\
\code{clearTranslationState(tr)} & Reset translation slot.\\
\code{clearTranscriptState(tx)} & Reset transcript record.\\
\code{clearTRNAState(t)} & Reset tRNA slot.\\
\code{clearProteinState(p)} & Reset protein slot.\\
\code{findFreeTranscriptSlot()} & First inactive slot index.\\
\code{findFreeSpliceosome()} & First free spliceosome.\\
\code{findFreeTRNA()} & First free tRNA slot.\\
\code{findFreeProteinSlot()} & First free protein slot.\\
\code{getTranscript(idx)} & Safe lookup.\\
\code{findTranscriptForTranslation()} & First mature exported transcript.\\
\code{totalCodingCodons()} & HBB-specific length for progress UI.\\
\code{isBaseReplicated(bp)} & Fork-coverage query.\\
\code{replicatedBaseCount()} & Aggregate.\\
\code{countActiveReplicationForks()} & Diagnostic.\\
\code{countFiredOrigins()} & Diagnostic.\\
\code{findFreeForkSlot()} & Available fork slot.\\
\code{findFreeReplicationErrorSlot()} & Available error slot.\\
\code{replicationReadyForM()} & G2/M gate input.\\
\code{geneSequence(idx, outLen)} & Phase 8A routed read.\\
\code{geneLength(idx)} & Phase 8A routed length.\\
\code{pickGeneForTranscription()} & Phase 8A promoter-weighted pick.\\
\code{registerViralGene(gene)} & Phase 8A append to registry.\\
\code{resetReplicationProgram()} & Post-division reset.\\
\code{clearMismatchState(err)} & Reset one mismatch record.\\
\code{countActiveTranscriptions()} & Diagnostic.\\
\code{countActiveTranslations()} & Diagnostic.\\
\code{countActiveProteins()} & Diagnostic.\\
\code{countActiveTRNAs()} & Diagnostic.\\
\bottomrule
\end{longtable}

Plus non-member helpers in the same file:
\code{transcribe}, \code{rnaComplement}, \code{anticodonForCodon},
\code{translateCodon}, \code{aminoAcidByLetter},
\code{isExon}, \code{buildHBBCodingMRNA}, \code{findStartCodon},
\code{findStopCodon}, \code{homopolymerRunLengthAt},
\code{localGCFraction}, \code{bpToMRNAHeuristic}.

\chapter{Every ParticleType Documented}
\label{ch:particle-types}

\begin{longtable}{p{4.5cm}p{1.5cm}p{8.0cm}}
\toprule
\textbf{ParticleType} & \textbf{Radius (\um)} & \textbf{Biological referent}\\ \midrule
\code{PT\_GLUCOSE}        & 0.0005 & Single glucose molecule; rendered at MOL\_VIS\_BOOST\_SMALL scale.\\
\code{PT\_ATP}            & 0.0006 & ATP; very small, heavily visual-boosted.\\
\code{PT\_MITO}           & 1.0    & One mitochondrion (rendered via GLB mesh).\\
\code{PT\_ER\_ROUGH}      & 1.5    & Rough ER segment (with bound ribosomes).\\
\code{PT\_ER\_SMOOTH}     & 1.0    & Smooth ER segment.\\
\code{PT\_GOLGI}          & 1.3    & Golgi apparatus stack.\\
\code{PT\_NUCLEUS}        & 4.0    & Nucleus mesh.\\
\code{PT\_LYSOSOME}       & 0.25   & One lysosome.\\
\code{PT\_DNA\_A}         & 0.002  & DNA backbone node side A.\\
\code{PT\_DNA\_B}         & 0.002  & DNA backbone node side B.\\
\code{PT\_HISTONE}        & 0.010  & Histone octamer.\\
\code{PT\_RIBOSOME\_SMALL}& 0.010  & 40S subunit.\\
\code{PT\_RIBOSOME\_LARGE}& 0.012  & 60S subunit.\\
\code{PT\_TRNA}           & 0.004  & Charged tRNA.\\
\code{PT\_POLYPEPTIDE}    & 0.002  & Nascent peptide.\\
\code{PT\_SPLICEOSOME}    & 0.015  & Spliceosome.\\
\code{PT\_NUCLEAR\_PORE}  & 0.050  & NPC.\\
\code{PT\_RECEPTOR\_GPCR} & 0.005  & 7-pass TM receptor.\\
\code{PT\_RECEPTOR\_RTK}  & 0.005  & Single-pass RTK.\\
\code{PT\_RECEPTOR\_CYTOKINE} & 0.005 & JAK-STAT-linked.\\
\code{PT\_RECEPTOR\_DEATH} & 0.005 & Fas, TNF-R.\\
\code{PT\_RECEPTOR\_TLR}   & 0.006 & Toll-like.\\
\code{PT\_RECEPTOR\_NOTCH} & 0.004 & Notch (proteolysis-released).\\
\code{PT\_RECEPTOR\_FRIZZLED} & 0.005 & Wnt receptor.\\
\code{PT\_RECEPTOR\_PATCHED} & 0.005 & Hedgehog receptor.\\
\code{PT\_RECEPTOR\_STK}   & 0.005 & Ser/Thr kinase receptor (TGF-$\beta$).\\
\code{PT\_ION\_CHANNEL}    & 0.004 & Voltage-gated Na/K/Ca/Cl.\\
\code{PT\_PUMP}            & 0.005 & Na$^+$/K$^+$ ATPase etc.\\
\code{PT\_CHAPERONE}       & 0.007 & HSP70, BiP.\\
\code{PT\_VESICLE\_COPII}  & 0.030 & ER $\to$ Golgi vesicle.\\
\code{PT\_VESICLE\_SECRETORY} & 0.040 & Outgoing vesicle.\\
\code{PT\_MOLECULE}        & 0.002 & Generic small molecule.\\
\bottomrule
\end{longtable}

\chapter{MolTag Values and Their Meaning}
\label{ch:moltag}

\code{MolTag} is a lightweight tag on each \code{IntraParticle} that
lets attraction fields filter. Unlike \code{ParticleType} (which is
about render geometry), \code{MolTag} is about chemical identity.

\begin{longtable}{p{3.5cm}p{10.5cm}}
\toprule
\textbf{Tag} & \textbf{Meaning}\\ \midrule
\code{TAG\_NONE}        & Untagged / generic.\\
\code{TAG\_ATP}         & ATP molecule.\\
\code{TAG\_GLUCOSE}     & Glucose molecule.\\
\code{TAG\_NADH}        & NADH cofactor.\\
\code{TAG\_AMINO\_ACID} & Amino acid (for translation).\\
\code{TAG\_CAMP}        & cAMP second messenger.\\
\code{TAG\_WATER}       & Water.\\
\code{TAG\_CALCIUM}     & Ca$^{2+}$ ion.\\
\code{TAG\_LIPID}       & Membrane lipid.\\
\code{TAG\_LIGAND\_GF}       & Growth factor (generic).\\
\code{TAG\_LIGAND\_HORMONE}  & GPCR ligand.\\
\code{TAG\_LIGAND\_CYTOKINE} & Immune cytokine.\\
\code{TAG\_CAMP\_SIGNAL}     & cAMP burst visualisation.\\
\code{TAG\_BIOAGENT} (Phase 7) & Drug / virion / bacterium particle.\\
\bottomrule
\end{longtable}

\chapter{MolJob State Machine}
\label{ch:moljob}

\begin{longtable}{p{3.5cm}p{10.5cm}}
\toprule
\textbf{Job} & \textbf{Effect on particle}\\ \midrule
\code{JOB\_IDLE}         & No target; pure Brownian + attraction-field response.\\
\code{JOB\_GOTO\_TARGET} & Extra pull toward \code{attachTarget}; arrival triggers next state.\\
\code{JOB\_ATTACH}       & Held at target; reaction dwell counting down.\\
\code{JOB\_CONSUMED}     & Marked for garbage collection at next tick.\\
\bottomrule
\end{longtable}

State transitions are driven by:
\begin{itemize}
  \item Manual \code{assignAttach(particle, target)} at the start of
        a reaction (e.g.\ glucose pulled to hexokinase).
  \item Distance check inside the Langevin step: on arrival within
        \code{ATTACH\_RADIUS}, switch to \code{JOB\_ATTACH}.
  \item Dwell expiry: switch to \code{JOB\_CONSUMED}.
  \item Garbage collection at end of tick: erase consumed
        particles; optionally spawn a product particle at the same
        position.
\end{itemize}

\chapter{The ApoptoticBody Lifecycle}
\label{ch:apobody-deep}

\file{src/main.mm} defines \code{ApoptoticBody} around line 3258.
It carries its own mass ledger (cytosol, membrane, receptor), a
phase (\code{ABODY\_DRIFT, ABODY\_SWELL, ABODY\_BURST}), age, spin
phase, and kind (0 = generic, 1 = nuclear fragment, 2 = mito
fragment).

\section{Birth}

When a cell transitions into FRAGMENTATION, \code{spawnApoptoticBodies(c)}
partitions its final biomass / membrane / receptor pool across
5--15 bodies. The partition is proportional but includes a
stochastic noise factor. Specialised bodies get lower biomass but
carry distinctive fragments: nuclear bodies hold chromatin fragments
(coloured deep red); mito bodies hold cytochrome c residue (coloured
green-tinged).

\section{Drift}

Each body is an independent Brownian particle in the dish fluid.
Its position advects with the curl-noise fluid current. Drag is
proportional to body radius; small bodies drift faster.

\section{Swell}

After \code{SECONDARY\_NECROSIS\_PORE\_START\_BIOSEC} = 4\,\bioh
the body's outer membrane develops nanopores (Chen 2016 mechanism).
Water floods in osmotically; the body swells. The engine tracks
\code{osmoticSwell} which grows linearly with time in this phase.

\section{Burst}

At \code{osmoticSwell >= LYSIS\_RUPTURE\_SWELL = 1.5} the body
ruptures. Contents spill into the medium at the body's final
position via \code{MediumField::exchange}. The body itself fades
to zero alpha over \code{LYSIS\_DUMP\_BIOSEC} = 10\,\biomin.

\section{Efferocytosis}

Live neighbour cells within 4\,\um of an ApoBody consume it via
\code{applyEfferocytosis}. The body's ledger moves into the
consuming cell's \code{lysosomalLoad\_{cyto,mem,rec}} pools;
\code{updateLysosomalDigestion} then converts $\sim$85\% back into
the cell's own biomass, with the remaining 15\% effectively
metabolised for energy (ATP gain). The efferocytosis rate is
calibrated from Ravichandran 2015.

\chapter{MediumChemical Swarm --- Visible Metabolism}
\label{ch:mediumchem-deep}

\code{MediumChemical} in \file{src/main.mm} is the rendered particle
swarm that represents the medium's floating molecules. One to a few
hundred particles per species float in the dish fluid; each has its
own Brownian drift and job state machine
(\code{MC\_FREE, MC\_ATTRACTED, MC\_BINDING, MC\_TRANSPORT}).

\section{Init}

\code{initMediumChemicals()} spawns:
\begin{itemize}
  \item $\sim$80 glucose particles.
  \item $\sim$20 ATP / ADP / NAD particles (visible cofactors).
  \item $\sim$40 water bubbles.
  \item $\sim$10 Ca$^{2+}$ ions.
  \item $\sim$10 O$_2$ / CO$_2$ gas bubbles.
\end{itemize}

\section{Per-tick update}

\code{updateMediumChemicals(dt)} runs the state machine. Free
particles drift. Occasionally a free particle within a cell's
uptake radius flips to ATTRACTED (pulled toward cell surface).
After arrival it enters BINDING (short pulse), then TRANSPORT
(fades into interior), finally DESPAWN.

\section{Rendering}

Each particle is drawn as a ball-and-stick molecule using the
\code{MoleculeCache} entry for that species. The shader fades the
alpha in TRANSPORT so the molecule appears to dissolve into the
cell. Together with the drug particle system (same pattern, just
with \code{DRUG\_SPECIES\_BASE = 1000} offset) this produces the
``visible chemistry'' view that differentiates this simulator from
abstract spreadsheet models.

\chapter{IntracellularPhysics --- Complete Listing}
\label{ch:intraphysics-full}

\section{IntraParticle record}

\begin{lstlisting}
struct IntraParticle {
    simd_float3 position, velocity;
    simd_float3 home;
    simd_float3 confineCenter;
    float radius;
    ParticleType type;
    MolTag tag;
    MolJob job = JOB_IDLE;
    int stateIndex = 0;
    bool active = true;
    int attachTargetIdx = -1;
    float attachTimer = 0.0f;
    uint32_t uid = 0;
};
\end{lstlisting}

\section{AttractionField record}

\begin{lstlisting}
struct AttractionField {
    simd_float3 center;
    float radius;
    float strength;
    MolTag attractedTag = TAG_NONE;
    ParticleType attractedType = PT_MOLECULE;
};
\end{lstlisting}

\section{PhysicsParams}

\begin{lstlisting}
struct PhysicsParams {
    float cellRadius;
    float dragCoef = 4.0f;      // s^-1
    float brownSigma = 0.01f;   // wu per bio-s half
    float hertzStiffness = 800.0f;
    float thermalEnergy = 1.0f; // scaled kT
};
\end{lstlisting}

\section{Key step() structure}

\begin{lstlisting}
void step(float dt) {
    buildSpatialHash();
    for (auto& p : particles) {
        if (!p.active) continue;
        // 1. Deterministic force: Hertz + attraction fields + job.
        simd_float3 F = computeForce(p);
        // 2. Integrate BBK Verlet.
        p.velocity = (1.0f - params.dragCoef * dt * 0.5f) * p.velocity + F * (0.5f * dt);
        p.position += p.velocity * dt;
        p.velocity *= (1.0f - params.dragCoef * dt * 0.5f);
        // 3. Langevin kick.
        simd_float3 noise = {ipRandGauss(), ipRandGauss(), ipRandGauss()};
        float sigma = sqrtf(2.0f * params.dragCoef * params.thermalEnergy * dt);
        p.velocity += noise * sigma;
        // 4. Confinement: project back inside membrane if over.
        intraProjectInsideConfinement(p);
        // 5. Job state transitions.
        updateJob(p, dt);
    }
}
\end{lstlisting}

\chapter{The Telemetry Log}
\label{ch:telemetry}

\file{src/simulation/TelemetryLog.h} (277 lines) is the per-tick
CSV logger used by \code{cellsim\_headless}. It is a small,
self-contained class that writes two files:
\file{logs/telemetry\_population\_<session>.csv} and
\file{logs/telemetry\_division\_<session>.csv}.

\section{Sampling strategy}

Every \code{populationSampleBiosec} (default 60\,\biosec), the
logger writes one row of population summary: cell count, phase
counts, average CDK readings, ATP, biomass, damage, etc.
Division events are logged immediately on occurrence (not on a
schedule), capturing parent/daughter UIDs, positions, and split
stats.

\section{Schema}

\paragraph{Population CSV} (one row per sample):
\code{sample\_idx, bio\_sec, cell\_count, phase\_G1, phase\_S,
phase\_G2, phase\_M, proliferating, quiescent, apoptotic, necrotic,
glycolytic, avg\_ATP, avg\_biomass, avg\_damage,
avg\_replProgress, avg\_CycD, avg\_CycE, avg\_CycA, avg\_CycB,
avg\_Rb, avg\_p21}.

\paragraph{Division CSV} (one row per event):
\code{bio\_sec, parent\_uid, parent\_clone, parent\_gen,
daughter\_A\_uid, daughter\_B\_uid, split\_biomass\_frac\_A,
split\_damage\_frac\_A, telomere\_A\_bp, telomere\_B\_bp,
pos\_A\_x, pos\_A\_y, pos\_A\_z, pos\_B\_x, pos\_B\_y, pos\_B\_z}.

\section{Usage}

\begin{lstlisting}
TelemetryLog tlog;
tlog.populationSampleBiosec = 60.0f;
tlog.open("headless_20260420_120000");
// ... run loop ...
tlog.sample(sim);
tlog.close();
\end{lstlisting}

\chapter{Specific Simulation Scenarios}
\label{ch:scenarios}

This chapter walks through seven representative scenarios, each
tracing what the user would see and which code paths drive it.

\section{Scenario 1: HeLa in DMEM (default template)}

Setup: 5 cells, DMEM + 10\% FBS, 5\% CO$_2$, pH 7.40. Run 72\,\bioh.

Expected outcome: doubling $\sim$every 24\,\bioh; cell count grows
from 5 to $\sim$40--80 depending on crowding. Almost no deaths.
Cycle phases settle to approximately 45\% G1, 35\% S, 15\% G2, 5\% M
at steady state (matching BrdU pulse-chase distribution).

Relevant code: \code{updateCellCycle}, \code{updateMetabolism},
\code{processDivisions}. Watch: CDK readouts should oscillate with
period $\sim$24\,\bioh; \code{replicationProgress} saw-tooths from 0
to 1; biomass grows then halves.

\section{Scenario 2: Hypoxia 3\% O$_2$}

Setup: 10 cells, O$_2$ = 0.04\,mM (3\% atm).

Expected outcome: cells enter Warburg mode (glycolytic = true)
within 30\,\biomin. Lactate builds up in medium; pH drops from 7.40
toward 7.00. Cycle progression slows (HIF-1$\alpha$-driven CDK
inhibition; models Pouys{\'e}gur 2006). Some cells enter necrosis
after 2\,\bioh of O$_2$ $<$ 0.035\,mM.

Watch: \code{glycolytic} flag flips true; \code{mitoHealth} drops
below 0.4; \code{lactate} mM in medium rises.

\section{Scenario 3: Starvation}

Setup: 5 cells, glucose 1\,mM, glutamine 0.5\,mM, AA pool 1\,mM.

Expected outcome: ATP falls below 50\%. G1/S gate closes
(\code{evalG1S} fails with ATP\_LOW). Cells enter quiescence; some
eventually commit to apoptosis via \code{survival\_loss} trigger
when GF-signalling is low and ATP crash-ed.

Watch: \code{fate} = SIM\_FATE\_QUIESCENT dominates; then
\code{fate} = SIM\_FATE\_APOPTOTIC rises after $\sim$10\,\bioh.

\section{Scenario 4: Dense Colony (400 cells)}

Setup: 400 cells in the default medium.

Expected outcome: fast initial division phase until
\code{MONOLAYER\_CAPACITY = 600} is approached; then crowding
pressure rises, \code{p21} rises (via \code{MECH\_P21\_COUPLING}),
G1/S gate closes, cells enter quiescence. The famous
contact-inhibition phenotype.

Watch: \code{localPressure} field grows; \code{p21} climbs across
the population.

\section{Scenario 5: CTC HeLa reference (seq02)}

Setup: 125 cells, CTC medium (10\% CO$_2$ rather than 5\%; more
glutamine; more AA).

Expected outcome: $\sim$252 divisions over 46\,\bioh; final count
$\sim$377 (ref: 377). Mean $|\text{rel err}|$ = 5.1\%.

This is the primary calibration benchmark, run by
\code{cellsim\_headless 46 60 125}.

\section{Scenario 6: Cisplatin 5\,\um M}

Setup: 5 cells, default medium; apply 5\,\um M cisplatin via Drug
Lab.

Expected outcome: cisplatin diffuses into cells (Lipinski
permeability low, so slow); binds DNA intercalation target;
damageLevel rises; p53 activates; Puma + Bax activate; MOMP; death
around 24--36\,\bioh in most cells. The characteristic cisplatin
death curve.

Watch: \code{damageLevel} rising; cells transitioning to PRIMED
then MOMP in sequence.

\section{Scenario 7: Flu infection (Phase 7)}

Setup: 1 cell (test light mode); inoculate 50 flu\_h1n1 virions.

Expected outcome: virions drift in; $\sim$45--50 bind receptors;
all 50 enter by 10\,\biomin; viral protein count rises; budding
starts; $\sim$300 new free virions released by 1\,\bioh;
eventually host cell apoptoses (cytotoxic load); releasePathogens
spills virions and bacteria to medium; since there's only 1 cell,
the chain stops there.

If run in colony mode with 5 cells, the infection would propagate
to neighbours --- the primary demonstration of the physics-driven
infection chain.

\chapter{Parameter Sweep Strategies}
\label{ch:sweeps}

\section{Why sweeps matter}

No single constant is ``right''; every constant in
\file{Constants.h} comes with a literature range and a choice
within it. A sensitivity sweep quantifies how much each constant
matters.

\section{Standard sweep: SLOW vs MEDIUM dt scales}

\begin{verbatim}
# scripts/calibrate_rate.sh
for slow in 0.050 0.052 0.055 0.058 0.061; do
  for med in 0.180 0.190 0.201 0.210 0.220; do
    for seed in 1 2 3 4 5; do
      CSVS="$CSVS cal_${slow}_${med}_${seed}.csv"
      SLOW_DT_SCALE=$slow MEDIUM_DT_SCALE=$med \
        ./build/cellsim_headless 46 60 125 > log_${slow}_${med}_${seed}.txt
      cp logs/headless_doubling.csv cal_${slow}_${med}_${seed}.csv
    done
  done
done
python scripts/analyze_sweep.py $CSVS
\end{verbatim}

The current optimum: $\code{SLOW}=0.055, \code{MED}=0.20096$.
Sensitivity: a 10\% change in either scales the final cell count by
$\sim$5\%; a 20\% change pushes into the 10-15\% error range.

\section{Cycle-ODE sensitivity}

Each rate in Ch.~\ref{ch:cycle-deep} can be varied $\pm 20\%$. The
most sensitive are $k_{D,\text{synth}}$ (G1 duration) and
$k_{B,K}$ (M-exit timing). The least sensitive are
$k_{E,\text{deg}}$ and $k_{P,\text{synth}}$. Full sensitivity
matrix is an upcoming Phase L task.

\section{Apoptosis sweep}

Cisplatin dose--response can be swept by running
\code{cellsim\_headless} with different drug concentrations and
measuring time-to-death distribution. A matching curve to Sia 2008
would validate the apoptosis engine calibration.

\chapter{Versus Other Cell Simulators}
\label{ch:versus}

\textsc{CellSim} occupies a specific niche. This chapter compares
it to other tools so the reader can place it.

\section{Versus PhysiCell}

PhysiCell (Ghaffarizadeh 2018) is agent-based, BSD-licensed, runs
in C++ with XML configuration. Strength: many cells (10$^5$+) in a
single simulation; mature community. Weakness: little
intracellular detail per cell; no visible chemistry. CellSim trades
population scale ($\sim$2000 cells) for per-cell detail (visible
ribosomes, transcription, drug flux, infection chain).

\section{Versus COPASI / VCell}

COPASI (ODE-focused) and VCell (spatial PDE + ODE) are excellent
for published SBML pathway models. They do not do 3D rendering,
agent-based cell-cell interactions, or user-interactive drug
application. CellSim offers SBML export in Phase 7U; the idea is
that a user can hand off to VCell for analysis after running a
scenario in CellSim.

\section{Versus E-Cell / WholeCellKB}

E-Cell (Karr 2012) pioneered whole-cell modelling with every gene
and reaction for \emph{M.\ genitalium}. This is the gold standard
of per-cell completeness. CellSim is far less complete but more
interactive --- a whole-cell E-Cell simulation of HeLa would be a
several-engineer-year effort; CellSim is accessible to a single
developer.

\section{Versus Folding@Home / GROMACS}

Molecular dynamics tools model protein structure at atomic
resolution. CellSim does not. Chemistry Tier 4 of the roadmap
(Ch.~\ref{ch:phase8-roadmap}) optionally delegates to OpenMM for
protein--ligand binding MD; this keeps CellSim's interactive speed
while gaining gold-standard $K_d$ estimates offline.

\section{Verdict}

CellSim is best viewed as an \emph{interactive biology} tool ---
closer to an educational / hypothesis-generation platform than a
production research simulator. Its 3D rendering, chemistry-first
drug system, and infection chain give it a distinctive role.

\chapter{Debugging Notes and Known Issues}
\label{ch:known-issues}

\section{G2 stall diagnostic}

The headless run prints a G2 stall diagnostic when cells pile up
in G2:
\begin{verbatim}
[headless] G2-stall diagnostic (first 5 G2 cells):
  cell uid=1 bio=2.30 ATP=100 p21=1.00 CycB=0.93 repl=1.00 ...
\end{verbatim}
This is informational, not a fault. Cells in dense colonies pile
up in G2 due to p21 from contact inhibition; this is the expected
biology.

\section{Inflated HBB sequence}

The 1558\,bp HBB sequence includes the full 852\,bp intron 2 which
was previously truncated to 120\,bp for performance. The correction
(back to full length) adds negligible CPU cost and fixes a subtle
bias in replication-error statistics that was visible in the
mismatch sweep.

\section{The ``cells don't move at 100x'' report}

Addressed in Ch.~\ref{ch:tips}. Actual cells in a petri dish don't
move much. User-visible dynamics are in the pathogens and drugs
drifting, not in cell motion.

\section{Metal duplicate-library warning}

\begin{verbatim}
ld: warning: ignoring duplicate libraries:
  '/Users/henry/vcpkg/installed/arm64-osx/lib/libglfw3.a'
\end{verbatim}
Harmless. vcpkg's GLFW is pulled in once by the CMake config and
once by a target-level link. Both resolve to the same static
archive.

\section{pdflatex not installed on the dev machine}

The \file{docs/} directory contains the LaTeX sources and a
pre-compiled \file{cellsim\_equations.pdf}. If the reader wants to
rebuild the PDF locally, install a LaTeX distribution (BasicTeX,
MacTeX, TeX Live, or MiKTeX) and run \texttt{pdflatex
cellsim\_full\_reference.tex} twice (first pass generates the TOC;
second pass resolves references).

\chapter{The Logs Directory}
\label{ch:logs}

At every run, CellSim writes a session-named set of logs into
\file{logs/}. Keeping them around is small overhead; they are the
post-mortem trail.

\begin{longtable}{p{5.5cm}p{8.0cm}}
\toprule
\textbf{Log file} & \textbf{Contents}\\ \midrule
\file{mitosis\_<session>.log}    & Every mitosis event: phase transitions, SAC events, furrow progression.\\
\file{diag\_<session>.log}       & Per-120-frame diagnostic: cell position, visible DNA, intracellular physics stats, mitosis flags.\\
\file{centrosome\_<session>.log} & Per-mitosis centrosome tracking: duplication, separation, pole orientation.\\
\file{headless\_doubling.csv}    & Time-series population + CDK averages.\\
\file{compare\_vs\_reference.csv}& Sim vs reference per-row error.\\
\file{telemetry\_population\_<session>.csv} & Full population telemetry.\\
\file{telemetry\_division\_<session>.csv}   & Every division event.\\
\bottomrule
\end{longtable}

\chapter{On The Name ``CellSim''}
\label{ch:name}

The project name is a deliberate understatement. Many things named
\emph{FooSim} are toys; this is intentional. The engine is a
working biology simulator that happens to look simple from the
outside --- one double-click and you're watching cells divide, die,
get infected, get rescued, in 3D. The complexity is inside, in the
twenty biochem modules, the calibrated constants, the
mass-conservation invariant, the emergent drug phenotypes, the
physics-driven infection chain. A new user finds it approachable;
an advanced user finds it rigorous.

The alternative name considered was \emph{Cytos} (Greek for
``cell''). \emph{CellSim} won because it's googlable and has no
trademark conflicts.

\chapter{Future Work}
\label{ch:future}

\section{Immediate (next session)}

Phase 8E: wire the Phase 8A gene registry into actual virus
replication. When a virion uncoats, register its genome as a
GeneTemplate; let host Pol II transcribe it; let host ribosomes
translate it; gate assembly on \code{viralProteinCount}. This is
the first test of whether the chemistry discipline extends from
drugs to pathogens.

\section{Short-term (Part 8 continuation)}

Phases 8B--8L: endosome mechanics, stoichiometric assembly,
Caspar-Klug capsid rendering, T3SS effector injection, real
bacterium wall meshes, conjugation, immune response.
Approximately 16 sessions.

\section{Mid-term (Part 7 continuation)}

Multi-cell-type populations, full GRN with $\sim$200 edges, FBA
metabolism, 3D O$_2$ diffusion, quorum sensing, biofilm,
conjugation, mutation / clonal evolution, immune cells as SimCell
subtypes, ECM fibres, cell-cell adhesion edges. Approximately 30
sessions.

\section{Long-term (Parts 2--6)}

Chemistry tiers: OpenFF, xTB, GNINA, OpenMM, AMOEBA, MACE-OFF,
MARTINI. Ray-traced rendering. VR. Time-scrubber. Multi-cell
focus split. Approximately 60 sessions.

\section{Optional-but-intriguing}

\begin{itemize}
  \item A ``calibrate mode'' that automatically finds the
        constants minimising error against any uploaded reference
        CSV.
  \item An automated literature-mining pipeline that hashes every
        rate in \file{Constants.h} to its literature source and
        produces a citation report on demand.
  \item A patient-specific mode where users upload RNA-seq and
        the simulator re-tunes receptor / expression profiles.
\end{itemize}

The last item is the platonic ideal of personalised biology
simulation: you bring your own cells, the simulator bends to fit.
It is years of work away but worth stating aloud.

% ═════════════════════════════════════════════════════════════════════
\part{Exhaustive Reference Tables}
% ═════════════════════════════════════════════════════════════════════

\chapter{Every Enum in the Codebase}
\label{ch:enums}

\section{Simulation core}

\begin{longtable}{p{4.2cm}p{10.0cm}}
\toprule
\code{SimMode} & MODE\_SINGLE\_CELL, MODE\_COLONY.\\
\code{SimFate} & SIM\_FATE\_UNDETERMINED, SIM\_FATE\_PROLIF, SIM\_FATE\_QUIESCENT, SIM\_FATE\_APOPTOTIC.\\
\code{MitosisPhase} & MITO\_NONE, MITO\_PROPHASE, MITO\_PROMETAPHASE, MITO\_METAPHASE, MITO\_ANAPHASE, MITO\_TELOPHASE, MITO\_CYTOKINESIS, MITO\_COMPLETE.\\
\code{CDogmaPhase} & CD\_IDLE, CD\_REPLICATION, CD\_TRANSCRIPTION, CD\_SPLICING, CD\_EXPORT, CD\_TRANSLATION, CD\_FOLDING, CD\_TRANSPORT.\\
\code{MolJob} & JOB\_IDLE, JOB\_GOTO\_TARGET, JOB\_ATTACH, JOB\_CONSUMED.\\
\bottomrule
\end{longtable}

\section{Apoptosis}

\begin{longtable}{p{4.2cm}p{10.0cm}}
\toprule
\code{Apoptosis::Phase} & ALIVE, PRIMED, MOMP, EXECUTION, FRAGMENTATION, BODIES, CLEARED.\\
\bottomrule
\end{longtable}

\section{Pathogens}

\begin{longtable}{p{4.2cm}p{10.0cm}}
\toprule
\code{VirionShape}     & ICOSAHEDRAL, HELICAL, COMPLEX, PLEOMORPHIC.\\
\code{GenomeKind}      & SSRNA\_POS, SSRNA\_NEG, DSRNA, SSDNA, DSDNA, RETRO.\\
\code{VirionStage}     & FREE, BOUND, ENTERING, UNCOATING, REPLICATING, ASSEMBLING, BUDDING, SPILLED, CLEARED.\\
\code{BacteriumShape}  & COCCUS, ROD, VIBRIO, SPIRILLUM, SPIROCHETE.\\
\code{GramType}        & GRAM\_POSITIVE, GRAM\_NEGATIVE, ACID\_FAST.\\
\code{BacteriumStage}  & FREE, ADHERING, INVADED, REPLICATING, LYSED.\\
\code{PathogenType}    & PATH\_NONE, PATH\_GRAM\_NEG\_BACTERIA, PATH\_GRAM\_POS\_BACTERIA, PATH\_DNA\_VIRUS, PATH\_RNA\_VIRUS, PATH\_RETROVIRUS, PATH\_FUNGUS, PATH\_PARASITE.\\
\code{InfectionStage}  & INF\_NONE, INF\_ATTACHMENT, INF\_ENTRY, INF\_INTRACELLULAR, INF\_ASSEMBLY, INF\_EGRESS, INF\_CLEARED.\\
\code{PAMPtype}        & PAMP\_LPS, PAMP\_PEPTIDOGLYCAN, PAMP\_FLAGELLIN, PAMP\_dsRNA, PAMP\_ssRNA, PAMP\_CpG\_DNA, PAMP\_MANNAN.\\
\bottomrule
\end{longtable}

\section{Chemistry}

\begin{longtable}{p{4.6cm}p{9.6cm}}
\toprule
\code{ChemicalEntityKind} & UNKNOWN, METABOLITE, DRUG, VIRUS, BACTERIUM, ORGANELLE, ENZYME, RECEPTOR, MEMBRANE\_LIPID, NUCLEIC\_ACID, COFACTOR, PEPTIDE, ANTIBODY, ION.\\
\code{PharmacophoreFeature::Kind} & HBOND\_DONOR, HBOND\_ACCEPTOR, HYDROPHOBIC, AROMATIC, POS\_CHARGE, NEG\_CHARGE, METAL\_CHELATOR.\\
\code{EntryPathway}       & NONE, PASSIVE\_DIFFUSION, RECEPTOR\_ENDOCYTOSIS, MEMBRANE\_FUSION, PHAGOCYTOSIS, INVASION.\\
\code{GenePolymerase}     & HOST\_POLII, VIRAL\_RdRp, VIRAL\_RT, VIRAL\_POL.\\
\bottomrule
\end{longtable}

\section{Vesicular traffic}

\begin{longtable}{p{4.2cm}p{10.0cm}}
\toprule
\code{Compartment} & COMP\_ER, COMP\_CIS\_GOLGI, COMP\_MEDIAL\_GOLGI, COMP\_TRANS\_GOLGI, COMP\_TGN, COMP\_EARLY\_ENDOSOME, COMP\_LATE\_ENDOSOME, COMP\_LYSOSOME, COMP\_PLASMA\_MEMBRANE, COMP\_CYTOSOL, COMP\_NUCLEUS, COMP\_COUNT.\\
\code{CoatType}    & COAT\_NONE, COAT\_COPII, COAT\_COPI, COAT\_CLATHRIN.\\
\code{CargoType}   & CARGO\_SECRETORY, CARGO\_REGULATED\_SECRETORY, CARGO\_MEMBRANE\_PROTEIN, CARGO\_LYSOSOMAL\_ENZYME, CARGO\_RECEPTOR\_RECYCLING, CARGO\_RECEPTOR\_DEGRADATION, CARGO\_AUTOPHAGY\_CARGO, CARGO\_ER\_RETRIEVAL.\\
\bottomrule
\end{longtable}

\section{Medium}

\begin{longtable}{p{4.2cm}p{10.0cm}}
\toprule
\code{MediumSpecies} & MS\_O2, MS\_CO2, MS\_GLUCOSE, MS\_GLUTAMINE, MS\_PYRUVATE, MS\_LACTATE, MS\_AA\_POOL, MS\_GROWTH\_F, MS\_IONS, MS\_CALCIUM, MS\_HPLUS, MS\_WATER, MS\_COUNT.\\
\code{MediumChemState} & MC\_FREE, MC\_ATTRACTED, MC\_BINDING, MC\_TRANSPORT, MC\_DESPAWN.\\
\bottomrule
\end{longtable}

\section{Apoptotic body}

\begin{longtable}{p{4.2cm}p{10.0cm}}
\toprule
\code{ApoBodyPhase} & ABODY\_DRIFT, ABODY\_SWELL, ABODY\_BURST.\\
\bottomrule
\end{longtable}

\chapter{Every Struct in Simulation Core}
\label{ch:structs}

Below is an exhaustive list, with their purpose and defining file.

\begin{longtable}{p{4.4cm}p{3.0cm}p{7.0cm}}
\toprule
\textbf{Struct} & \textbf{File} & \textbf{Purpose}\\ \midrule
\code{SimSetup}           & main.mm                             & Startup setup overlay configuration.\\
\code{MolRenderData}      & main.mm                             & GPU buffers for atom + bond pipelines.\\
\code{MediumChemical}     & main.mm                             & Rendered molecule particle in medium.\\
\code{MediumChemSpec}     & main.mm                             & Rendering spec for a medium species.\\
\code{GPUAtomInstance}    & gpu/ShaderTypes.h                   & Per-instance atom for Metal pipeline.\\
\code{GPUBondInstance}    & gpu/ShaderTypes.h                   & Per-instance bond cylinder.\\
\code{CellInstance}       & gpu/ShaderTypes.h                   & Per-instance cell shell.\\
\code{Uniforms}           & gpu/ShaderTypes.h                   & Camera + light + time for pipeline.\\
\code{FluidUniforms}      & gpu/ShaderTypes.h                   & Fluid shader uniforms.\\
\code{ApoptoticBody}      & main.mm                             & Floating apoptotic body with ledger.\\
\code{DampEvent}          & main.mm                             & Damage-associated molecular pattern event.\\
\code{Centrosome}         & main.mm                             & Per-cell centriole pair + pole vectors.\\
\code{MitosisState}       & CellCycleProgram.h                  & Per-cell mitosis sub-state.\\
\code{CellCycleProgram}   & CellCycleProgram.h                  & Wraps cdogma + mitosis + events.\\
\code{SimCell}            & Simulation.h                        & Per-cell primary record.\\
\code{CDKState}           & Simulation.h                        & Cyclins, Rb, p21, p53, E2F, APC/C.\\
\code{SplitStats}         & Simulation.h                        & Division asymmetry parameters.\\
\code{NutrientField}      & Simulation.h                        & Legacy nutrient grid scaffold.\\
\code{CheckpointResult}   & Simulation.h                        & Check pass + reason string.\\
\code{Simulation::DivisionEvent} & Simulation.h                & Queued division to commit.\\
\code{Simulation::AppliedDrug} & Simulation.h                  & Active drug with binding affinities.\\
\code{MediumField}        & MediumField.h                       & 2D reaction-diffusion grid.\\
\code{TelemetryLog}       & TelemetryLog.h                      & CSV logger for headless runs.\\
\code{IntracellularPhysics}& IntracellularPhysics.h             & Particle simulator.\\
\code{IntraParticle}      & IntracellularPhysics.h              & One intracellular particle.\\
\code{AttractionField}    & IntracellularPhysics.h              & Tag-filtered attractor.\\
\code{PhysicsParams}      & IntracellularPhysics.h              & Per-cell physics tuning.\\
\code{SpatialHash}        & IntracellularPhysics.h              & 3D uniform grid index.\\
\code{TranscriptionState} & CentralDogma.h                      & Per-Pol II state.\\
\code{SplicingState}      & CentralDogma.h                      & Per-spliceosome state.\\
\code{TranslationState}   & CentralDogma.h                      & Per-ribosome state.\\
\code{TranscriptState}    & CentralDogma.h                      & Per-mRNA record.\\
\code{ChargedTRNAState}   & CentralDogma.h                      & Per-tRNA slot.\\
\code{ProteinProductState}& CentralDogma.h                      & Per-released-protein slot.\\
\code{ReplicationOriginState}&CentralDogma.h                    & Per-origin state.\\
\code{ReplicationForkState}&CentralDogma.h                      & Per-fork state.\\
\code{ReplicationMismatchState}&CentralDogma.h                  & Per-error record.\\
\code{GeneTemplate}       & CentralDogma.h (Phase 8A)           & Registered gene entry.\\
\code{CentralDogmaState}  & CentralDogma.h                      & Whole cdogma wrapper.\\
\code{ApoptosisState}     & biochem/Apoptosis.h                 & Bcl-2/Bax/caspase pool.\\
\code{ApoTriggers}        & biochem/Apoptosis.h                 & 11-input trigger vector.\\
\code{ApoptosisEngine}    & biochem/Apoptosis.h                 & Runs the cascade ODE.\\
\code{ChemicalEntity}     & biochem/ChemicalEntity.h            & Universal structure record.\\
\code{PharmacophoreFeature}&biochem/ChemicalEntity.h            & One pharmacophore point.\\
\code{BindingAffinity}    & biochem/ChemicalEntity.h            & Drug$\times$target kinetics.\\
\code{TargetProfile}      & biochem/BindingMatcher.h            & Target descriptor preference.\\
\code{TargetDefinition}   & biochem/TargetLibrary.h             & Name + profile + modulator.\\
\code{TargetLibrary}      & biochem/TargetLibrary.h             & Catalogue.\\
\code{BioagentRegistry}   & biochem/Bioagent.h                  & Drug catalogue.\\
\code{VirionSpec}         & biochem/Virion.h                    & Virion catalogue entry.\\
\code{Virion}             & biochem/Virion.h                    & Per-instance virion.\\
\code{BacteriumSpec}      & biochem/Bacterium.h                 & Bacterium catalogue entry.\\
\code{Bacterium}          & biochem/Bacterium.h                 & Per-instance bacterium.\\
\code{ReceptorProfile}    & biochem/PathogenKinetics.h          & Receptor descriptor baseline.\\
\code{PathogenRegistry}   & biochem/PathogenRegistry.h          & YAML-loaded pathogen catalogue.\\
\code{PathogenState}      & biochem/Pathogens.h                 & Per-cell infection state.\\
\code{MicrobiotaState}    & biochem/Pathogens.h                 & Commensal mix.\\
\code{PathogenEngine}     & biochem/Pathogens.h                 & Intracellular infection driver.\\
\code{InnateImmuneState}  & biochem/Immunity.h                  & TLRs + sensors per cell.\\
\code{AdaptiveImmuneState}& biochem/Immunity.h                  & MHC + TCR/BCR repertoire.\\
\code{IonState}           & biochem/MembraneTransport.h         & Ion pools + gating.\\
\code{MembraneTransportEngine}&biochem/MembraneTransport.h      & Step function.\\
\code{Vesicle}            & biochem/VesicularTraffic.h          & One trafficked vesicle.\\
\code{CompartmentState}   & biochem/VesicularTraffic.h          & Per-compartment pH + loads.\\
\code{UPRState}           & biochem/VesicularTraffic.h          & UPR arms.\\
\code{AutophagyState}     & biochem/VesicularTraffic.h          & Autophagy pipeline.\\
\code{VesicularTrafficEngine}& biochem/VesicularTraffic.h       & Step function.\\
\code{MetabolismState}    & biochem/Metabolism.h                & Flux pool.\\
\code{MetabolismEngine}   & biochem/Metabolism.h                & Step function.\\
\code{CadherinJunction}   & biochem/CellJunctions.h             & Per-edge cadherin bundle.\\
\code{Desmosome}          & biochem/CellJunctions.h             & Desmosome bundle.\\
\code{TightJunction}      & biochem/CellJunctions.h             & Tight junction bundle.\\
\code{GapJunction}        & biochem/CellJunctions.h             & Gap junction bundle.\\
\code{Integrin}           & biochem/CellJunctions.h             & Integrin record.\\
\code{FocalAdhesion}      & biochem/CellJunctions.h             & FA maturation record.\\
\code{ECMComponent}       & biochem/CellJunctions.h             & ECM fibre.\\
\code{Selectin}           & biochem/CellJunctions.h             & Transient rolling.\\
\code{CellJunctionsEngine}& biochem/CellJunctions.h             & Step function.\\
\code{ActinState}         & biochem/Cytoskeleton.h              & Actin pool + filaments.\\
\code{MicrotubuleState}   & biochem/Cytoskeleton.h              & MT pool + polarity.\\
\code{IFState}            & biochem/Cytoskeleton.h              & Intermediate filaments.\\
\code{Motors}             & biochem/Cytoskeleton.h              & Myosins + kinesins + dyneins.\\
\code{CytoskeletonEngine} & biochem/Cytoskeleton.h              & Step function.\\
\code{SignalingState}     & biochem/Signaling.h                 & All pathway species.\\
\code{SignalingEngine}    & biochem/Signaling.h                 & Step function.\\
\code{DNADamageState}     & biochem/DNARepair.h                 & Lesion ledger.\\
\code{DNARepairEngine}    & biochem/DNARepair.h                 & BER/NER/MMR/HR/NHEJ/TLS.\\
\code{ProteomeState}      & biochem/Proteome.h                  & Chaperones + PTMs.\\
\code{ProteomeEngine}     & biochem/Proteome.h                  & Step function.\\
\code{GenomeState}        & biochem/GenomeModel.h               & Ploidy + chromatin.\\
\code{GeneExpressionState}& biochem/GeneRegulation.h            & TF activities.\\
\code{CancerState}        & biochem/Cancer.h                    & Transformation profile.\\
\code{DevelopmentState}   & biochem/Development.h               & Morphogen gradients.\\
\code{StemCellState}      & biochem/StemCells.h                 & Stemness core.\\
\code{CellBiologyEngine}  & biochem/CellBiologyEngine.h         & Umbrella wrapper.\\
\bottomrule
\end{longtable}

\chapter{ImGui Panel Visual Mockups}
\label{ch:ui-mockups}

This chapter shows the on-screen layout of each ImGui panel at
default resolution (1440$\times$900). Each mockup is drawn with
Tikz at approximate proportion.

\paragraph{Cell Mode --- Setup (modal, top-centre).}
\begin{verbatim}
+-------------------------------------------------+
| Cell Mode - Setup                               |
+-------------------------------------------------+
| Templates:                                      |
|  [HeLa Standard] [CTC HeLa]  [Hypoxia 3% O2]    |
|  [Starvation]    [Rich Med]  [Single Cell]      |
|  [Dense Colony]  [Test Mode (Lite 16GB)]        |
+-------------------------------------------------+
| Configuration:                                  |
|  Initial cells     [------|-----] 125           |
|  Glucose (mM)      [------|-----] 25.0          |
|  Glutamine (mM)    [------|-----] 4.0           |
|  AA pool (mM)      [------|-----] 5.0           |
|  O2 (mM)           [------|-----] 0.20          |
|  CO2 (mM)          [------|-----] 1.20          |
|  pH                [------|-----] 7.40          |
|  Growth factors    [------|-----] 50.0 ng/mL    |
|  Initial timeScale [------|-----] 5.0x          |
|  [x] Light / Test mode (1 cell, 16GB-friendly)  |
+-------------------------------------------------+
|            |>  Start simulation                 |
+-------------------------------------------------+
\end{verbatim}

\paragraph{Drug Lab panel (right, below Export).}
\begin{verbatim}
+------------------------------------+
| Drug Lab (chemistry-driven)        |
+------------------------------------+
| MOA-free: structure -> emergent    |
+------------------------------------+
| Compound:   [cisplatin        v]   |
|   MW=300.1  logP=-2.3  TPSA=96.5   |
|   HBD=3     HBA=8      rings=0     |
|                                    |
| log[C] uM:  [-----+----]           |
| Concentration: 5.000 uM            |
|                                    |
| [       APPLY UNIFORMLY        ]   |
+------------------------------------+
| Active: cisplatin @ 5.00 uM        |
| Top affinities (score > 0.50):     |
|   TGT_DNA_INTERCALATION  Kd=12 uM  |
|   TGT_TOPO2              Kd=140 uM |
|   TGT_HSP70              Kd=2.1 mM |
+------------------------------------+
\end{verbatim}

\paragraph{Test Actions panel (left, light mode only).}
\begin{verbatim}
+----------------------------------------+
| [!] Test Actions (one click)           |
+----------------------------------------+
| Watch each action play out in 3D.      |
+----------------------------------------+
| Viruses                                |
|  [ Inoculate Flu (H1N1) x 50        ]  |
|  [ Inoculate HBV x 30               ]  |
|  [ Inoculate Adenovirus x 40 lytic  ]  |
+----------------------------------------+
| Bacteria                               |
|  [ Release E. coli x 20             ]  |
|  [ Release Listeria x 15            ]  |
|  [ Release S. aureus x 15 (toxin)   ]  |
+----------------------------------------+
| Drugs (chemistry-driven)               |
|  [ Cisplatin 5 uM (DNA crosslinker) ]  |
|  [ Doxorubicin 2 uM (TopII poison)  ]  |
|  [ Paclitaxel 0.1 uM (tubulin)      ]  |
|  [ Staurosporine 1 uM (pan-kinase)  ]  |
+----------------------------------------+
| Cell state                             |
|  [ Force apoptosis on focused cell  ]  |
|  [ Heal cell (clear damage)         ]  |
|  [ Clear pathogens + drugs + bodies ]  |
+----------------------------------------+
\end{verbatim}

\chapter{The Shader Source --- Cell Render}
\label{ch:shader-cell}

An annotated walkthrough of \file{assets/shaders/CellRender.metal}.
The shader is one vertex + one fragment function that draws every
cell shell in one instanced draw call.

\section{Vertex function}

\begin{lstlisting}[language=C,morekeywords={vertex,fragment,constant,device,metal}]
vertex VertexOut cell_vertex(uint vid [[vertex_id]],
                             uint iid [[instance_id]],
                             device const CellInstance* cells [[buffer(1)]],
                             constant Uniforms& uni [[buffer(2)]]) {
    // Extract one sphere vertex from the shared mesh.
    // ... (sphere mesh is pre-generated at startup)
    const CellInstance c = cells[iid];
    // Apply per-instance translation, scale, furrow deformation.
    float3 localPos = sphereVertex[vid].position;
    // Furrow: pinch radius on equatorial disk.
    float lat = asin(localPos.y);
    float furrow = c.furrowDepth * exp(-10.0 * lat * lat);
    float r = (1.0 - 0.35 * furrow);
    localPos.xz *= r;
    // Apoptosis bleb: small radial perturbation at high apoProgress.
    // (Baked by CPU into furrow via the same channel.)
    float3 worldPos = c.position + localPos * c.radius;
    VertexOut o;
    o.clipPos = uni.viewProjection * float4(worldPos, 1.0);
    o.worldPos = worldPos;
    o.normal = localPos;  // sphere is its own normal in local frame
    o.color = c.color;
    o.glow = c.glowIntensity;
    o.phase = c.phase;
    return o;
}
\end{lstlisting}

\section{Fragment function}

\begin{lstlisting}[language=C,morekeywords={vertex,fragment,constant,device}]
fragment float4 cell_fragment(VertexOut in [[stage_in]],
                              constant Uniforms& uni [[buffer(2)]]) {
    float3 N = normalize(in.normal);
    float3 V = normalize(uni.cameraPos - in.worldPos);
    float3 L = normalize(uni.lightDir);
    // Body term: Lambertian with forward scatter.
    float body = max(0.0, dot(N, L));
    // Fresnel rim: pow(1 - N.V, 3).
    float rim = pow(1.0 - max(0.0, dot(N, V)), 3.0);
    float3 col = in.color.rgb * body
                 + in.color.rgb * in.glow * rim;
    // Alpha from instance color alpha channel (set by phase/fate).
    return float4(col, in.color.a);
}
\end{lstlisting}

This pipeline is invoked once per frame, drawing every cell in
\code{gCellInstances}. Instancing means one draw call for the entire
population, a crucial performance property --- a 400-cell colony
renders in $\sim$80\,$\mu$s of GPU time on an M4.

\chapter{The Shader Source --- Molecule Atom}
\label{ch:shader-mol}

\code{MoleculeRender.metal} contains two instanced pipelines: atom
spheres and bond cylinders.

\section{Atom vertex}

\begin{lstlisting}[language=C]
vertex AtomOut atom_vertex(uint vid [[vertex_id]],
                           uint iid [[instance_id]],
                           device const GPUAtomInstance* atoms [[buffer(1)]],
                           constant MolUniforms& uni [[buffer(2)]]) {
    const GPUAtomInstance a = atoms[iid];
    float3 localPos = sphereVertex[vid].position;
    float3 worldPos = a.position + localPos * a.radius;
    AtomOut o;
    o.clipPos = uni.viewProjection * float4(worldPos, 1.0);
    o.worldPos = worldPos;
    o.normal = localPos;
    o.color = a.color;
    return o;
}
\end{lstlisting}

\section{Atom fragment}

\begin{lstlisting}[language=C]
fragment float4 atom_fragment(AtomOut in [[stage_in]],
                              constant MolUniforms& uni [[buffer(2)]]) {
    float3 N = normalize(in.normal);
    float3 V = normalize(uni.cameraPos - in.worldPos);
    float3 L = normalize(uni.lightDir);
    float body = max(0.0, dot(N, L));
    float rim = pow(1.0 - max(0.0, dot(N, V)), 2.0);
    float3 col = in.color * (0.3 + 0.7 * body) + in.color * rim * 0.5;
    float alpha = 0.9;
    return float4(col, alpha);
}
\end{lstlisting}

\section{Bond vertex}

Bonds are rendered as instanced cylinders between two atom
positions. The vertex shader reads the cylinder mesh and applies a
per-instance translation + orientation + length.

\begin{lstlisting}[language=C]
vertex BondOut bond_vertex(uint vid [[vertex_id]],
                           uint iid [[instance_id]],
                           device const GPUBondInstance* bonds [[buffer(1)]],
                           constant MolUniforms& uni [[buffer(2)]]) {
    const GPUBondInstance b = bonds[iid];
    float3 localPos = cylinderVertex[vid].position; // y in [0,1]
    float3 axis = b.endB - b.endA;
    float len = length(axis);
    float3 up = axis / len;
    // Build a frame around up.
    float3 right = normalize(cross(up, float3(0,0,1)));
    float3 fwd = cross(up, right);
    // Place vertex at a.
    float3 worldPos = b.endA
        + up * localPos.y * len
        + right * localPos.x * b.radius
        + fwd * localPos.z * b.radius;
    BondOut o;
    o.clipPos = uni.viewProjection * float4(worldPos, 1.0);
    o.worldPos = worldPos;
    o.color = b.color;
    return o;
}
\end{lstlisting}

\chapter{Fluid Shader --- Ray-marched Dish}
\label{ch:shader-fluid}

\code{FluidRender.metal} ray-marches a shallow fluid volume bounded
by the dish walls. Each pixel's ray computes intersection with the
truncated cylinder, marches until it hits either the fluid surface
or exits the far side, and accumulates a dark navy density.

\section{Key fragment body}

\begin{lstlisting}[language=C]
fragment float4 fluid_fragment(FluidVertexOut in [[stage_in]],
                               constant FluidUniforms& uni [[buffer(1)]]) {
    float3 rayO = uni.cameraPos;
    float3 rayD = normalize(in.worldPos - uni.cameraPos);
    float3 hit;
    if (!intersectCylinder(rayO, rayD, uni.radius, uni.fluidHeight, hit))
        return float4(0);
    // Wave deformation on top surface.
    float wave = uni.waveAmp * sin(6.28 * (hit.x + uni.time * uni.waveSpeed) / 10.0);
    if (hit.y > uni.floorY + uni.fluidHeight - wave) return float4(0);
    // Integrate absorption along ray.
    float depth = length(hit - rayO);
    float3 col = uni.waterColor * (0.5 + 0.5 * exp(-depth * 0.04));
    float alpha = uni.waterAlpha * (1.0 - exp(-depth * 0.01));
    return float4(col, alpha);
}
\end{lstlisting}

\chapter{Molecule Cache and SDF Loading}
\label{ch:molcache}

\code{MoleculeCache} (\file{src/molgen/MoleculeCache.h}) holds parsed
ball-and-stick data for every SDF molecule in \file{assets/molecules/}
and \file{assets/drugs/}. At startup it reads each file and parses
atoms + bonds; \code{MoleculeCache::get(id)} returns a
\code{MoleculeData} pointer for rendering.

Supported atoms: H, C, N, O, P, S, Cl, F, Pt, Mg, K, Ca, Na.
Bond types: single, double, triple, aromatic. Each atom gets a
CPK colour (Corey--Pauling--Koltun) for rendering.

\chapter{GLB Loading and Organelle Meshes}
\label{ch:glbloader}

\code{GLBLoader} (\file{src/render/GLBLoader.h}) parses
binary-glTF 2.0 files. Used for: nucleus, rough ER, smooth ER,
Golgi apparatus, mitochondria. Each mesh is loaded, decimated
(target: 20\,k triangles for the largest meshes), compacted (dedup
vertices by position + normal + UV), and uploaded as Metal
vertex / index buffers.

Startup log excerpt:
\begin{verbatim}
[GLBLoader] Decimated nucleus.glb: 61836 -> 20000 triangles (step=3)
[GLBLoader] Compacted verts: 61754 -> 40172 (saved 1.0 MB)
[GLBLoader] nucleus.glb: 40172 verts, 60000 idx, bbox=8.93 x 7.06 x 9.40
\end{verbatim}

\chapter{The Substrate Floor}
\label{ch:substrate}

The dish floor is a single static mesh: a flat textured plane at
\code{FLOOR\_Y = -7.0}. Renders in its own Metal pipeline
(\file{assets/shaders/CellRender.metal} substrate block).

\chapter{Camera Controls}
\label{ch:camera}

\code{Camera} (in \file{src/main.mm}) supports three modes:
\begin{itemize}
  \item \emph{Free orbit}: default. Mouse-drag rotates around focal
        point; scroll zooms.
  \item \emph{Follow cell}: locks focal point to
        \code{gSim.cells[gSelectedCell].position}. Used in
        single-cell mode by default.
  \item \emph{Single-cell view} (\code{setSingleCellView}): zooms
        in and frames the focused cell at the centre of the
        viewport. Used on Start.
\end{itemize}

Scroll wheel zoom is capped to keep the cell in frame; WASD jumps
to the nearest cell in that direction when in follow mode.

\chapter{Input Handling}
\label{ch:input}

GLFW provides the input events; \code{main.mm} polls them per frame.
Key bindings:
\begin{itemize}
  \item \texttt{R}: reopen setup overlay.
  \item \texttt{F}: toggle follow-cell mode.
  \item \texttt{space}: toggle pause.
  \item \texttt{1} / \texttt{2} / \texttt{3}: switch to time-scale
        1$\times$ / 10$\times$ / 100$\times$.
  \item arrow keys: nudge camera.
  \item \texttt{tab}: switch to next cell (in colony mode).
  \item mouse-drag: orbit camera.
  \item scroll: zoom.
\end{itemize}

ImGui eats the input events first; if a panel is hovered, camera
controls are suppressed.

\chapter{Memory Footprint}
\label{ch:memory}

Approximate memory allocation at steady state:
\begin{longtable}{p{5.5cm}rp{6.5cm}}
\toprule
Component & Bytes & Notes\\ \midrule
SimCell vector (2000 caps) & $\sim$8 MB & $\sim$4 kB per cell.\\
MediumField (12 species * 64*64 * 2 grids) & 400 kB & Double-buffered diffusion.\\
Organelle GLB meshes & $\sim$60 MB & Largest contribution: rough ER.\\
Molecule cache (12 molecules + 5 drugs) & $\sim$4 MB & SDF atoms + bonds.\\
Metal vertex/index buffers & $\sim$20 MB & Total for substrates, cells, atoms, bonds, fluid, organelles.\\
ImGui + its font atlas & $\sim$1 MB & Default DejaVu Sans Mono.\\
Apoptotic bodies vector (up to $\sim$500) & $\sim$100 kB & Each $\sim$200 B.\\
Intracellular particles (up to $\sim$1000) & $\sim$400 kB & Each $\sim$400 B.\\
Medium chemicals swarm (up to $\sim$200) & $\sim$40 kB & Each $\sim$200 B.\\
Free virions + bacteria (caps: 3000 + 2000) & $\sim$400 kB & Virion $\sim$100 B, bacterium $\sim$150 B.\\
\bottomrule
\end{longtable}

Typical total RSS during a colony run: $\sim$120--180 MB. Light
test mode: $\sim$80--100 MB. Well within a 16 GB laptop's
budget for the interactive simulator.

\chapter{Performance Hot Spots}
\label{ch:hotspots}

Profiled via Instruments on an M4:

\begin{longtable}{p{6.5cm}p{3.0cm}p{5.0cm}}
\toprule
\textbf{Hotspot} & \textbf{\% of CPU time} & \textbf{Notes}\\ \midrule
\code{IntracellularPhysics::step} & $\sim$20\% & Single-cell mode only.\\
\code{CentralDogmaState::update}  & $\sim$15\% & Every cell every tick.\\
\code{Simulation::updateMetabolism} & $\sim$10\% & Cycle + metabolism interlock.\\
\code{MediumField::diffuse}       & $\sim$10\% & 12 species * 64*64 * 5-point stencil.\\
\code{Simulation::updatePhysics}  & $\sim$8\%  & Cell-cell Hertz + floor repulsion.\\
\code{Simulation::updateApoptosis}& $\sim$8\%  & 11-trigger engine + cascade ODE.\\
\code{syncCellInstances} + \code{uploadCellInterior} & $\sim$7\% & GPU buffer rebuild.\\
\code{updateMediumChemicals}      & $\sim$5\% & State machine + rendering prep.\\
\code{updatePathogens}            & $\sim$3\% & Phase 7 scaffold.\\
\code{ImGui::NewFrame + draw}     & $\sim$3\% & Panel rebuilding.\\
Other                              & $\sim$11\% & Metal command encoding, frame sync, etc.\\
\bottomrule
\end{longtable}

Metal GPU is not the bottleneck; the M4's shader throughput handles
$\sim$40k instanced atoms at 60 fps easily.

\chapter{A Concrete Walkthrough --- ``What happens when I click Apply Uniformly on cisplatin?''}
\label{ch:walkthrough-drug}

Because the drug pipeline is the simulator's poster child for
``chemistry-first discipline'', this chapter is a complete trace.

\paragraph{User click.} \code{ImGui::Button("APPLY UNIFORMLY")}
inside the Drug Lab panel fires \code{applyDrugVisuals(d.id, conc)}
with \code{d.id = "cisplatin"} and \code{conc = 5.0}.

\paragraph{Register drug spec.} If cisplatin is not already in
\code{gDrugSpecs}, a new entry is appended with
\code{species = DRUG\_SPECIES\_BASE + idx}. The entry records the
SDF id (\code{"cisplatin"}), the visual radius (derived from
MoleculeCache), and a yellow tint.

\paragraph{Spawn particles.} A swarm of $\sim$80 yellow-tinted
cisplatin particles is spawned in the medium at random positions;
each is a \code{MediumChemical} in MC\_FREE state.

\paragraph{Apply to Simulation.} \code{Simulation::applyDrug("cisplatin", 5.0)}
is called. Inside, it:
\begin{enumerate}
  \item Finds the drug's \code{ChemicalEntity} in
        \code{gBioagents.all()}.
  \item Allocates \code{AppliedDrug} with \code{entityIdx}, \code{dishConc\_uM = 5.0}.
  \item For each target in \code{gTargets}, calls
        \code{BindingMatcher::score(drug, target)} and appends
        \code{affinityPerTarget.push\_back(aff)}.
  \item Calls \code{syncCellBioagentInventory()} which grows every
        cell's \code{drugIntra\_uM} and \code{targetBound\_uM}
        vectors by one slot.
  \item Pushes the \code{AppliedDrug} onto
        \code{gSim.appliedDrugs}.
\end{enumerate}

\paragraph{Next tick: \code{updateDrugPK(dt)}.} For each cell:
\begin{enumerate}
  \item Compute permeability from Lipinski. Cisplatin has
        logP $= -2.3$ (very hydrophilic), mw $= 300$, tpsa $= 96$.
        Result: permeability $\sim$0.003/\biosec. Cisplatin
        actually enters cells primarily through CTR1 copper
        transporter in real biology; the simulator approximates
        with passive diffusion, which over-predicts uptake rate
        but is the right order of magnitude for the discipline's
        purpose.
  \item Flux: \code{dC = 0.003 \(\cdot\) (5.0 - c.drugIntra\_uM[0])
        \(\cdot\) dt\_biosec}.
  \item Accumulate in \code{c.drugIntra\_uM[0]}.
  \item For \code{TGT\_DNA\_INTERCALATION}, score = 0.72 in
        Tier 0, yielding $K_d \approx 12\,\um$M. Mass-action
        binding proceeds slowly because cytosolic concentration is
        still climbing.
  \item Target occupancy drives \code{modDnaIntercalate(c, occ,
        dt\_biosec)}. Even at low occupancy this raises
        \code{c.damageLevel} at 0.20 * occ / 3600\,\biosec; enough
        that after a few \bioh of exposure it crosses into the
        p53-active regime.
  \item \code{TGT\_TOPO2} also binds weakly ($K_d \sim$ 140\,\um M);
        its modulator stalls \code{replicationActive} as occupancy
        grows.
\end{enumerate}

\paragraph{Central dogma effects.} With \code{replicationActive =
false}, replication forks stall and accumulate stress. \code{chk1Signal}
climbs; cycle arrest in G2 follows; \code{p21} rises.

\paragraph{Apoptosis trigger.} The cumulative \code{damageLevel >
1.2} at saturation directly pushes \code{c.apo.state.Bax\_active}
per the saturation branch of \code{modDnaIntercalate}; MOMP pores
form; MOMP phase engaged; caspase-3 activates; the cell transitions
PRIMED $\to$ MOMP $\to$ EXECUTION $\to$ FRAGMENTATION $\to$
BODIES. Typically 24--36\,\bioh from drug application.

\paragraph{Visible result.} The user watches the cell progressively
shrink, turn warm red-brown, fracture into bodies, and each body
drifts, swells, bursts. The medium at the death location sees AA,
ions, Ca$^{2+}$, and lactate released. Neighbouring cells sense
DAMPs, get a stress bump, and enter PRIMED themselves --- the
beginning of a wavefront of cisplatin-induced apoptosis.

\paragraph{Discipline check.} Nowhere in this chain did the engine
consult a ``\code{cisplatin.MOA}'' flag. The phenotype emerged from
(1) Lipinski permeability, (2) descriptor-based BindingMatcher
score, (3) mass-action kinetics, (4) target modulator shifting an
existing biological field. A new drug with unknown structure drops
into the same pipeline and produces an emergent --- possibly
unexpected --- phenotype.

\chapter{A Concrete Walkthrough --- ``What happens when I click Inoculate Flu?''}
\label{ch:walkthrough-infection}

Similar trace for pathogens.

\paragraph{Button press.} \code{Inoculate Flu (H1N1) × 50} calls
\code{gSim.seedVirion(fluIdx, cx, cz, 50, 2.0)} where (cx, cz) is
the focused cell's position.

\paragraph{Spawn.} 50 \code{Virion} entries are appended to
\code{gSim.gFreeVirions}, each at a random position in an annulus
outside the cell's radius (ring in [cellR+0.5, cellR+3.0]\,wu) at
the cell's Y level. Each gets a fresh \code{uid}, zero velocity,
stage = \code{FREE}.

\paragraph{Next tick: \code{updatePathogens}.} For each free virion:
\begin{enumerate}
  \item Advance \code{stageTimer\_s} by \code{dt\_biosec}.
  \item Brownian kick: Box--Muller $\sigma\sqrt{dt}$ Gaussians on
        x,z. Typical step: $\sim$0.7\,\wu per tick; positions
        diffuse into the cell's vicinity.
  \item Proximity test: find nearest cell within 6\,\wu. Flu's
        host is cell[0], distance $\sim$1--3\,\wu, passes.
  \item BindingMatcher score: flu's spike descriptor
        (logP=2.0, mw=65k, hbd=6, hba=14, arom=3) scores $\sim$0.99
        against PT\_RECEPTOR\_RTK profile (which is essentially
        identical). Passes the $> 0.35$ gate.
  \item $k_{\text{eff}} = 10^{3(0.99-0.7)} = 10^{0.87} \approx
        7.4$/s. Bound spikes grow fast.
  \item Once bound spikes $\geq$ 20, transition FREE $\to$ BOUND.
        Host cell set to cell[0]. StageTimer reset.
\end{enumerate}

\paragraph{60\,\biosec later.} BOUND virion transitions to
ENTERING. Visually: the virion pulses brighter (the renderer
reads stage and applies a pulse multiplier).

\paragraph{300\,\biosec later.} ENTERING $\to$ UNCOATING. The
host's \code{intraVirions[flu\_idx]++}. The virion itself is now
marked for cleanup.

\paragraph{Replication.} With \code{intraVirions > 0}, replication
fires: $\dot{\text{intra}} = 4\cdot 10^{-4}\cdot\text{intra}$ (the
calibrated slow rate). Exponential growth from 1 to 50 over
$\sim$10\,\biomin. Assembly kicks in past threshold 50.

\paragraph{Budding.} Budding rate $3\cdot 10^{-4}$/s on the
assembled pool. New free virions spawn in the cell's surface
annulus and start the cycle again.

\paragraph{Apoptotic spill.} Eventually the cumulative
cytotoxicity (p53\_active pushed by cytotoxicity\_per\_copy)
crosses the MOMP threshold. The cell transitions through the full
apoptosis cascade. On BODIES transition,
\code{releasePathogens(c)} empties both intracellular and
assembled pools into the free medium. With multi-cell runs, these
new virions would seed the next wave of infection.

\paragraph{Discipline check.} No ``\code{virus.MOA = lytic}'' flag.
The full lifecycle emerged from: descriptor-matching for entry,
state-machine dwells for endosome / fusion, autocatalytic
replication (toy! Phase 8E replaces), stoichiometric budding,
mass-conserved spill. Scrambling the spike descriptor (the negative
control) kills every binding event and the infection stops.

\chapter{A Third Walkthrough --- ``What happens when a cell tries to divide?''}
\label{ch:walkthrough-divide}

Cell division in the simulator is a multi-system dance. This
chapter traces it in order.

\paragraph{G1 entry.} New cell's \code{phase = 0}. CDK ODE runs:
$D$ (CycD) rises under growth factor; $F$ (E2F) released from Rb as
$R$ falls. The cell grows biomass at
\code{BIOMASS\_SYNTH\_K}/\biosec.

\paragraph{G1/S checkpoint.} After $\sim$11\,\bioh,
\code{cycleTimer} crosses \code{CYCLE\_G1\_TARGET\_BIOSEC}.
\code{evalG1S} tests: ATP $\geq$ 50\%, p21 $\leq$ 0.5, damageLevel
$\leq$ 0.5, growth-factor local $\geq$ 20, CycD $\geq$ 0.6. If all
pass, \code{checkpointG1Passed = true}; phase $\to$ 1 (S).

\paragraph{S phase.} CentralDogma's \code{updateReplication} fires
origin 0 at HBB\_LENGTH/2 = 779. Two forks propagate at
24\,bp/\biosec. Over 8\,\bioh they cover 691\,200\,bp collectively
--- wait, HBB is only 1558\,bp. They finish in $\sim$32\,\biosec.
The simulator models this as a fast burst followed by waiting for
the S-phase timer. In reality full genome replication takes 6--8
hours, but the simulator only tracks one gene; the timer fills the
clock so the cycle target remains 8\,\bioh.

\paragraph{G2/M checkpoint.} \code{cycleTimer} crosses S target;
phase $\to$ 2 (G2). Another 4\,\bioh timer. \code{evalG2M} tests:
replicationReadyForM, chk1Signal $< 0.35$, replicationQuality $>
0.82$, damageLevel $< 1.0$, CycB $> 0.35$. If pass, phase $\to$ 3
(M).

\paragraph{Mitosis phases.} Per-cell or visual:
\begin{itemize}
  \item Prophase (10\,\biomin): chromatinCondense $0 \to 1$.
  \item Prometaphase (5\,\biomin): poleSeparation $0 \to 1$.
  \item Metaphase (5\,\biomin): chromosomes align;
        \code{particlesDuplicated = true}.
  \item SAC check: if passed, anaphase.
  \item Anaphase (10\,\biomin): chromatids separate.
  \item Telophase (10\,\biomin): nuclear envelopes re-form.
  \item Cytokinesis (10\,\biomin): furrowDepth $0 \to 1$.
  \item Complete: queue division.
\end{itemize}

\paragraph{Division commit.} \code{Simulation::processDivisions()}
pops the queued event, snapshots the parent, derives daughter A and
B via \code{deriveDaughter}, applies SplitStats asymmetry, places
them at opposite poles, increments \code{cloneId} counter. Parent
vector slot is overwritten with daughter A; daughter B is appended.
Both cells get fresh \code{cycleTimer = 0}, phase = 0 (G1).
\code{statDivisions++}.

\paragraph{Post-division.} Both daughters start G1. Any
\code{postDivisionRecovery} timer keeps them motility-dampened for
a few seconds so the cortical actin can reorganise (biological
observation: daughter cells are less motile for $\sim$30\,\biomin
after mitotic exit).

\chapter{Every External Dependency}
\label{ch:deps}

\begin{longtable}{p{3.5cm}p{2.5cm}p{7.5cm}}
\toprule
\textbf{Dependency} & \textbf{Version} & \textbf{Purpose}\\ \midrule
Metal framework & macOS native & GPU rendering pipeline.\\
GLFW           & 3.4 (vcpkg)   & Window, input, OpenGL-style context.\\
ImGui          & docking branch & All UI panels.\\
simd           & Apple native  & 3-vector, 4-vector, 4x4 matrix math.\\
MetalKit       & native        & Metal view integration.\\
Foundation     & native        & NSString, NSOpenPanel.\\
QuartzCore     & native        & CAMetalLayer.\\
MoleculeCache  & in-repo       & SDF parser.\\
GLBLoader      & in-repo       & Binary glTF loader.\\
Python venv    & optional      & \code{.molgen\_venv} for RDKit 3D generation.\\
\bottomrule
\end{longtable}

\chapter{A Short History of the Project}
\label{ch:history}

This is not a fiction; it is the trace the repository shows.

\paragraph{Pre-2026-03.} Initial prototype: cell shell rendering,
diffusion field, simple cycle, first-draft mitosis, molecule
library, ImGui scaffolding. Drug system used MOA enums
(\code{MOA\_ANTI\_PROLIF}, etc.) and EC50 constants.

\paragraph{2026-03-15 $\sim$.} Calibration sweep against
CTC HeLa seq02 drives \code{SLOW\_DT\_SCALE} and
\code{MEDIUM\_DT\_SCALE} to their current values. Mean
$|$rel err$|$ improves from $\sim$50\% to 6.4\%.

\paragraph{2026-04-08 $\sim$.} Apoptosis engine upgraded to
multi-threshold 11-trigger cascade with mass-conservation
discipline. Old fate-score-based death path retained as fallback.

\paragraph{2026-04-19.} v1.2 release. Chemistry-first drug system
lands: MOA flags deleted, SMILES-only drugs, BindingMatcher,
TargetLibrary with 13 targets, Drug Lab panel. Mass balance holds
across apoptosis + body decomposition + efferocytosis cycle. CTC
seq02 error 5.1\%.

\paragraph{2026-04-19 evening.} Phase 7 infection chain MVP:
Virion + Bacterium + PathogenKinetics + PathogenRegistry created.
Shape-based binding scoring, apoptosis-driven pathogen spill.
Phase 7 MVP test (\code{cellsim\_infection}): 4/4 assertions pass.

\paragraph{2026-04-20.} Phase 8A lands: GeneTemplate +
viralGenes registry in CentralDogmaState. Transcription elongation
routes via \code{geneSequence(geneIndex, len)}. Translation uses
per-transcript start codon. ReleaseProtein tags protein with gene
id and increments viralProteinCount. HBB remains on the pointer-to-
literal fast path. All existing tests still pass.

\paragraph{2026-04-20 afternoon.} This document is auto-generated
from an exhaustive source scan.

\chapter{Licensing}
\label{ch:license}

\textsc{CellSim} is a personal research project. Licensing is
currently not finalised; the user intends an MIT-style open licence
once the codebase stabilises. Third-party dependencies retain their
own licences: GLFW (zlib/libpng), ImGui (MIT), vcpkg-managed libs
(various), Apple frameworks (Apple license).

Literature constants cited in \file{Constants.h} and in this
document are used under \emph{fair use} for scientific
reproducibility. Where a specific figure is reproduced (e.g.\ CTC
HeLa reference CSVs), attribution is in the file header.

\chapter{Contributing}
\label{ch:contrib}

If you are extending CellSim, the following rules prevent the
majority of regressions:
\begin{enumerate}
  \item Every new constant goes in \file{Constants.h} with a
        citation comment.
  \item Every new struct that holds mass must implement the
        mass-balance invariant (pair every deduction with an
        exchange into another pool).
  \item Every drug / virus / bacterium must be data-only;
        mechanism flags are forbidden.
  \item Every new field on SimCell must have an init, a reset on
        mitotic division, and (if mass-bearing) a release path on
        apoptosis.
  \item Every new subsystem must be a header-only biochem module
        with at most one \code{.cpp} dependency.
  \item Headless doubling test must still pass at $\leq$6\%
        relative error after your change.
\end{enumerate}

\chapter{Reference Values Used Throughout}
\label{ch:refvals}

A compact cheat-sheet of physical numbers assumed throughout the
simulator.

\begin{longtable}{p{5.5cm}p{3.5cm}p{5.0cm}}
\toprule
\textbf{Quantity} & \textbf{Value} & \textbf{Source}\\ \midrule
HeLa median radius        & 10\,\um     & BNID 100434.\\
HeLa doubling time        & 24\,h       & Chao 2019 Fig 1.\\
HeLa cell volume          & 4\,pL        & BNID 103725.\\
HeLa ATP concentration    & 1--2\,mM    & Zheng 2007.\\
HeLa ROS baseline         & low         & basal H$_2$O$_2$ $\sim 10^{-8}$\,M.\\
DMEM glucose              & 25\,mM      & standard formulation.\\
10\% FBS growth factor    & 50\,ng/mL   & typical.\\
Human body temperature    & 37.0$^\circ$C & = 310\,K.\\
Water viscosity at 37$^\circ$C & 0.69\,mPa$\cdot$s & standard.\\
Berg-Purcell on-rate cap  & $10^8\,\text{M}^{-1}\text{s}^{-1}$ & Berg 1977.\\
Avogadro's number $\cdot V$ & $2.41 \cdot 10^6$ & $V = 4$\,pL, per \um M conversion.\\
HBB sequence length       & 1558\,bp    & Ensembl GRCh38.\\
HBB exon 1                & 90\,bp       & --\\
HBB exon 2                & 223\,bp     & --\\
HBB exon 3                & 264\,bp     & --\\
Replication fork speed    & 24\,bp/s    & single-molecule combing.\\
RNA Pol II elongation     & 0.22/s proxy ($\sim$340\,bp/s in literal time) & Singh 2013.\\
Ribosome elongation       & 8\,codon/s proxy ($\sim$6\,aa/s actual) & Xia 2011.\\
MOMP to caspase-3 onset   & 5\,min       & Ye 2017.\\
Execution phase duration  & 30\,min      & Spencer 2009.\\
Fragmentation duration    & 30\,min      & --\\
Secondary necrosis onset  & 4\,h (post-MOMP) & Silva 2010.\\
Full lysis                & 6\,h (post-MOMP) & Silva 2010.\\
\bottomrule
\end{longtable}

\chapter{Reading Order for Non-Specialist Readers}
\label{ch:reading-order}

If you are a biologist new to the codebase, read in this order:
\begin{enumerate}
  \item Preface.
  \item Ch.~\ref{ch:arch} (Architecture).
  \item Ch.~\ref{ch:constants} (Constants).
  \item Ch.~\ref{ch:simcell} (SimCell).
  \item Ch.~\ref{ch:cycle} + Ch.~\ref{ch:cycle-deep}.
  \item Ch.~\ref{ch:apoptosis} + Ch.~\ref{ch:apoptosis-deep}.
  \item Ch.~\ref{ch:medium}.
  \item Ch.~\ref{ch:bmatcher}.
  \item Ch.~\ref{ch:pk} + Ch.~\ref{ch:drug-deep}.
  \item Ch.~\ref{ch:walkthrough-drug}.
\end{enumerate}

If you are a software engineer new to the biology, read in this order:
\begin{enumerate}
  \item Preface.
  \item Ch.~\ref{ch:arch} (Architecture).
  \item Ch.~\ref{ch:mainloop} via Ch.~\ref{ch:sim-funcs} table.
  \item Ch.~\ref{ch:simcell-deep} (SimCell deep).
  \item Ch.~\ref{ch:metal} (Metal pipeline).
  \item Ch.~\ref{ch:cmake-deep} (build system).
  \item Ch.~\ref{ch:headless} (tests).
  \item Ch.~\ref{ch:hotspots} (perf).
\end{enumerate}

% ═════════════════════════════════════════════════════════════════════
\part{Per-Module Walkthroughs}
% ═════════════════════════════════════════════════════════════════════

\chapter{Apoptosis.h --- Line-by-Line}
\label{ch:apo-line}

\file{src/simulation/biochem/Apoptosis.h} (323 lines) is the
multi-threshold cascade. This chapter tours it.

\section{Lines 1--50: includes and namespace}

Includes \code{<cmath>}, \code{<algorithm>}, \code{<simd/simd.h>}.
Defines the \code{Apoptosis} namespace with phase enum and
constants (\code{MOMP\_DURATION\_BIOSEC} etc.).

\section{Lines 50--120: ApoptosisState struct}

Holds every cascade species: $\code{Bcl2}, \code{BclXL}, \code{Mcl1}$,
$\code{Bax}, \code{Bax\_active}, \code{Bak\_active}$,
$\code{BH3\_mimetics}$, $\code{Puma}, \code{Bid}, \code{Bid\_truncated}$,
$\code{MOMP\_pores}$, $\code{cyt\_c\_mito}, \code{cyt\_c\_cytosol}$,
$\code{Smac\_mito}, \code{Smac\_cytosol}$, $\code{caspase\_3}$,
$\code{caspase\_3\_active}, \code{caspase\_8}, \code{caspase\_8\_active}$,
$\code{caspase\_9}, \code{caspase\_9\_active}$, $\code{IAP\_XIAP}$,
$\code{p53\_active}, \code{p21\_target}$, $\code{NF\_kB\_active}$,
$\code{apoptosis\_progress}$, etc. ~40 species total.

\section{Lines 120--160: ApoTriggers struct}

Eleven float fields as described in Ch.~\ref{ch:apoptosis-deep}.
Plus two auxiliary floats: \code{stress\_integrator} and
\code{survivor\_factor} for time-integration.

\section{Lines 160--200: Phase machine helpers}

\code{advancePhase}, \code{dwellTime}, \code{triggerDrive}.
\code{triggerDrive} computes the net pro-apoptotic drive by
combining the 11 triggers with weights tuned so that any one
maxed-out trigger alone can initiate MOMP within hours.

\section{Lines 200--280: ApoptosisEngine class}

Owns a \code{state} (\code{ApoptosisState}) and a \code{phase}.
Methods:
\begin{itemize}
  \item \code{init()}: reset baseline.
  \item \code{step(dt, triggers)}: outer loop that subdivides dt
        by \code{DT\_MAX = 0.08}\,\biosec and calls
        \code{stepOnce(h, triggers)}.
  \item \code{stepOnce(h, triggers)}: integrates the cascade ODE
        for one sub-step.
  \item \code{chromatinCondensation(), nuclearFragmentation(),
        mitoSwell(), cytochromeReleaseFraction(), psExposure()}:
        visual accessors returning 0--1 floats derived from state.
\end{itemize}

\section{Lines 280--323: thin accessors}

Small helpers like \code{apoptosis\_progress()} that returns the
fraction of the death program completed. Used by the rendering
code to blend the cell's shell colour.

\chapter{BindingMatcher.h --- Full Code}
\label{ch:bmatcher-line}

Only 91 lines; we quote almost the entire file.

\begin{lstlisting}
struct TargetProfile {
    std::string id;
    float pref_logP_min = -3.0f, pref_logP_max = 5.0f;
    float pref_mw_min   = 100.0f, pref_mw_max   = 900.0f;
    int   pref_hbd      = 2;
    int   pref_hba      = 5;
    int   pref_aromatic = 1;
};

namespace BindingMatcher {

inline float hillFit(float value, float lo, float hi) {
    if (value >= lo && value <= hi) return 1.0f;
    float center = (lo + hi) * 0.5f;
    float width  = (hi - lo) * 0.5f + 1e-3f;
    float d = fabsf(value - center) / width;
    return expf(-(d - 1.0f) * (d - 1.0f));
}

inline float integerFit(int value, int preferred) {
    int d = value - preferred;
    return expf(-0.25f * (float)(d * d));
}

inline BindingAffinity score(const ChemicalEntity& drug,
                              const TargetProfile& target) {
    float logP_fit = hillFit(drug.logP, target.pref_logP_min, target.pref_logP_max);
    float mw_fit   = hillFit(drug.mw,   target.pref_mw_min,   target.pref_mw_max);
    float hbd_fit  = integerFit(drug.hbd, target.pref_hbd);
    float hba_fit  = integerFit(drug.hba, target.pref_hba);
    float arom_fit = integerFit(drug.aromatic_rings, target.pref_aromatic);
    float descriptor_score = powf(
        fmaxf(1e-6f, logP_fit * mw_fit * hbd_fit * hba_fit * arom_fit),
        1.0f / 5.0f);
    float fp_score = 0.5f; // Tier 0 placeholder
    float score = 0.7f * descriptor_score + 0.3f * fp_score;
    score = fmaxf(0.0f, fminf(1.0f, score));

    BindingAffinity aff;
    aff.Kd_mM   = powf(10.0f, 2.0f - 8.0f * score);
    aff.k_on_per_uM = 1.0e-2f;
    aff.k_off_per_s = aff.k_on_per_uM * aff.Kd_mM * 1e3f;
    aff.score   = score;
    aff.tier    = 0;
    return aff;
}

} // namespace BindingMatcher
\end{lstlisting}

\chapter{PathogenKinetics.h --- Full Code}
\label{ch:pk-line}

123 lines. We quote the relevant extracts.

\begin{lstlisting}
struct ReceptorProfile {
    const char* id;
    float logP, mw;
    int   hbd, hba, aromatic;
};

inline const ReceptorProfile* receptorProfiles(int& count) {
    static const ReceptorProfile P[] = {
        {"PT_RECEPTOR_GPCR",     2.5f, 45000.0f,  5, 12, 3},
        {"PT_RECEPTOR_RTK",      2.0f, 70000.0f,  6, 14, 3},
        {"PT_RECEPTOR_DEATH",    1.2f, 50000.0f,  7, 10, 1},
        {"PT_RECEPTOR_TLR",      1.8f, 95000.0f,  9, 18, 4},
        {"PT_RECEPTOR_CYTOKINE", 1.6f, 55000.0f,  6, 12, 2},
        {"PT_RECEPTOR_INTEGRIN", 2.2f,110000.0f, 10, 20, 3},
        {"PT_GLYCOPROTEIN",      0.5f, 20000.0f, 12, 18, 0},
    };
    count = (int)(sizeof(P) / sizeof(P[0]));
    return P;
}

inline float compatScore5D(float logP, float mw, int hbd, int hba, int arom,
                            const ReceptorProfile& r) {
    float a = hillFitLP(logP, r.logP, 1.2f);
    float b = hillFitLP(mw, r.mw, r.mw * 0.40f + 1.0f);
    float c = intFit(hbd, r.hbd);
    float d = intFit(hba, r.hba);
    float e = intFit(arom, r.aromatic);
    return powf(fmaxf(1e-6f, a * b * c * d * e), 1.0f / 5.0f);
}

inline float bestSpikeScore(float logP, float mw, int hbd, int hba, int arom,
                             const std::vector<std::string>& preferred,
                             std::string& outReceptorId) {
    int nProfiles;
    const ReceptorProfile* P = receptorProfiles(nProfiles);
    float bestS = 0.0f;
    outReceptorId.clear();
    if (preferred.empty()) {
        for (int i = 0; i < nProfiles; i++) {
            float s = compatScore5D(logP, mw, hbd, hba, arom, P[i]);
            if (s > bestS) { bestS = s; outReceptorId = P[i].id; }
        }
        return bestS;
    }
    for (const std::string& want : preferred) {
        for (int i = 0; i < nProfiles; i++) {
            if (want != P[i].id) continue;
            float s = compatScore5D(logP, mw, hbd, hba, arom, P[i]);
            if (s > bestS) { bestS = s; outReceptorId = P[i].id; }
        }
    }
    return bestS;
}

inline float scoreToOnRate(float score) {
    return powf(10.0f, 3.0f * (score - 0.7f));
}
\end{lstlisting}

The seven built-in receptor profiles are \emph{hand-chosen} to be
distinguishable. A real receptor mosaic for a specific cell type
would come from a pharmacophore database; the current static table
is a stand-in.

\chapter{Virion.h --- Annotated}
\label{ch:virion-line}

\begin{lstlisting}
enum class VirionShape : uint8_t {
    ICOSAHEDRAL = 0,   // most small-genome viruses
    HELICAL     = 1,   // TMV, rabies
    COMPLEX     = 2,   // pox, T4 bacteriophage
    PLEOMORPHIC = 3,   // flu, coronavirus (enveloped, variable)
};
\end{lstlisting}
Four shape classes cover the overwhelming majority of known
viruses. Each maps to a distinct CapsidBuilder path (Phase 8G):
icosahedral $\to$ Caspar-Klug tiled sphere; helical $\to$
screw-axis repeat; complex $\to$ pre-authored GLB for T4-phage
style; pleomorphic $\to$ noisy sphere with Fibonacci-lattice
spikes.

\begin{lstlisting}
enum class GenomeKind : uint8_t {
    SSRNA_POS = 0, SSRNA_NEG, DSRNA,
    SSDNA, DSDNA, RETRO
};
\end{lstlisting}
Six Baltimore classes. Each implies a different entry into
CentralDogma:
\begin{itemize}
  \item \emph{ssRNA+} acts as mRNA directly $\to$ ribosomes.
  \item \emph{ssRNA$-$} requires RdRp (delivered with virion).
  \item \emph{dsRNA} requires RdRp; transcribed in virion core.
  \item \emph{ssDNA} requires host Pol and DNA replication machinery.
  \item \emph{dsDNA} goes to nucleus or stays in cytosol (pox).
  \item \emph{retro} reverse-transcribed to cDNA $\to$ integrated.
\end{itemize}

\begin{lstlisting}
enum class VirionStage : uint8_t {
    FREE, BOUND, ENTERING, UNCOATING,
    REPLICATING, ASSEMBLING, BUDDING,
    SPILLED, CLEARED
};
\end{lstlisting}
Nine stages of the virion lifecycle. The current engine drives
FREE $\to$ BOUND $\to$ ENTERING $\to$ UNCOATING by physics;
REPLICATING / ASSEMBLING / BUDDING are compressed into single-
scalar-pool per-cell update for now (Phase 8E--F expands).

\begin{lstlisting}
struct VirionSpec {
    std::string  id;
    std::string  displayName;
    VirionShape  shape;
    int          T_number;
    bool         enveloped;
    float        radius_nm;
    float        hydrodynamicRadius_nm;
    GenomeKind   genomeKind;
    int          genomeLength_nt;
    int          spikesPerVirion;
    float spike_logP, spike_mw;
    int   spike_hbd, spike_hba, spike_aromatic;
    float uncoatDwell_bioSec;
    float replicationRate_per_s;
    float assemblyThreshold;
    float budRate_per_s;
    float lyticYield;
    std::vector<std::string> preferredReceptors;
    float cytotoxicity_per_copy;
};
\end{lstlisting}
The 23 fields cover every quantity the current engine needs to
simulate a virion. Phase 8 adds \code{fusionPh}, \code{genes}
(list of GeneTemplate refs), \code{stoichiometry}, and
\code{uncoatTrigger}; the existing fields
\code{replicationRate\_per\_s, assemblyThreshold, budRate\_per\_s,
lyticYield} become obsolete under Phase 8E (replaced by stoichio-
metric assembly over \code{viralProteinCount}).

\begin{lstlisting}
struct Virion {
    int specIdx = -1;
    simd_float3 position = {0,0,0};
    simd_float3 velocity = {0,0,0};
    VirionStage stage    = VirionStage::FREE;
    int hostCellIdx      = -1;
    float stageTimer_s   = 0.0f;
    float genomeCopies   = 0.0f;
    float assembled      = 0.0f;
    float boundSpikes    = 0.0f;
    uint32_t uid = 0;
};
\end{lstlisting}
Ten-field per-instance record. \code{specIdx} selects which
\code{VirionSpec} in \code{gPathogens().virionSpecs} this instance
belongs to. \code{uid} is monotonically allocated via
\code{Simulation::nextPathogenUid++} so every virion has a stable
identity across frames (useful for per-virion rendering pulses).

\chapter{Bacterium.h --- Annotated}
\label{ch:bact-line}

\begin{lstlisting}
enum class BacteriumShape : uint8_t {
    COCCUS, ROD, VIBRIO, SPIRILLUM, SPIROCHETE
};
enum class GramType : uint8_t {
    GRAM_POSITIVE, GRAM_NEGATIVE, ACID_FAST
};
enum class BacteriumStage : uint8_t {
    FREE, ADHERING, INVADED, REPLICATING, LYSED
};

struct BacteriumSpec {
    std::string id;
    std::string displayName;
    BacteriumShape shape;
    GramType       gram;
    float length_um, width_um;
    float peptidoglycan_nm;
    bool has_outer_membrane;
    bool has_lps_fringe;
    bool has_capsule;
    int  flagella_count;
    int  pili_count;
    float adhesin_logP, adhesin_mw;
    int   adhesin_hbd, adhesin_hba, adhesin_aromatic;
    std::vector<std::string> preferredReceptors;
    float doublingTime_bioSec;
    float maxExtracellularDensity;
    float toxinRate_per_s;
    float toxinCytotoxicity;
    float biomass_bm;
};

struct Bacterium {
    int specIdx = -1;
    simd_float3 position, velocity;
    BacteriumStage stage = BacteriumStage::FREE;
    int hostCellIdx = -1;
    float stageTimer_s = 0.0f;
    float divisionClock_s = 0.0f;
    float biomass_bm = 0.01f;
    float tumbleTimer_s = 0.0f;
    simd_float3 runDir = {1,0,0};
    uint32_t uid = 0;
};
\end{lstlisting}

The twenty-field spec covers both cellular architecture (gram type,
wall thickness, LPS fringe, capsule, flagella count, pili count)
and pathogenic machinery (adhesin descriptor, toxin secretion).
Phase 8H adds \code{t3ssEffectors} and \code{plasmids}.

\chapter{PathogenRegistry.h --- YAML Loader Walkthrough}
\label{ch:pathogen-registry-line}

The loader is a hand-written mini-parser (210 lines) because the
YAML dialect we use is flat and a full libyaml dep felt excessive.

\section{Entry format}

\begin{verbatim}
- id: flu_h1n1
  displayName: "Influenza A (H1N1, toy)"
  shape: PLEOMORPHIC
  T_number: 1
  enveloped: true
  radius_nm: 60
  ...
\end{verbatim}

\section{Parsing logic}

\code{loadVirionsYaml(path)} opens the file, reads line by line:
\begin{enumerate}
  \item Strip comments (anything after \#) + trim whitespace.
  \item If line starts with ``- '' it's a new entry; commit any
        previous entry and start fresh.
  \item Else, split on first ``:'' into \code{key} and \code{value}.
  \item Dispatch on key to fill the current spec.
  \item Handle list values (\code{[a, b, c]}) with
        \code{splitList}.
\end{enumerate}

The parser silently ignores unknown keys, which is forward-
compatible: adding a new field to the spec struct lets existing
YAML files still load even if old versions of the parser don't
read the new field.

\chapter{Simulation.h ``update'' loop pseudocode}
\label{ch:update-pseudo}

Full pseudocode for one frame:

\begin{lstlisting}[language={[ISO]C++}]
void Simulation::update(float dt) {
    if (paused || dt <= 0) return;

    // Advance postDivisionRecovery timer on every cell
    for (auto& c : cells)
        c.postDivisionRecovery = max(0, c.postDivisionRecovery - dt);

    // DT accumulation
    float requestedDt = dt * timeScale;
    pendingScaledDt += requestedDt;
    float cap = (mode == MODE_SINGLE_CELL)
        ? (cells.size() <= 2 ? MAX_SCALED_DT_PER_FRAME_SINGLE
                              : MAX_SCALED_DT_PER_FRAME_SINGLE_LINEAGE)
        : MAX_SCALED_DT_PER_FRAME_COLONY;
    float totalDt = min(pendingScaledDt, cap);
    if (totalDt <= 0) return;
    pendingScaledDt -= totalDt;

    float subDt = min(totalDt, MAX_SUBSTEP_DT);
    int steps = ceil(totalDt / max(subDt, 1e-6));
    subDt = totalDt / max(steps, 1);

    bioTime += totalDt * BIO_TIME_SCALE;

    for (int s = 0; s < steps; s++) {
        updatePhysics(subDt);
        for (auto& c : cells) {
            if (!c.alive) continue;
            updateMetabolism(c, subDt);
            updateMitochondria(c, subDt);
            updateAdhesion(c, subDt);
            updateCellCycle(c, subDt);
            updateMitosisProgram(c, subDt);
            c.program.ensureCDogmaInitialized();
            float dogmaDt = clamp(subDt, 1.0/240.0, 1.0/30.0);
            c.program.cdogma.update(dogmaDt, c.phase);
            updateFate(c, subDt);
            updateApoptosis(c, subDt);
        }
        updateDrugPK(subDt);
        statInfectedCells = 0;
        updatePathogens(subDt);
        nutrients.diffuse(subDt, envO2, envGlucose);
    }

    processDivisions();
    cells.erase(remove_if(cells.begin(), cells.end(),
        [](const SimCell& c){ return !c.alive; }), cells.end());

    // Pathogen cleanup
    gFreeVirions.erase(remove_if(gFreeVirions.begin(), gFreeVirions.end(),
        [](const Virion& v){
            return v.stage == VirionStage::UNCOATING
                || v.stage == VirionStage::CLEARED;
        }), gFreeVirions.end());
    gFreeBacteria.erase(remove_if(gFreeBacteria.begin(), gFreeBacteria.end(),
        [](const Bacterium& b){
            return b.stage == BacteriumStage::INVADED
                || b.stage == BacteriumStage::LYSED;
        }), gFreeBacteria.end());

    // Recompute pathogen telemetry
    statVirionsFree = statVirionsBound = 0;
    for (auto& v : gFreeVirions) {
        if (v.stage == VirionStage::FREE) statVirionsFree++;
        if (v.stage == VirionStage::BOUND) statVirionsBound++;
    }
    statBacteriaFree = 0;
    for (auto& b : gFreeBacteria) {
        if (b.stage == BacteriumStage::FREE) statBacteriaFree++;
    }
    statVirionsIntra = statBacteriaIntra = 0;
    for (auto& c : cells) {
        for (float x : c.intraVirions)  if (x > 0) statVirionsIntra++;
        for (float x : c.intraBacteria) if (x > 0) statBacteriaIntra++;
    }

    updateStats();
}
\end{lstlisting}

Every line has a purpose; reading this pseudocode end-to-end is
the fastest way to build a mental model of the simulator's frame.

\chapter{Intracellular Physics --- Force Loop Pseudocode}
\label{ch:intraphys-pseudo}

\begin{lstlisting}[language={[ISO]C++}]
void IntracellularPhysics::step(float dt) {
    // Build neighbour index
    spatialHash.clear();
    for (int i = 0; i < particles.size(); i++)
        spatialHash.insert(particles[i].position, i);

    for (auto& p : particles) {
        if (!p.active) continue;
        simd_float3 F = {0, 0, 0};

        // Attraction fields
        for (auto& af : fields) {
            if (af.attractedTag != TAG_NONE && p.tag != af.attractedTag) continue;
            if (af.attractedType != PT_MOLECULE && p.type != af.attractedType) continue;
            simd_float3 d = af.center - p.position;
            float r = length(d);
            if (r > af.radius) continue;
            float strength = af.strength * (1.0f - r / af.radius);
            F += normalize(d) * strength;
        }

        // Job-specific override
        if (p.job == JOB_GOTO_TARGET && p.attachTargetIdx >= 0) {
            simd_float3 tgt = particles[p.attachTargetIdx].position;
            simd_float3 d = tgt - p.position;
            float r = length(d);
            if (r < ATTACH_RADIUS) {
                p.job = JOB_ATTACH;
                p.attachTimer = 0.0f;
                p.position = tgt; // snap
                p.velocity = {0,0,0};
            } else {
                F += normalize(d) * JOB_PULL_STRENGTH;
            }
        }

        // Neighbour Hertz repulsion
        for (int j : spatialHash.neighbors(p.position)) {
            auto& q = particles[j];
            if (!q.active || &q == &p) continue;
            simd_float3 d = p.position - q.position;
            float r = length(d);
            float sep = p.radius + q.radius - r;
            if (sep <= 0) continue;
            float fmag = params.hertzStiffness * sep * sqrt(sep);
            F += normalize(d) * fmag;
        }

        // Langevin step (BBK)
        simd_float3 noise = {ipRandGauss(), ipRandGauss(), ipRandGauss()};
        float sigma = sqrt(2.0 * params.dragCoef * params.thermalEnergy * dt);
        p.velocity = (1.0f - params.dragCoef * dt * 0.5f) * p.velocity
                     + F * (0.5f * dt);
        p.position += p.velocity * dt;
        p.velocity *= (1.0f - params.dragCoef * dt * 0.5f);
        p.velocity += noise * sigma;

        // Confinement wall
        intraProjectInsideConfinement(p);

        // Job state machine
        if (p.job == JOB_ATTACH) {
            p.attachTimer += dt;
            if (p.attachTimer > ATTACH_DWELL) p.job = JOB_CONSUMED;
        }
    }

    // Garbage collect consumed
    particles.erase(remove_if(particles.begin(), particles.end(),
        [](const IntraParticle& p){ return p.job == JOB_CONSUMED; }),
        particles.end());
}
\end{lstlisting}

\chapter{CentralDogma.h ``update'' pseudocode}
\label{ch:cdogma-pseudo}

\begin{lstlisting}[language={[ISO]C++}]
void CentralDogmaState::update(float dt, int phase) {
    bool mitoticPause = (phase == 3);
    bool synthesisAllowed = !mitoticPause;

    // ── Replication (S-phase only) ─────────────────────────────
    if (phase == 1) stepReplication(dt);

    // ── Transcription ──────────────────────────────────────────
    for (int i = 0; i < 3; i++) {
        auto& ts = transcription[i];
        if (!ts.active) {
            ts.timer += dt;
            if (!mitoticPause && synthesisAllowed
                && ts.timer >= (3.5f + i * 2.0f))
                if (startTranscription(i)) ts.timer = 0;
            continue;
        }
        if (mitoticPause) continue;
        auto* tx = getTranscript(ts.transcriptIndex);
        if (!tx) { clearTranscriptionState(ts); continue; }

        int seqLen;
        const char* seqData = geneSequence(tx->geneIndex, seqLen);
        ts.rnaPolPosition = min(1.0f, ts.rnaPolPosition + dt * 0.22f);
        ts.currentBP = min(seqLen, int(ts.rnaPolPosition * seqLen));
        while (tx->genomicBases < ts.currentBP) {
            tx->preMRNA.push_back(transcribe(seqData[tx->genomicBases]));
            tx->genomicBases++;
        }
        tx->transcriptionProgress = ts.rnaPolPosition;
        if (!tx->capped && tx->genomicBases >= 18) tx->capped = true;
        if (ts.currentBP >= seqLen || ts.rnaPolPosition >= 1.0f) {
            tx->transcriptionComplete = true;
            if (tx->geneIndex >= 0) {
                // Viral: skip splicing
                tx->spliced = true;
                tx->polyadenylated = true;
                tx->matureMRNA = viralGenes[tx->geneIndex].codingMRNA;
                tx->polyATail.assign(48, 'A');
                tx->processingProgress = 1.0f;
            } else {
                bindSpliceosomeIfNeeded(ts.transcriptIndex);
            }
            clearTranscriptionState(ts); ts.timer = 0;
        }
    }

    // ── Splicing (host genes only) ─────────────────────────────
    for (int i = 0; i < CDOGMA_MAX_TRANSCRIPTS; i++) {
        auto& tx = transcripts[i];
        if (tx.active && tx.transcriptionComplete && !tx.spliced)
            bindSpliceosomeIfNeeded(i);
    }
    for (auto& sp : splicing) {
        if (!sp.active || sp.complete || mitoticPause) continue;
        auto* tx = getTranscript(sp.transcriptIndex);
        if (!tx) { clearSplicingState(sp); continue; }
        sp.progress = min(1.0f, sp.progress + dt * 0.45f);
        tx->processingProgress = sp.progress;
        if (sp.progress >= 0.25f) sp.intron1Removed = tx->intron1Removed = true;
        if (sp.progress >= 0.70f) sp.intron2Removed = tx->intron2Removed = true;
        if (sp.progress >= 1.0f) {
            sp.complete = true;
            tx->spliced = true; tx->polyadenylated = true;
            tx->matureMRNA = codingMRNA;
            tx->polyATail.assign(48, 'A');
            clearSplicingState(sp);
        }
    }

    // ── Nuclear export ─────────────────────────────────────────
    for (auto& tx : transcripts) {
        if (!tx.active) continue;
        tx.age += dt;
        if (tx.spliced && !tx.exported && !mitoticPause) {
            tx.exportProgress = min(1.0f, tx.exportProgress + dt * 0.60f);
            if (tx.exportProgress >= 1.0f) tx.exported = true;
        }
    }

    // ── Translation ────────────────────────────────────────────
    for (int i = 0; i < 6; i++) {
        auto& tr = translation[i];
        if (!tr.active) {
            tr.timer += dt;
            if (!mitoticPause) {
                int txIdx = findTranscriptForTranslation();
                if (txIdx >= 0 && tr.timer >= (0.8f + i * 0.25f))
                    startTranslation(i, txIdx);
            }
            continue;
        }
        if (mitoticPause) continue;
        tr.timer += dt;
        auto* tx = getTranscript(tr.transcriptIndex);
        if (!tx || !tx->exported) { finishTranslation(i); continue; }
        if (tr.stopReached) { releaseProtein(tr); finishTranslation(i); continue; }
        // tRNA engagement + shuttle + elongation (see ch.~\ref{ch:translation})
        // ...
    }

    // ── Protein folding + PTMs ─────────────────────────────────
    for (auto& p : proteins) {
        if (!p.active) continue;
        p.releaseAge += dt;
        if (!p.folded) {
            p.foldProgress = min(1.0f, p.foldProgress + dt * 0.30f);
            p.folded = p.foldProgress >= 0.999f;
        } else if (!p.mature) {
            p.modificationProgress = min(1.0f, p.modificationProgress + dt * 0.18f);
            p.mature = p.modificationProgress >= 0.999f;
        }
    }
}
\end{lstlisting}

\chapter{Metabolism.h High-Level Step}
\label{ch:metab-pseudo}

\begin{lstlisting}[language={[ISO]C++}]
void MetabolismEngine::step(float dt_s) {
    // Compute all fluxes from current pool state
    Flux f = computeFluxes(pool);

    // Apply stoichiometric updates
    pool.glucose  -= f.glycolysis * dt_s;
    pool.pyruvate += f.glycolysis * 2 * dt_s;
    pool.ATP      += f.glycolysis * 2 * dt_s;  // net ATP from glycolysis
    pool.NADH     += f.glycolysis * 2 * dt_s;

    pool.pyruvate -= (f.LDH + f.PDH) * dt_s;
    pool.lactate  += f.LDH * dt_s;
    pool.acetylCoA+= f.PDH * dt_s;
    pool.NADH     += f.PDH * dt_s - f.LDH * dt_s;

    pool.acetylCoA -= f.TCA * dt_s;
    pool.citrate   += f.TCA * dt_s;
    // ... eight TCA reactions

    pool.O2       -= f.oxPhos * 0.5 * dt_s;
    pool.ATP      += f.oxPhos * 2.5 * dt_s;   // P/O = 2.5 for NADH
    pool.NADH     -= f.oxPhos * dt_s;

    // Glutamine anaplerosis
    pool.glutamine -= f.glutaminolysis * dt_s;
    pool.glutamate += f.glutaminolysis * dt_s;
    pool.alphaKG   += f.glutaminolysis * dt_s;

    // Biomass synthesis
    float biomass_rate = biomassSynth(pool);
    pool.ATP  -= biomass_rate * ATP_PER_BIOMASS * dt_s;
    pool.AA   -= biomass_rate * AA_PER_BIOMASS * dt_s;
    // ... other precursors
    biomass   += biomass_rate * dt_s;

    // Update redox ratio
    pool.NADratio = pool.NAD_plus / max(pool.NADH, 1e-6);
}
\end{lstlisting}

\chapter{MediumField::diffuse Pseudocode}
\label{ch:medium-pseudo}

\begin{lstlisting}[language={[ISO]C++}]
void MediumField::diffuse(float dt_s, float envO2, float envGlucose) {
    for (int s = 0; s < MS_COUNT; s++) {
        float D = D_um2_per_s[s];
        for (int iz = 0; iz < N; iz++) {
            for (int ix = 0; ix < N; ix++) {
                int idx = iz * N + ix;
                float C = c[s][idx];
                float L = 0;
                int neighbors = 0;
                if (ix > 0)     { L += c[s][idx - 1];     neighbors++; }
                if (ix < N - 1) { L += c[s][idx + 1];     neighbors++; }
                if (iz > 0)     { L += c[s][idx - N];     neighbors++; }
                if (iz < N - 1) { L += c[s][idx + N];     neighbors++; }
                float laplacian = L - neighbors * C;  // 5-point stencil
                cNext[s][idx] = C + D * dt_s * laplacian;
            }
        }
    }
    // Swap buffers
    for (int s = 0; s < MS_COUNT; s++)
        std::swap(c[s], cNext[s]);

    // Decay auxiliary fields
    float keep = exp(-PROTEASE_DECAY_RATE * dt_s);
    for (int i = 0; i < CELLS; i++)
        extracellularProteaseLevel[i] *= keep;
}
\end{lstlisting}

\chapter{MediumField::exchange Pseudocode}
\label{ch:medium-exchange-pseudo}

\begin{lstlisting}[language={[ISO]C++}]
void MediumField::exchange(float wx, float wz,
                           const float flux[MS_COUNT],
                           float scale) {
    int ix = (int)((wx / SCENE_BOUND + 1.0f) * 0.5f * N);
    int iz = (int)((wz / SCENE_BOUND + 1.0f) * 0.5f * N);
    ix = clamp(ix, 0, N-1);
    iz = clamp(iz, 0, N-1);
    int idx = iz * N + ix;
    for (int s = 0; s < MS_COUNT; s++) {
        c[s][idx] += flux[s] * scale;
    }
}
\end{lstlisting}

Mass goes into one grid cell, then diffuses out. This is why the
exchange function does not balance itself against an ``outgoing''
pool --- the outgoing pool is whatever called exchange (a
releaseCytosol from a dying cell, a drug uptake, etc.).

\chapter{Cross-module Interaction Matrix}
\label{ch:interactions}

\begin{longtable}{p{3.3cm}p{10.5cm}}
\toprule
\textbf{Module} & \textbf{Depends on}\\ \midrule
Simulation       & MediumField, CentralDogma, IntracellularPhysics, every biochem module via SimCell.\\
CentralDogma     & Constants (HBB, rates), Phase 8A: GeneTemplate.\\
Apoptosis        & Constants (phase dwells, release fractions), triggers from Simulation.\\
Metabolism       & Constants (nutrients), state feedback from Simulation.\\
MembraneTransport& Ion concentrations + SimCell.ATP.\\
VesicularTraffic & UPR $\to$ CHOP $\to$ Apoptosis; autophagy $\to$ AA pool.\\
Cytoskeleton     & RhoA from Signaling; drives adhesion strength.\\
Signaling        & Inputs from receptors (BindingMatcher binding), outputs to cycle, metabolism, apoptosis.\\
GeneRegulation   & TF activities from Signaling; gene on/off $\to$ cycle proteins.\\
DNARepair        & damageLevel from replication errors + drugs; feeds back to p53.\\
CellJunctions    & Adhesion strength, contact inhibition, anoikis.\\
Immunity         & PAMPs from pathogens; cytokines to medium; MHC presentation.\\
Pathogens        & BindingMatcher for entry, CentralDogma for replication (Phase 8E), Apoptosis for spill.\\
Virion/Bacterium & PathogenKinetics for scoring, Simulation::updatePathogens dispatch.\\
ChemicalEntity   & MoleculeCache (geometry), BindingMatcher (affinity).\\
TargetLibrary    & ChemicalEntity (targets), SimCell fields (modulators).\\
\bottomrule
\end{longtable}

\chapter{Lifecycle of a Drug Particle (Rendered)}
\label{ch:drug-particle-life}

Every \code{MediumChemical} entry goes through states
\code{MC\_FREE $\to$ MC\_ATTRACTED $\to$ MC\_BINDING $\to$
MC\_TRANSPORT $\to$ MC\_DESPAWN}. This chapter documents each.

\paragraph{MC\_FREE.} Brownian drift in the fluid. Every tick,
check: is the particle within the uptake radius of any cell?
If yes and the cell can absorb (has the right receptor class),
assign target and switch to \code{MC\_ATTRACTED}.

\paragraph{MC\_ATTRACTED.} Velocity is damped-spring pulled toward
the cell's surface. When distance $<$ attach radius, switch to
\code{MC\_BINDING}.

\paragraph{MC\_BINDING.} Hold at membrane; pulse briefly (shader
reads state and brightens). Dwell $\sim$0.3\,\biosec, then switch
to \code{MC\_TRANSPORT}.

\paragraph{MC\_TRANSPORT.} Interpolate position from surface
into cell centre over $\code{TRANSPORT\_BIOSEC}$ (default
1.0\,\biosec). Shader fades alpha. At the end, switch to
\code{MC\_DESPAWN}.

\paragraph{MC\_DESPAWN.} Swept up at the end of the tick.

This state machine serves the rendering; the actual chemistry of
drug uptake is in \code{updateDrugPK}. The two run in parallel:
chemistry is always happening regardless of particle rendering;
particles are just a visual trace.

\chapter{Lifecycle of a Virion (Rendered)}
\label{ch:virion-life}

Each free \code{Virion} renders as a coloured sphere with 8
spike points (if the spec has $> 60$ spikes per virion). The
colour is hashed from the spec id. Shader pulse brightens the
sphere when the stage is BOUND or ENTERING.

\paragraph{FREE.} Brownian drift; baseline brightness.
\paragraph{BOUND.} On contact with host cell; pulse 1.15--1.4$\times$.
\paragraph{ENTERING.} Same pulse. Virion slowly pulled toward
cell centre for visual effect.
\paragraph{UNCOATING.} Marked for cleanup; usually invisible
because erased before next render.
\paragraph{Intracellular stages} (REPLICATING, ASSEMBLING,
BUDDING) are not rendered as distinct virions; they live inside
the cell. The cell shell tints cyan-green based on viral load
(summed \code{intraVirions} + \code{assembledVirions}).

\chapter{Lifecycle of a Bacterium (Rendered)}
\label{ch:bacterium-life}

Each free \code{Bacterium} renders as three spheres in a row
(dumbbell body) plus two trailing small spheres (flagellum
trails). Colour hashed from spec id, warm green hues.

\paragraph{FREE.} Run-and-tumble motility. Tumble every $\sim$1
\biosec. Run at 0.6\,\wu/\biosec along \code{runDir}.
\paragraph{ADHERING.} On contact; pulse 1.2--1.5$\times$. Stageed
for 120\,\biosec.
\paragraph{INVADED.} Marked for cleanup. Cell surface tints
yellow-orange.
\paragraph{REPLICATING.} Intracellular; not individually
rendered.
\paragraph{LYSED.} Cell dies; pathogens spill; bacteria drift in
medium anew.

\chapter{A Short History of Key Design Decisions}
\label{ch:decisions}

\paragraph{Why 1\,\wu = 5\,\um rather than 1\,\um?} To fit a
dish FOV of a few hundred micrometres into a few tens of world
units, keeping camera coordinate math numerically stable.

\paragraph{Why Metal rather than OpenGL?} Native on macOS; Apple
Silicon performance is excellent; instanced draws are
first-class. OpenGL 4.1 (macOS's max) lacks some instancing
features.

\paragraph{Why ImGui?} Immediate-mode GUIs let the simulator
expose dozens of panels without boilerplate. They are the right
fit for scientific tools.

\paragraph{Why CentralDogma-as-one-gene rather than
full-genome?} Cost/value tradeoff. Simulating every one of
$\sim$20\,000 genes per cell is excessive for the visual +
emergent-drug + emergent-infection scenarios we target. HBB is
the visible reference; Phase 8A adds viral genes transparently;
the GRN layer (Phase 7L) adds regulatory edges that operate on
whichever genes are active.

\paragraph{Why not use SystemsBiology-style SBML
internally?} Export is planned (Phase 7U); internal
representation stays native C++ for runtime speed + clarity.

\paragraph{Why a ``Test Mode'' template?} 16\,GB machines
struggle with the single-cell-mode intracellular particle view.
Users need a lighter path that still exposes the new pathogen +
drug systems. MODE\_COLONY with 1 cell + the opt-in
\code{lightMode} flag does exactly that.

\paragraph{Why no AI/ML anywhere?} The user's core ask was
``no black boxes''. Every decision in the simulator can be
traced through lines of C++. This rules out ML surrogates for
chemistry or fate decisions, but keeps the simulator auditable.

\paragraph{Why header-only biochem modules?} Encourages tight
coupling only through SimCell. Each module can be grepped
independently, compiled independently, replaced independently.

\paragraph{Why Phase 8A before Phase 8B-L?} 8A is the
structural prerequisite: without multi-gene transcription, there
is no place for viral mRNA to go. Every downstream phase of Part
8 assumes the gene registry exists.

\chapter{Open Problems}
\label{ch:open}

\begin{enumerate}
  \item The autocatalytic replication in Part 7 is still a toy.
        Phase 8E replaces it with real transcription through the
        Phase 8A gene registry. Required work: $\sim$3 sessions.
  \item The entry machinery (clathrin, V-ATPase, endosome fusion)
        is a state-machine dwell rather than a physics event.
        Phase 8C+D replaces with real VesicularTraffic events.
  \item The receptor decoration is a Fibonacci shell layer, not
        real IntraParticle entries. Phase 8B fixes this.
  \item The infection chain does not currently propagate in test
        mode because there is only one cell. Running in colony
        mode would demonstrate it, but test mode is the
        16\,GB-friendly path; ideally we'd have a hybrid where a
        small cluster of cells (5--10) runs in colony mode with
        the light rendering.
  \item Mass-balance discipline is currently enforced by
        convention. A future \code{checkBalance()} overhaul could
        emit a stack trace identifying which function broke the
        invariant --- a significant debugging aid.
\end{enumerate}

\chapter{Reviewer Questions}
\label{ch:reviewer}

Likely questions from a reviewer who wants to trust or extend
this codebase:

\paragraph{Q: Can I run this on Linux?}

Currently no --- Metal is macOS-only. A Vulkan backend would be
straightforward given the instance-based draw calls; $\sim$2
sessions of work.

\paragraph{Q: How do I add a new drug?}

Add a row to \file{data/bioagents/drugs.csv} with (id, SMILES,
name). Run the optional RDKit preproc (\file{scripts/chem\_worker.py})
to generate Tier 0 descriptors. Restart. The drug appears in the
Drug Lab dropdown.

\paragraph{Q: How do I add a new virion?}

Add an entry to \file{data/pathogens/virions.yaml} with the 23
fields. Restart. Virion appears in Pathogen Lab.

\paragraph{Q: How do I change the cell type from HeLa to
something else?}

Currently the simulator is hard-coded to HeLa-like parameters
(radius, cycle time, metabolism). A CellType catalogue is planned
(Phase 7A) that would let you pick a fibroblast, macrophage, or
stem cell at setup. Implementation $\sim$1 session.

\paragraph{Q: How do I validate against my own reference data?}

Replace \file{data/reference/growth\_curves/ctc\_hela\_cellcount\_seq02.csv}
with your own CSV (same column format). Run
\code{cellsim\_headless T S N} where T is bio-hours target,
S is seconds-per-step, N is initial cell count. Read
\file{logs/compare\_vs\_reference.csv}.

\paragraph{Q: Is the apoptosis cascade calibrated against real
time-course data?}

Qualitatively yes (Ye 2017 MOMP onset, Spencer 2009 execution
duration). Quantitatively no --- the cascade rate constants are
from the literature ballpark but not refit to specific
microfluidic single-cell time traces. Phase L of the plan
includes a rigorous refit.

\paragraph{Q: How do I run thousands of cells?}

Raise \code{MAX\_CELLS} in \file{Constants.h} (current cap 2000)
and increase your memory headroom. Performance should scale
linearly up to the point where \code{MediumField::diffuse} (fixed
at 12$\cdot$64$\cdot$64 = 49\,152 voxel updates per tick) becomes
negligible relative to per-cell work.

\paragraph{Q: How trust-worthy is the drug phenotype?}

Qualitatively trustworthy: cisplatin arrests in G2 + kills,
paclitaxel induces mitotic catastrophe, doxorubicin damages DNA.
Quantitatively uncertain: the Tier 0 BindingMatcher is a crude
descriptor fit; Tier 3 GNINA docking would tighten the $K_d$
estimates by an order of magnitude. Phase H of the plan adds
Tier 3.

\paragraph{Q: The source code has 32 compilation warnings. Should
I worry?}

All are \emph{unused function/variable} warnings in scaffolded
headers (DNA repair, Proteome, etc.) that aren't yet driven at
runtime but are present for future phases. Zero are
correctness-affecting.

\paragraph{Q: Why C++20?}

For \code{std::unordered\_map}, ranged for, structured bindings,
\code{constexpr} lambdas, and the convenience of using
\code{simd\_float3} in aggregates. Could compile as C++17 with
minor changes.

\chapter{A Plea for Biological Realism}
\label{ch:realism}

The entire thrust of \textsc{CellSim} is \emph{don't lie}. A lie
in this context is a line of code that implements a specific
biological outcome without tracing that outcome to a cause. When
you see code like
\begin{verbatim}
  if (drug == CISPLATIN) cell.die();
\end{verbatim}
that is a lie; it is not a model of cisplatin.

The chemistry-first discipline rejects this. Cisplatin is a
structure: SMILES \texttt{N.N.Cl[Pt]Cl}. The structure determines
its permeability (Lipinski). The structure determines which target
it fits (BindingMatcher). The bound target alters a cellular field
(modulator). The altered field propagates through the biology
networks the cell already has (CDK, damage, p53). The cell decides
to die (apoptosis cascade). The cell spills mass into the medium
(release\_*). Neighbouring cells sense DAMPs and are stressed.

At every step, the line of code has a cause and an effect. No flag
says \texttt{MOA = KILLS\_IN\_S\_PHASE}. The killing \emph{emerges}
from the sequence. If you change cisplatin's structure to something
else, the emergence changes. This is what it means to be a
\emph{simulator} rather than an \emph{animation}.

Part 8 extends this discipline to pathogens. A virion is a
structure: a capsid shape, an envelope flag, a genome sequence, a
spike descriptor. Its entry, replication, assembly, and egress all
trace to lines of code that have causes and effects. The toy
replication in Part 7 is a violation of this discipline --- it
says ``replicate at rate $r$'' without asking the cell for its
ribosomes. Phase 8E fixes this.

\chapter{Summary}
\label{ch:summary}

You have now read a complete reference to \textsc{CellSim}. You
know:

\begin{itemize}
  \item The coordinate system (1\,\wu = 5\,\um) and time base
        (180\,\biosec/\realsec).
  \item The update loop (physics, metabolism, cycle, mitosis,
        central dogma, fate, apoptosis, drug PK, pathogens,
        diffusion).
  \item The SimCell structure (every field, every update site).
  \item The cell cycle ODE (nine species, rate constants).
  \item The mitosis pipeline (eight phases, SAC gate).
  \item The apoptosis cascade (11 triggers, seven phases, Bcl-2
        / Bax / MOMP / caspase).
  \item Replication (3 origins, 4 forks, mismatch repair, post-
        replicative forced completion).
  \item Transcription (3 Pol II slots, Phase 8A routing through
        \code{geneSequence}).
  \item Splicing (2 spliceosomes) + nuclear export.
  \item Translation (6 ribosomes, 12 tRNAs, per-transcript start
        codon since Phase 8A).
  \item The vesicular traffic engine (11 compartments, 4 coats,
        pH gradient, UPR, autophagy).
  \item The membrane transport engine (ion pools, Nernst, Goldman,
        Hodgkin-Huxley).
  \item Metabolism (glycolysis, TCA, OxPhos, Warburg, biomass
        synthesis).
  \item Cytoskeleton (actin, MT, IF, myosins, RhoA).
  \item MediumField (64$\times$64 PDE, 12 species, mass-balance
        invariant).
  \item IntracellularPhysics (Langevin BBK Verlet, Hertz contact,
        attraction fields, spatial hash).
  \item ChemicalEntity (universal structure record), BindingMatcher
        (descriptor score, Kd mapping), TargetLibrary (13 targets +
        modulators).
  \item Virion, Bacterium, PathogenKinetics (shape-based binding),
        PathogenRegistry (YAML loader).
  \item Apoptotic bodies (spawn, drift, swell, burst, efferocytosis).
  \item Rendering (Metal pipelines, cell shell, atom, bond, fluid).
  \item UI (setup overlay, drug lab, pathogen lab, test actions,
        per-cell panels).
  \item Export (6-file CSV bundle).
  \item Headless test suite (CTC validator, infection smoke test).
  \item CMake build (three targets, vcpkg dependencies).
  \item Calibration (mean $|$rel err$|$ vs CTC seq02 = 5.1--7.4\%
        across 5 seeds; seed-swept mean $\approx 6.2\%$).
  \item Phase 8A (multi-gene Central Dogma, landed 2026-04-20).
  \item Phase P1 (determinism audit; \code{--seed} CLI + seeded
        SimRng wrapper; bit-identical headless CSVs under pinned
        seed, 2026-04-20).
  \item Phase G2 (TP53↔MDM2 stress-response loop with
        literature-calibrated Kd + half-lives; \code{cellsim\_p53\_pulsing}
        validator with 4 assertion gates, 2026-04-20).
  \item Phase G3 (MDM2 mRNA intermediate + Hill(n=8) transactivation;
        Purvis-2012-style sustained p53 pulsing under damage;
        validator extended with 4 additional gates for peak count,
        mean period, stress-scaling, and control specificity;
        zero drift on seq02 preserved at 5.1\%, 2026-04-20).
  \item Phase 8B-L roadmap (16 remaining sessions).
\end{itemize}

You can extend the simulator, write new biochem modules, add new
drugs / pathogens / cell types, and validate against your own
reference data, with full traceability from constant to phenotype.

Thank you for reading.

\chapter{Colophon}
\label{ch:colophon}

This document was typeset in \LaTeX{} using the \texttt{book} class
at 11\,pt, Latin Modern fonts, and \texttt{microtype} for
light hyphenation tweaks. The listings environment uses a
slightly darkened Solarized-adjacent palette. Code is rendered in
a monospaced font at footnote size.

The document is not pre-typeset in the repository; the user is
expected to compile it locally with any \LaTeX{} distribution
(\texttt{pdflatex cellsim\_full\_reference.tex} twice will resolve
the TOC and cross-references).

\vspace{2em}
\begin{center}
\emph{CellSim: Computational Cell Biology Simulator.}\\
\emph{Full Technical Reference, 2026-04-20 revision.}\\[1em]
\emph{Compiled from an exhaustive source-code scan.}
\end{center}

% ═════════════════════════════════════════════════════════════════════
\appendix

\chapter{File-by-file index}
\label{app:files}

\begin{longtable}{lrp{8cm}}
\toprule
\textbf{File} & \textbf{Lines} & \textbf{Role}\\ \midrule
\file{src/main.mm}                             & 9159 & Metal shell, ImGui UI, rendering, input handling, export\\
\file{src/simulation/Simulation.h}             & 2826 & Core Simulation class, SimCell struct, per-tick update pipeline\\
\file{src/simulation/CentralDogma.h}           & 2321 & Transcription, replication, splicing, translation, Phase 8A gene registry\\
\file{src/simulation/IntracellularPhysics.h}   & 514  & Langevin + attraction + confinement particle simulator\\
\file{src/simulation/MediumField.h}            & 337  & Dish-level 2D diffusion PDE + mass-balance invariant\\
\file{src/simulation/CellCycleProgram.h}       & 121  & Cycle + mitosis program bundle on each cell\\
\file{src/simulation/TelemetryLog.h}           & 277  & Per-tick CSV logging for headless validation\\
\file{src/simulation/biochem/Apoptosis.h}      & 323  & 11-trigger multi-threshold death cascade\\
\file{src/simulation/biochem/BindingMatcher.h} &  91  & Tier-0 descriptor-based drug $\times$ target affinity\\
\file{src/simulation/biochem/Bioagent.h/.cpp}  &  54  & Drug catalogue loader\\
\file{src/simulation/biochem/Cancer.h}         & 189  & Oncogene/TSG + metastasis scaffolding\\
\file{src/simulation/biochem/CellBiologyEngine.h} & 253 & Umbrella wrapper around metabolism + proteome\\
\file{src/simulation/biochem/CellJunctions.h}  & 208  & Cadherin/desmosome/tight-junction/integrin\\
\file{src/simulation/biochem/ChemicalEntity.h} & 139  & Universal structure record\\
\file{src/simulation/biochem/Cytoskeleton.h}   & 271  & Actin, microtubule, myosin, GTPases\\
\file{src/simulation/biochem/DNARepair.h}      & 435  & NER, BER, HR, NHEJ, TLS\\
\file{src/simulation/biochem/Development.h}    & 225  & Morphogen gradients, positional info\\
\file{src/simulation/biochem/GeneRegulation.h} & 235  & TF--target graph scaffold\\
\file{src/simulation/biochem/GenomeModel.h}    & 259  & Ploidy, chromatin, epigenetics placeholders\\
\file{src/simulation/biochem/Immunity.h}       & 271  & TLRs, PAMPs, MHC, V(D)J repertoire\\
\file{src/simulation/biochem/MembraneTransport.h} & 179 & Ion pools, pumps, H--H gating\\
\file{src/simulation/biochem/Metabolism.h}     & 962  & Central-carbon + TCA + OxPhos + glutaminolysis\\
\file{src/simulation/biochem/PathogenKinetics.h}&123  & Virion/bacterium binding scoring\\
\file{src/simulation/biochem/PathogenRegistry.h}&210  & YAML loader for virion + bacterium specs\\
\file{src/simulation/biochem/Pathogens.h}      & 232  & Per-cell infection state + engine\\
\file{src/simulation/biochem/Proteome.h}       & 380  & Chaperones, HSF1, PTMs, UPR\\
\file{src/simulation/biochem/Signaling.h}      & 588  & Wnt, Notch, TGF-$\beta$, MAPK, PI3K--AKT\\
\file{src/simulation/biochem/StemCells.h}      & 253  & Self-renewal vs differentiation\\
\file{src/simulation/biochem/TargetLibrary.h/.cpp}& 59 & Drug-target catalogue + modulators\\
\file{src/simulation/biochem/VesicularTraffic.h} & 281 & Compartments, coats, SNAREs, UPR, autophagy\\
\file{src/simulation/biochem/Virion.h}         & 111  & Virion spec + instance\\
\file{src/simulation/biochem/Bacterium.h}      &  97  & Bacterium spec + instance\\
\file{src/core/Constants.h}                    & 523  & Every constant: geometry, time, cycle, apoptosis, drugs\\
\file{tests/headless\_doubling.cpp}            & 285  & CTC growth-curve validator\\
\file{tests/headless\_infection.cpp}           & 104  & Phase 7 infection smoke test\\
\bottomrule
\end{longtable}

\chapter{Glossary}
\label{app:glossary}

\begin{description}[style=nextline]
\item[bio-second (\biosec)] One second in the simulated biology clock.
  Default mapping: $1\,\realsec = 180\,\biosec$ at \code{timeScale = 1}.
\item[MOMP] Mitochondrial outer membrane permeabilisation.
  Point of no return in intrinsic apoptosis.
\item[HBB] Human beta-globin gene (NCBI Gene ID 3043). The real 1558-bp
  sequence embedded in \code{CentralDogma.h} as \code{HBB\_SEQUENCE}.
\item[SAC] Spindle assembly checkpoint. Holds anaphase entry until every
  chromosome is bi-oriented.
\item[CTC seq02] Cell Tracking Challenge, Fluo-N2DL-HeLa, sequence 02.
  The primary validation reference (46 bio-h, 125-cell start).
\item[world unit (\wu)] Distance unit of the simulator; 1\,\wu $= 5\,\um$.
\item[apoptotic body] Membrane-bounded fragment of a dying cell,
  carrying cytosol and/or organelle mass, 5--15 per death.
\item[MediumField] The dish-level reaction-diffusion grid.
\item[IntraParticle] A discrete particle inside a focused cell
  (ribosome, receptor, metabolite molecule, etc.).
\item[Tier 0--4] Chemistry-pipeline tiers: 0 = RDKit descriptors,
  1 = OpenFF force field, 2 = xTB semi-empirical QM, 3 = GNINA docking,
  4 = OpenMM MD.
\item[Phase 8A] The multi-gene Central Dogma upgrade landed 2026-04-20.
\end{description}

\chapter{Bibliography}
\label{app:bib}

\begin{thebibliography}{99}
\bibitem{alberts} B.\ Alberts et al.\ \emph{Molecular Biology of the Cell},
  7th edition, Garland Science, 2022.
\bibitem{biona} BioNumbers database, \url{https://bionumbers.hms.harvard.edu}.
\bibitem{ensembl} Ensembl REST API, \url{https://rest.ensembl.org}.
\bibitem{chao2019} M.~Chao et al.\ \emph{HeLa cell cycle dynamics.}
  \emph{eLife}, 2019.
\bibitem{park2016} J.~Park et al.\ \emph{HeLa metabolomics.}
  \emph{Cell Rep.}, 2016.
\bibitem{ye2017} L.~Ye et al.\ \emph{MOMP kinetics.}
  \emph{Sci. Rep.}, 2017.
\bibitem{spencer2009} S.~Spencer et al.\ \emph{Non-genetic origins of
  cell-to-cell variability.} \emph{Nature}, 2009.
\bibitem{silva2010} J.~Silva et al.\ \emph{Secondary necrosis.}
  \emph{FEBS Lett.}, 2010.
\bibitem{chen2016} S.~Chen et al.\ \emph{Nanopore-driven osmotic lysis.}
  \emph{Curr.\ Biol.}, 2016.
\bibitem{ravichandran2015} K.~Ravichandran. \emph{Efferocytosis.}
  \emph{Nat.\ Rev.\ Immunol.}, 2015.
\bibitem{appelqvist2013} H.~Appelqvist et al.\ \emph{Lysosomal digestion.}
  \emph{J.\ Cell Biol.}, 2013.
\bibitem{aldridge2006} B.~Aldridge et al.\ \emph{Apoptosis signalling.}
  \emph{Nat.\ Cell Biol.}, 2006.
\bibitem{kholodenko2006} B.~Kholodenko. \emph{MAPK cascade.}
  \emph{Nat.\ Rev.\ Mol.\ Cell Biol.}, 2006.
\bibitem{berg1972} H.~Berg, D.~Brown. \emph{Chemotaxis in \emph{E.\ coli}
  analysed by three-dimensional tracking.} \emph{Nature}, 1972.
\bibitem{gillespie1977} D.~Gillespie. \emph{Exact stochastic simulation
  of coupled chemical reactions.} \emph{J.\ Phys.\ Chem.}, 1977.
\bibitem{tewey1984} K.~Tewey et al.\ \emph{Doxorubicin and Topo II
  poisoning.} \emph{Science}, 1984.
\bibitem{ramsay1997} R.~Ramsay. \emph{Complex-I inhibitors.}
  \emph{J.\ Biol.\ Chem.}, 1997.
\bibitem{harrison2015} S.~Harrison. \emph{Viral membrane fusion.}
  \emph{Annu.\ Rev.\ Biochem.}, 2015.
\bibitem{mercer2010} J.~Mercer, U.~Greber, A.~Helenius. \emph{Virus entry
  by endocytosis.} \emph{Annu.\ Rev.\ Biochem.}, 2010.
\bibitem{huotari2011} J.~Huotari, A.~Helenius. \emph{Endosome maturation.}
  \emph{EMBO J.}, 2011.
\bibitem{pinto1992} L.~Pinto et al.\ \emph{Influenza M2 proton channel.}
  \emph{Cell}, 1992.
\bibitem{ganser2012} B.~Ganser-Pornillos et al.\ \emph{Retroviral capsid
  assembly.} \emph{Curr.\ Opin.\ Virol.}, 2012.
\bibitem{casparklug1962} D.~Caspar, A.~Klug. \emph{Physical principles in
  virus construction.} \emph{Cold Spring Harb.\ Symp.}, 1962.
\bibitem{galancollmer} J.~Gal{\'a}n, A.~Collmer. \emph{Type III secretion
  systems.} \emph{Science}, 1999.
\bibitem{ronald2010} A.~Ronald, B.~Beutler. \emph{Innate immune
  system---RIG-I / Type-I IFN.} \emph{Science}, 2010.
\bibitem{crank1975} J.~Crank. \emph{The Mathematics of Diffusion.}
  Oxford Univ.\ Press, 2nd ed., 1975.
\bibitem{cornish2012} A.~Cornish-Bowden. \emph{Fundamentals of Enzyme
  Kinetics.} 4th ed., Wiley-VCH, 2012.
\bibitem{brenda} BRENDA enzyme database,
  \url{https://www.brenda-enzymes.org}.
\bibitem{berg1977} H.~Berg, E.~Purcell. \emph{Physics of chemoreception.}
  \emph{Biophys.\ J.}, 1977.
\bibitem{singh2013} J.~Singh, R.~Padgett. \emph{Rates of in vivo
  transcription elongation.} \emph{Nat.\ Struct.\ Mol.\ Biol.}, 2013.
\bibitem{sia2008} J.~Sia et al.\ \emph{Cisplatin death kinetics in HeLa.}
  \emph{Chem.\ Biol.}, 2008.
\end{thebibliography}

\printindex
\end{document}
